% =====================================================================
%  R. Nielsen, DEN SPECULATIVE LOGIK I DENS GRUNDTRÆK.  Hefter 1--3.
%  Kjøbenhavn, 1841--42.  192 pp.  Fraktur.  (Hefte 4 is not in the scanned volume.)
%
%  TRANSCRIPTION -- GENERATED; DO NOT EDIT BY HAND.  Rebuild, chained (&&), from the repo root:
%    python3 ebook-method/speculative-logik/draft.py --emit --force
%      this template + the scan's text layer (KB's ABBYY OCR in the Royal Library's scan,
%      https://www.kb.dk/e-mat/dod/11140802669D.pdf; ~/bibliotek/Nielsen, Rasmus/speculative-logik-1841.pdf, see SCANS.tsv);
%    python3 ebook-method/collate/check.py speculative-logik --witness ocr --apply
%      the decisions in decisions.tsv, taken against tesseract's independent reading of the images;
%    python3 ebook-method/speculative-logik/oepass.py --apply
%      the ø pass (lexicon + tight word crops read at the image, oe-reads.tsv);
%    python3 ebook-method/speculative-logik/finish.py --apply
%      footnotes, headings, emphasis, pp. 177-178, corrections across markup, PRINTED AS IS marks.
%
%  WORD-ACCURACY (2026-10-07).  No e-book, no earlier transcription, no Wikisource text exists.  The draft
%  (KB's ABBYY layer) was collated word by word, with punctuation and word division, against tesseract
%  (tessdata_best script/Fraktur), with two further independent readings (deu_latf; UB Mannheim frak2021)
%  as votes: 9370 disagreements.  7026 were settled by rule (rules.py, each decision named: each OCR's
%  known Fraktur faults, counted on this book, undone; the layer's split words joined; agreement of the
%  readings; attested in >= 2 other books of the corpus or a compound of such words); 2077 at the image,
%  in BLIND crops read by Sonnet readers who saw neither OCR reading (9 readers, 2 rounds + a residue round;
%  18 of 18 planted controls read right).  A blind audit of 140 rule decisions, mixed into the crops, found
%  1 rule error (that rule narrowed, its 26 decisions withdrawn) and 4 missing ø (caught by the ø pass).
%  The ø (read o by every OCR): lexicon (where only the ø form is a word) + 313 tight word crops read at the
%  image (for/før, bor/bør ...); the layer's s for ø also taken as ø where the word is attested.
%  Footnotes: all 54 notes read whole at the image and set at the marks read there (5 run over a page);
%  the notes of pp. 34 and 38-39 are in body-size type and were set from the collated text, formulas
%  from the image.  Headings: all 19 heading blocks read at the image.  pp. 177-178: the layer lost
%  half-lines; the text there is the full reading at the image.  Emphasis: bold by stroke width (>= 1.5 x
%  the page median taken as bold: a blind audit of 60 found 53 bold, 6 read plain); 405 doubtful words read
%  at the image (bold / spaced / plain).  A final sweep read 535 unattested words at the image (82 mended).
%
%  NOT ESTABLISHED.  156 disagreements left open stand as the layer read them after the rules (mostly in
%  formulas, headings' surroundings and note joints); the words of the text not attested elsewhere and not
%  read in the sweep (the book's own terms: Exclusion, Qvanta, Hæcceiteten ...) were not all seen at the image;
%  emphasis is the least certain layer (see the audit).  A spot-read of pp. 60 and 86 (cropped at the right
%  edge) found no word fault in what it saw beyond spacing before punctuation (mended throughout).
%  PRINTED AS IS: Forudsætuinger (p. 138), Tankeb.vægelse (p. 143), esotoriske (p. 113), and the sweep's.
%
%  SECOND READING BY TRANSLATION (2026-10-07).  Translating the book (five translators) was a full second reading:
%  265 odd spots listed by the translators, and 137 more by prescreen.py (doubled words and marks left by --apply,
%  o: for ɔ:, letter-spaced words left split, heading debris, compound hyphens set as dashes).  241 were cropped and
%  read BLIND at the image (two Sonnet readers; .parts/tr/ts-spots.tsv, ts-reads.tsv); the rest are classes settled by
%  rule after a sample (o: and 9: for ɔ: 5 of 5; Følge for Folge 5 of 5; rc. for ꝛc. 3 of 3) or by both OCRs (doubled
%  words and marks; false paragraph breaks inside a sentence; sheet signatures and folios in the text).  Corrections:
%  finish.py step 7 (trpost.py: 204 by reading, about 190 by rule).  Found so: the thesis of § 20 (p. 81) dropped by the
%  notes step; the last line of p. 136 lost by the layer (restored from tesseract); the note on p. 104 set twice; the
%  antinomic note of p. 109 at the wrong mark; § titles set twice on pp. 1, 51, 58, 70; 12 more misprints (PRINTED AS IS).
%  Still NOT ESTABLISHED: what neither the translators nor the script flagged (the residue rate after this reading is
%  unknown; a sampled page check would measure it); emphasis remains the least certain layer.
%
%  PAGE MAP.  printed p = PDF page - 7 throughout (p. 1 = PDF 8, p. 192 = PDF 199), read off the folios
%  (layer and tesseract on 155 pages; the rest read at the image, 2026-10-07).  PDF 6: title of 1ste Hæfte;
%  PDF 7: the author's note on its verso (both set below); PDF 202: title of 3die Hæfte, bound at the end
%  of KB's copy (set after the text); no title of 2det Hæfte in the volume.
%
%  CONVENTIONS.  Orthography and punctuation as printed.  The Fraktur I/J (one sort) is set I before a
%  consonant, J before a vowel (Intet, Idee, Jeg).  Long s = s; ɔ: as printed.  The print's emphasis is
%  BOLD Fraktur (\textbf) and letter-spacing (\emph); Roman (antiqua) type within the Fraktur is \textit.
%  The compound's double hyphen is "-" (Væren-for-sig).  Footnotes keep the print's marks *) **) per page.
%  Misprints are kept and marked % PRINTED AS IS at the spot.
% =====================================================================
\documentclass[12pt,a4paper,openany]{book}
\usepackage[utf8]{inputenc}
\usepackage[T1]{fontenc}
\usepackage{libertinus}
\usepackage{libertinust1math}
\usepackage{textalpha}
\usepackage{graphicx}
\DeclareUnicodeCharacter{0254}{\reflectbox{c}}  % ɔ (U+0254) = reversed c
\DeclareUnicodeCharacter{221E}{\ensuremath{\infty}}
\DeclareUnicodeCharacter{2032}{\ensuremath{'}}
\DeclareUnicodeCharacter{2033}{\ensuremath{''}}
\DeclareUnicodeCharacter{2034}{\ensuremath{'''}}
\DeclareUnicodeCharacter{2057}{\ensuremath{''''}}
\DeclareUnicodeCharacter{2212}{\ensuremath{-}}
\DeclareUnicodeCharacter{2299}{\ensuremath{\odot}}
\DeclareUnicodeCharacter{A75B}{r}   % r rotunda (ꝛc. = etc.)
\usepackage[danish]{babel}
\usepackage[protrusion=true,expansion=true]{microtype}
\usepackage{setspace}
\usepackage[margin=1.3in]{geometry}
\usepackage{fancyhdr}
\usepackage{hyperref}
\hypersetup{pdftitle={Den speculative Logik i dens Grundtræk}, pdfauthor={Rasmus Nielsen}, hidelinks}
\setlength{\headheight}{15pt}
\pagestyle{fancy}
\fancyhf{}
\fancyhead[LE,RO]{\thepage}
\fancyhead[C]{\textit{R. Nielsen: Den speculative Logik i dens Grundtræk (1841--42)}}
\renewcommand{\thefootnote}{\ifcase\value{footnote}\or *)\or **)\or ***)\else
  \arabic{footnote})\fi}
\setstretch{1.3}
\usepackage{marginnote}
\newcommand{\opage}[1]{\marginnote{\footnotesize\textit{[#1]}}}
\newcommand{\apage}[1]{\marginnote{\footnotesize\textit{[#1]}}}
% headings as printed: Capitel (bold, with its title), Afsnit A./B./C. (bold), a)/b)/c) (spaced),
% the § number and its argument (bold), and the centred "Overgang." (spaced)
\newcommand{\capitel}[2]{\chapter*{\centering\textbf{#1}\\[0.3em]\textbf{#2}}\addcontentsline{toc}{chapter}{#1 #2}}
\newcommand{\afsnit}[1]{\section*{\centering\textbf{#1}}\addcontentsline{toc}{section}{#1}}
\newcommand{\underafsnit}[1]{\subsection*{\centering\emph{#1}}\addcontentsline{toc}{subsection}{#1}}
\newcommand{\paragraf}[2]{\begin{center}§ #1.\\[0.2em]\textbf{#2}\end{center}}
\newcommand{\overgang}{\begin{center}\emph{Overgang.}\end{center}}

\begin{document}

% --- front matter, read at the image 2026-10-07 (.parts/notes/r-struct.tsv, F6, F7).
% PDF 6, the title of 1ste Hæfte, in a Greek-key border (not reproduced); a short rule (drawn plain).
\begin{titlepage}
\begin{center}
\vspace*{4em}
{\Huge\textbf{Den speculative Logik}}\\[1em]
{\Large i dens Grundtræk}\\[1em]
af\\[1em]
{\Large\textbf{Lic. Rasmus Nielsen,}}\\[0.5em]
Prof. extr. i Philosophien ved Kjøbenhavns Universitet.\\[2em]
\rule{6em}{0.4pt}\\[2em]
{\large\textbf{1ste Hæfte.}}\\[1em]
\textit{1841.}
\end{center}
\end{titlepage}
% PDF 7, the verso: the author's note.
\thispagestyle{empty}
\vspace*{8em}
Disse Grundtræk ere at betragte som Fragment af en philosophisk Methodologi, hvis første Deel skulde
indeholde Logiken med en forudskikket Indledning. Nødvendigheden af at have en trykt Veiledning for det
mundtlige Foredrag har fremskyndet Udgivelsen. De øvrige Hæfter ville følge, efterhaanden som den i
Lectionscatalogen annoncerede Forelæsning rykker frem.

\medskip
Kjøbenhavn, den 10de November 1841.\hfill\textbf{R. Nielsen.}
\clearpage

% ===== TEXT (pp. 1--192) =====
% --- p. 1 ---
\setcounter{footnote}{0}%
\opage{1}\capitel{Første Capitel.}{Væren.}
\afsnit{A. Qvalitet.}
\underafsnit{a) Den rene Væren.}
\paragraf{11}{Væren, Intet, Vorden.}

\emph{Væren} er den meest omfattende og just derfor den
meest indholdsløse af alle positive Bestemmelser; den har
som saadan Intet til Grændse. Ifølge denne Begrændsning maae Væren og Intet udelukke hinanden,
men netop derfor ogsaa bestemme hinanden,
ɔ: sættes som hinandens Prædikater. Denne dobbelte
Modsigelse, hvormed ethvert af Leddene er behæftet,
ophæves i Vorden, som indeholder begges uendelige
Alternation.
 Copula „er“, der, skjult eller aabenbart, udtrykker Forbindelsen imellem ethvert Subject og dets Prædikater, antyder, naar vi nærmere betragte det, at alle tilværende Ting, alle concrete Bestemmelser, saavel de væsentlige som de uvæsentlige, have deres fælleds Medium i Væren. Væren forholder sig da, umiddelbart taget, til Alt, som Rummet til de sandselige Gjenstande. Tænke vi alle Gjenstande borte, saa bliver 
% --- p. 2 ---
\setcounter{footnote}{0}%
\opage{2}kun den tomme Plads, det rene Rum tilovers; saaledes
ogsaa, naar vi abstrahere fra alle Tilværelsens
specielle Bestemmelser, faae vi den rene Væren tilbage.
Forsaavidt Tingene altsaa falde sammen med
Væren, falde de derved sammen med hverandre indbyrdes; forsaavidt de derimod adskille sig fra hverandre,
falde de udenfor den rene Væren. Den rene Væren
maa man derfor ikke tænke sig som et Phantasibillede;
for at søge den behøver man ei at stirre ind i
den blaae Luft eller ned i Tankernes Afgrund; den er
ikke at betragte som en fremmanet Gestalt, der kun
viser sig i en mystisk Midnat, men som et aldeles simpelt
Tilværelsesmoment (momentum simplicissimum),
der overalt er os lige nær. Den rene Væren gaaer
saavel gjennem Træ og Steen som igjennem Bevidstheden
og dens Tanker; men Hovedsagen er, at Øiet ikke blændes af de heterogene Objecters utallige Forskjelsbestemmelser.

Hvad vi her maae fastholde, er da kun dette, at,
hvor forskjellige Tingene end iovrigt ere, saa gjælder
det dog om dem alle, at de \textbf{ere}.

Da der nu ingen Forskjelsbestemmelser ere i den
rene Væren, og da Forskjellighederne, hvis de som saadanne
falde udenfor Væren, maae sættes som ikke-værende:
maa al Forskjel forsvinde. Denne Forsvinden
har sin positive og negative Form. Fra den positive
Side fremstiller Forholdet sig saaledes, at Væren selv
breder sig ud over alle Forskjelsmomenter, assimilerer
dem alle med sig, idet den selv er det Urelement,
hvoraf de alle ere dannede; fra den negative Side derimod
saaledes, at enhver Forskjel som saadan falder
udenfor Væren og derfor ophører at være. Tale vi
saaledes om A"s og B"s Væren, da synes de begge
hver for sig at have et Moment, der gaaer op i Vær%
% --- p. 3 ---
\setcounter{footnote}{0}%
\opage{3}en, et andet derimod, der falder udenfor den. A kunne
vi da tænke os deelt i a og α, B i b og β; falde
nu a og b indenfor Væren, saa falde α og β udenfor. Da nu a i dets Forskjel fra b virkelig existerer,
saa maa man igjen kunne tale om a"s og b"s Væren, og paa Ny halvere a og b; den tilbageblevne Differents
indsvinder derved til det Halve, denne kan igjen
halveres, og saaledes i det Uendelige, indtil Væren
har udstrakt sin Magt over Alt og erobret al Forskjel.
I den negative Form see vi strax, at da a og b falde
udenfor Væren, saa ere de slet ikke; de tilintetgjøres
da; man sætter og erkjender dem i deres Intet.

I sin første Form maa altsaa Værens Udstrækning
fremstilles som: Væren + Væren + \dots{}; i sidste derimod som Væren + Intet.

I begge Former sættes Væren som ubegrændset. Men den sidste Udtryksmaade er langt bestemtere end
den første; thi medens denne blot antyder et Tankeforsøg,
fremstiller hiin derimod det endelige Resultat. \textbf{Væren har altsaa Intet} \textbf{udenfor sig}.

Men selv det herved udbragte Resultat er behæftet
med en stor Tvetydighed og det en saadan, der ligger
i Sproget selv. Den Setning: Intet begrændser Væren,
kan udtrykke: at Væren overstiger alle Grændser,
at der ikke er Noget, der kan sætte den Skranker, at den er positivt ubegrændset; at Intet virkelig
kan sætte Væren en Skranke, at hvor Intetheden
kommer til Herredømme, maa Væren ophøre. Her gives altsaa virkelig en Definition af Værens Ubegrændsethed. Den ubestemte Væren bliver tænkt i
dens hele Ubestemthed som saadan, ɔ: dens Ubestemthed
sættes som dens Bestemmelse. Men for at faae nogen
Bestemmelse ind i Væren, der absorberer alle Bestemmelser, maa man forbinde en bestemt Tanke  med det
% --- p. 4 ---
\setcounter{footnote}{0}%
\opage{4}ellers ubestemte Intet. Væren, der ophæver al Forskiel,
kan altsaa dog ikke faae Bugt med det almindeligste Forskjelsmoment, med Negationen som saadan,
med Ikke-Væren, eller med Intet. \textbf{Væren er da
ikke} Intet, \textbf{og} Intet \textbf{ikke Væren}. Men i sin rene Umiddelbarhed kan denne Forskjel
heller ikke holde sig. Prøver man f. Ex. paa at tænke
Væren som Indbegrebet af alle Realiteter, hvis Modsætninger
ere udjævnede, altsaa som den rene gedigne,
universelle Realitet selv, saa bliver Intet at forestille
som Tomhedens rene Afgrund. Tomheden har altsaa
sin Væren ved Siden af Realiteten; men det er characteristisk
for den tomme Væren, at den kun indeholder
sig selv, og Realitetens Væren indeholder ligeledes
kun sig selv; Realitet og Tomhed ophøre da her at
være Modsætninger: vi finde Intet i Væren og Væren
i Intet. Det Spørgsmaal: hvad er den rene Væren,
kræver, at Væren gjennem sit Andet skal affirmere sig
selv. Væren søger altsaa her sine forskiellige Bestemmelser:
x + y + z. Kunne vi finde disse, saa kunne
vi ogsaa definere Væren. Vi have allerede fundet
dem; thi det er aabenbart, at det, der søges i x + y + z, maa være Forskjelsmomenter, og da vi nu vide, at al
Forskjel er forsvunden, maae vi i deres Plads sætte: Intet + Intet + Intet. Da nu Intet her ganske i
Almindelighed betegner en savnet Forskiel, saa er Intet i denne almindelige Form = Intet. Altsaa: Væren
kommer gjennem Intet til sig selv: \textbf{Væren er
Intet}. Den Sætning: Væren er Intet, er det reneste
Udtryk for den Grundlov, at al Eenhed som saadan
er sin egen Modsætning, ɔ: at Modsætningen udspringer
af Eenheden selv. Fra det positive Synspunkt kunne vi nu sige, at den rene Væren her er i Begreb
% --- p. 5 ---
\setcounter{footnote}{0}%
\opage{5}med at restituere alle de Modsætninger, den forhen
absorberede, skjøndt disse endnu ikke kunne characteriseres
uden ved den Fælledsbestemmelse, at de overalt
skulle udgjøre den rene Værens Intet. Idet Væren sætter sit eget Intet, giver den derved alle Modsætninger
deres specifiske Væren, og sætter sig derved
som den \textbf{absolute} Væren. Fra det negative Standpunkt udtrykker Sætningen
derimod dette, at den rene Væren som saadan er et
Phantom, en indholdsløs Abstraction, og at den sande
Væren maa søges i de Modsætninger, hvorfra vi have
abstraheret, ɔ: i de værende Ting. Denne Tanke fastholder
Bevidstheden med Lethed. Naar alt Værende
tages bort, bliver der Intet tilbage. Den for alt Værende
blottede Væren er da det rene Intet. Vi kunne nu ogsaa vende Sætningen om og sige:
\textbf{Intet er den rene} \textbf{Væren}. Saa negativt det rene
Intet end er under alle Skikkelser, saa viser det sig
dog overalt som en Væren. Det forestillede Intet er
den tomme Plads, hvorfra Alt er taget bort; men forestillede
man sig nu Pladsen borttagen, saa maatte
man tillige forestille sig den Plads, hvor Pladsen
havde været. Intet vedbliver altsaa at fremstille sig
som den værende, skjøndt tomme, Plads, ɔ: som den
Væren, hvis concrete Bestemmelser ere bortfjernede,
ɔ: som den rene Væren. Denne Forestilling retfærdiggjøres
ogsaa af den strænge Tænkning. Det rene Intet har nemlig ikke hjemme i Sandselighedens, men i
Tankens Verden. Men Tænkningens negative Operation, der ophæver al Umiddelbarhed, kan ikke komme
til det rene Intet uden ved at negere Umiddelbarheden
i dennes universelleste Skikkelse, ɔ: den rene Væren. Det rene Intet er altsaa Reflexen af den rene
% --- p. 6 ---
\setcounter{footnote}{0}%
\opage{6}Væren selv, overført fra Umiddelbarhedens i Reflexionens
Form.

Herved ere vi komne til tvende hinanden modsigende
Setninger: a) Væren er ikke Intet (Intet er
ikke Væren), og b) Væren er Intet (Intet er Væren). Da vi ikke kunne fornegte nogen af disse Sætninger,
medens den ene dog ophæver den anden, saa er der
ingen anden Udvei end at urgere begges Overgang i
hinanden.

Da vi begyndte at tænke den rene Væren, blev
Tanken dreven over til det modsatte Extrem, det rene
Intet, og under Bestræbelsen for at fastholde dette,
førtes den igjen tilbage til Væren. Tanken var altsaa
virkelig i Bevægelse, fordi dens Object bestandig
bevægede sig; men den mærkede ikke Bevægelsen som
saadan. Den forsøgte derfor først ganske umiddelbart
at opfatte Væren som hvilende i sig, og da den kom til
Intet, troede den ligeledes at kunne fastholde dette i dets
Reenhed; nu veed den derimod, at den ikke skal slaae
sig til Ro i nogen af de nævnte Categorier. Den
nye Categorie, der opklarer Værens og Intets Overgang
i hinanden, er \textbf{Vorden}.

Naar Tanken skal fastholde Overgangen som saadan,
da maa den uafbrudt bevæge sig imellem de modsatte
Led; Tanken om Væren kan ikke være absolut simultan
med Tanken om Intet. I denne Oscillation er
begge Momenternes Forskjel bevaret. Saaledes see vi
nu, at Væren idelig maa afvexle med Intet; altsaa
have vi: Væren, Intet, Væren. Men Bevægelsen synes endnu blot at tilhøre Tanken, medens begge de
modsatte Led hvile roligt ved Siden af hinanden.

Den sande Overgang erkjendes først, naar vi fremhæve Momenternes Eenhed. Væren er Intet; idet
den sætter sig selv, sætter den sit Intet; dens Over%
% --- p. 7 ---
\setcounter{footnote}{0}%
\opage{7}gang i Intet beroer da derpaa, at den udvikler sit
skjulte Intet; ligesom Intets Overgang i Væren alene
følger deraf, at det har Væren i sig.

I Vordens Række maa da ethvert af de modsatte
Led \textbf{fordobles}: Væren (a), Intet (a) \textbar{} Intet (b), Væren (b). Ethvert af Leddene indeholder nemlig
baade et sat og et endnu ikke sat Moment, naar vi
fremstille det objectivt, eller, fra det subjective Synspunkt taget, en aabenbar og en skjult Side. Begynde vi saaledes med den umiddelbare Væren (a), da
synes den vel ved første Øiekast kun at ville sætte sig
selv Væren (a) = Væren (a) men betragte vi den
nærmere, da vender den strax sin anden Side frem,
nemlig sit skjulte Intet (a). Idet Intet (a) fremtræder, bliver det derved selv sat, og kan derfor som
Intet (b) blive et nyt og umiddelbart Begyndelsesled,
der, selv oprundet af Væren (a), nu ogsaa maa
sætte sit skjulte Værensmoment, Væren (b). Uagtet
nu Væren (a) selv er Intet (a), og Intet (a) igjen
Intet (b) ꝛc., saa er der dog ligesaavel en Forskjel
imellem Væren (a) og Intet (a) som imellem Intet (a)
og Intet (b); Eenheden og Forskjellen komme og
svinde i det Uendelige: det er en urolig Eenhed i
en urolig Forskjel. Paa Grund af den Modsætning, der stedse finder
Sted imellem det aabenbare og skjulte Moment, er det
ingenlunde ligegyldigt, om man begynder med Væren
som sat, for at den kan sætte sit Intet, eller med et
sat Intet, der skal ponere sin Væren.

Værens Fremgang af Intet er en \textbf{Opstaaen}; Intet
resulterer derimod af Væren ved en \textbf{Forgaaen}. Men
som Følge af Momenternes Identitet er enhver Opstaaen
igjen en Forgaaen, og omvendt.
% --- p. 8 ---
\setcounter{footnote}{0}%
\opage{8}Vorden er altsaa et Begreb, der virkelig resulterer
af Væren og Intet. Men idet alle Begrebets Momenter sættes og udvikles, komme de derved alle tilstede. Væren søger da ikke længere sit Intet, den har
det allerede; Intet gaaer heller ikke over i Væren, thi
da det har sat sin Væren, er der ingen Overgang at
giøre. Momenternes urolige Eenhed kommer derved til
Ro: Vorden ophæves.

\underafsnit{b) Tilværen.}
\paragraf{12}{Noget, Andet, Forandring.}

I \emph{Tilværen}, ɔ: den dobbelte Væren, er Intet med Væren og Væren med Intet. Væren, der
har ophævet sit Intet og integreret sig selv, bliver
derved til Noget, medens Intet, bestemt ved sin
Væren, nu ogsaa slutter sig sammen med sig, og viser sig som \emph{Andet}. Da det er den samme Væren og det samme Intet, der her have fordoblet deres
Eenhed og Forskjel, saa er Noget selv Andet, og Andet Noget: Sandheden ligger i begges \emph{Forandring}. Væren, Intet og Vorden kunne vi nu betragte som
de Grundelementer, hvoraf alle efterfølgende logiske
Categorier udvikle sig; disse følge nu paa hverandre
som fuldkomnere og fuldkomnere Udtryk for Specificationer
af den positive Væren og Negationen i disses
gjensidige Overgang, Begrændsning og Reflexion.

Væren, der gjennem sit Intet begrændser sig selv,
er Tilværen; denne Selvbegrændsnings Momenter
udvikle sig saaledes: A) I Vorden viste det sig, at
% --- p. 9 ---
\setcounter{footnote}{0}%
\opage{9}Intet fremkom af Væren, ligesom Væren af Intet; vi
erholde derved som første Led: Væren (a), Intet, Væren
(b). Ved det mellemtrædende Intet opstaaer en
Modsætningsnuance imellem Væren (a) og Væren (b);
thi der er Forskjel paa den Væren, der gaaer umiddelbart
forud for Negationen, og den, der resulterer af
den ponerede Negation. Vi have altsaa ikke her
simpelthen Væren + Væren i det Ubegrændsede, saaledes
som i § 11, hvor Væren i dens usande abstracte
Umiddelbarhed var = Væren, og derfor ogsaa Intet;
men Væren har her fordoblet sig selv: det er i Tilværen,
at Væren først bliver \textbf{til} Væren. B) Igjennem den fordoblede Væren glimter Intethedens Moment
frem og forsvinder. Intet forsvinder i Tilværen:
thi Negationen, der adskiller Væren (a) fra Væren (b),
skal her kun betegne den Forskjel, der ligger til Grund
for begge Momenters medierede Eenhed. Intet bliver
\textbf{derved selv til Intet}. Analysere vi denne
prægnante Sætnings Indhold, da ligger deri: α) at Intetheden ganske ophører, og overvindes af Tilværen,
β) at Intetheden (Negationen) selv faaer Tilværen,
ɔ: bliver et tilværende Intet. Denne Tilværen følger
ligefrem af § 11, hvor vi saae, at Intet, just fordi
det ikke var Væren, havde Væren i sig. Den eneste
Maade, hvorpaa Tilværen kan slippe fri fra det hule,
tomme Intet, er den, at Intet selv føres ind i Tilværen.
Ligesom Væren gjennem Intet gik sammen
med sig selv, saaledes gaaer nu Intet (a) gjennem
Væren (Intets Negation) sammen med Intet (b) og
bliver et i sig fordoblet, ɔ: et \textbf{tilværende} Intet. C) Herved er Tilværen selv bleven redupliceret; den
maa derfor nærmere bestemmes, som: 1) Tilværen
[Væren (a) Intet, Væren (b)] begrændset ved 2) Tilværen
[Intet (a) Væren, Intet (b) ;] første Led er her
% --- p. 10 ---
\setcounter{footnote}{0}%
\opage{10}at betragte som den gjennemførte Position; andet Led
derimod som den gjennemtænkte Negation; begges Identitet udgjør først den fulde, integrerede \textbf{Tilværen},
der bestemtere kan sættes som \textbf{Tilværelse}.

I Tilværelsen maa nu hvert af de udviklede Led
for sig tænkes i een Tanke, sees i eet Syn; derved
bliver den tilværende Væren at bestemme som \textbf{Noget},
det tilværende Intet som \textbf{Andet}. For Udviklingen
af Noget og Andet er Væren vor faste Basis: det abstracte Intet er forsvundet; vi have kun Værens egne
positive Modsætninger tilbage. I Noget træffe vi et
umiddelbart \textbf{Værende}, men i dets rene Umiddelbarhed er dette blot det første Tilværelse-punkt, der strax
henviser til et \textbf{andet} Værende. Hvis nu denne Henviisning
fortsattes i det Uendelige, vilde Tilværelsen
aldrig blive til Noget; naar derimod Andetheden retorqveres, naar den bliver et Moment, gjennem hvilket
det Værende fordobler og integrerer sig selv, erholder
det sin vedbørlige Grændse, og sættes som Tilværende\footnote{I Fornemmelsen optage vi en Mangfoldighed af Tilværelsens Momenter; men formedelst det Enes idelige Forsvinden i det Andet, kommer det i den blotte Fornemmelse aldrig til en bestemt Opfattelse af Noget.}.
Ligesom den i Tilværen reduplicerede
Væren har Intetheden som ophævet Moment, saaledes
har nu ogsaa det bestemte Noget Andethedens Moment
i sig; men saalænge Andet selv endnu ikke er udviklet, mangler Noget den fulde Bestemthed. Andet er
jo netop Nogets sande Grændse, bliver denne ikke som
saadan fixeret, da bortflyder Noget selv i Tilværelsens Strøm: \textbf{for at Noget skal} \textbf{kunne holde} sig,
\textbf{maa Andet blive} til \textbf{Noget}.
% --- p. 11 ---
\setcounter{footnote}{0}%
\opage{11}Denne Tanke indeholder en betydningsfuld Prægnants; thi den udtrykker: at Andet selv skal understøtte
det positive Noget, ɔ: at Noget maa sætte sit
skjulte Andet, for gjennem denne Andethed at bekræfte
sig selv. Altsaa høre Momenterne: Noget (a) (det begyndende
Noget), Andet (den bestemte Grændse), Noget
(b) (det af Andetheden resulterende Noget), sammen
i een eneste Umiddelbarhed, for at frembringe Nogets
Tilværelse, eller det fuldbaarne Noget; dernæst (b) at
Andet selv bliver Noget, at det holder sig i sin Andethed
som saadan, og soutenerer sin egen Tilværelse.

Tilværelsen af Andet bestemmes ved en videre Udvikling
af det tilværende Intet. At Andet indtager
Negationens Plads, indlyser strax deraf, at det tilhører
Tanken og Reflexionen. Alle de Mediationsmomenter,
der høre til, for at udvikle Noget, maae igjen
ophæves; i Noget prædominerer Umiddelbarheden; Bevidstheden
om et hvilketsomhelst Andet tilbagetrænges,
for at Noget kan sees i dets egen Belysning. Andetheden derimod fremstiller sig strax i et Dobbeltlys;
den vidner ikke ligefrem om sig selv, men tager stedse
Hensyn til det forudgangne Noget. Noget maa være
det \textbf{Første}, dernæst følger det \textbf{Andet} (secundum, af
sequi). Men naar Andet indtræder, er det virkelig
forbi med Noget; vel bliver Andet, derved at det umiddelbart
sætter sig selv, ogsaa til Noget; men dette er
jo \textbf{noget} ganske \textbf{Andet}.

Begge de modsatte Led synes saaledes at holde sig
rolige, hvert paa sin Side. Deres gjensidige Ligegyldighed
hidrører fra deres afsluttede Væren. Noget er
kun sig selv, og bryder sig derfor ikke om Andet. Andet
antyder vel en Exclusion af Noget; men da det
ikke selv gjør Fordring paa at være Alt, men kun et
Andet, saa er denne Exclusion ingen absolut Fortræn%
% --- p. 12 ---
\setcounter{footnote}{0}%
\opage{12}gen. Begge ere, saa at sige, skikkelige Naboer, der
leve i Fred, fordi hver holder sig paa sit Eget. Relationen
kan imidlertid ikke ophæves. Da det, hvorved Noget integreres, netop er et Andet,
og det, hvorved Andet soutenerer sig, er et ophævet Noget, saa resultere de idelig af hinanden. Noget
er det anticiperede Andet, og Andet det explicerede Noget.
Noget bliver derved idelig til Andet, og Andet
til Noget. At Noget selv er sit Andet, og Andet selv
sit Noget, er kun begribeligt derved, at det Andet,
der ligger i Noget, tillige er det Noget, der udvikler
sig af Andet. Saaledes kommer Noget til at begrændse
sig selv; det gjenfinder kun sig selv i ethvert Andet, og
ponerer sin Uendelighed.

ed. Men Andet er selv sit Noget derved, at det har
Noget til sit Andet. Idet nu Noget kommer frem
paa ethvert Punkt, gjenkjender ogsaa Andet overalt sin
egen Andethed, og begrændser ligeledes sig selv som
det Uendelige. Saaledes synes Tilværelsen at faae
Bugt med Intet; men det vil snart vise sig, at det
ligesaavel er Intet, der selv er blevet til, og derved
destruerer den positive Tilværelse. Som Relationens
bestemtere Udtryk faaer Andetheden allevegne Overhaand.
Modsætningen kan ikke forsvinde, uden at vi
igjen synke ned til den abstracte Væren. Vel er Noget
selv sit Andet, men ikke tautologisk. Sandheden er, at
Noget \textbf{forandrer} sig. Andet er det bestemte Intet
af Noget, ligesom Noget det positive Intet for Andet.
Jo skarpere Noget begrændser og udelukker Andet for
at vedligeholde sig selv, desto mere bevirker det sin
egen Opløsning, efterdi det selv er sit Andet. Selvopløsningen hører saaledes med til Begrebet af det tilværende
Noget. Ligeledes naar Andet har Noget til sit
Andet, er Meningen ikke, at Andet umiddelbart er
% --- p. 13 ---
\setcounter{footnote}{0}%
\opage{13}= Andet; men omvendt, at Andet er forskjelligt fra Andet.
I Forandringen bliver den positive Identitet
af Noget og Andet nedsat til at være Moment for en
uendelig Begrændsning og Nedbryden af alle Grændser:
det blivende Characteristicon er saaledes Endelighed
og Forgjængelighed.

Forandringen er den concrete Vorden, dog med den
Modification, at ligesom Vorden nærmest betegner en
Tilbliven, hvorved Tilværelsen \textbf{indledes}, saaledes udtrykker
Forandringen nærmest en Destruction, hvori
Tilværelsen \textbf{ophæves}.

Forandringens Bevægelse udtales i følgende Sætninger: a)
af Noget kommer Andet; b) af Intet (ɔ:
af det bestemte, tilværende Intet, nemlig Andet) kommer
Noget, og endelig c) af Intet kommer Intet; thi
i Forandringen er Noget og Andet hinandens gjensidige
Intet, og netop derfor et Intet hvert for sig.
I Forandringen ligger Tilværelsens uendelige Tilintetgørelse. Det er ingen fremmed Magt, for hvilken Tilværelsen her bukker under. Tilintetgørelsen er dens eget
allerinderste Væsen. Det er derfor saa langt fra, at
den sande Væren skulde gaae tilgrunde i Tilværelsens Intet, at den meget mere bereder sig sin Seier af
selve Alterationen og Forkrænkeligheden. Hvad der
gjælder om ethvert enkelt Moment i Tilværelsen, at
det nemlig gaaer over i sit Andet, maa nødvendigviis
ogsaa gjælde om den hele Tilværelse selv; i sin uendelige
Alteration henviser denne idelig til sit sande Intet. Men dette Intet kan i dets positive Bestemthed
ikke sættes som den fragmentariske Væren-for-Andet;
thi i denne Form er det netop Tilværelsen selv.
Hvis Tilværelsen endte i et saadant Intet, gik den
just derved sammen med sig selv, beseirede hele sin 
An%
% --- p. 14 ---
\setcounter{footnote}{0}%
\opage{14}dethed, og blev det Modsatte af, hvad den forhen havde
været. Det Intet, hvori Tilværelsen forsvinder, er
netop det allergedigneste Noget, der har gjennemgaaet
og overvundet enhver Alteration (Væren-for-Andet),
altsaa den egentlige Væren-for-sig.

\underafsnit{c) For-sig-væren.}
\paragraf{13}{Eet, Flere, Exclusion.}

Naar Tilværen har gjennemgaaet og overvundet
enhver muelig Forandring, bestemmer den sig
som positiv Væren-for-sig, ɔ: som medieret
Qvalitet. Da Andetheden er forsvunden, ponerer
Noget sig selv som det uopløselige Eet, som det
logiske Atom. Alt Andet er saaledes kun en Repetition
af Eet; Andetheden og Modsætningen maae
derfor søges i Fleerheden. Det med sig uendelig
identiske Eet adsplitter sig selv i de \emph{Flere}. Men
idet Ethvert af de Flere selv er et Eet, der fortrænger
alle andre, medens de alle absorberes af Eenhedens
Eet, saa er Qvalitetens Selvmodsigelse herved
stegen til det Høieste. Den for-sig-værende Negation
yttrer sig derfor som en alsidig \emph{Exclusion},
der destruerer al qvalitativ Grændse.
 Resultatet af Tilværelsens Selvudvikling er den fuldkomne Identitet af Noget og Andet. Men en saadan Væren er ikke længere Tilværen; thi dennes Criterium var jo netop en Dualitet, en endelig Modsætning imellem begge de værende Led. Det vil heraf 
% --- p. 15 ---
\setcounter{footnote}{0}%
\opage{15}kunne indsees, at den logiske Ophævelse af givne Modsætninger
baade er disses Tilintetgørelse og Opbevarelse.
I Tilværen gjaldt det om at finde og sætte
Andetheden paa alle Punkter; i For-sig-væren derimod
gjælder det at overvinde den i alle Skikkelser,
og vise den som Moment i den afsluttede Selvaffirmation.
I Tilværen vare alle Bestemmelser udvortes
og relative, fordi Noget og Andet idelig gave Anviisning
paa hinanden; i For-sig-væren er Bestemtheden absolut; her ligger Negationen og Begrændsningen
ikke længere udenfor, men i selve det, som er
for sig. Men Tilværelsen med dens Momenter er desuagtet
opbevaret; thi a) det er af disse Momenter, at alle Bestemmelser i For-sig-væren dannes.
Skjøndt alle Forskjelligheders og Forandringers hele
Uendelighed her synker sammen i umiddelbar Eenhed,
saa er vor Væren-for-sig derfor ikke den rene Væren.
I den rene Væren vare Forskjellene endnu
ikke logisk satte; i For-sig-væren ere de satte som
overvundne. For at tænke For-sig-værens Identitet, maa man først tænke Tilværens Forskjel.
b) er saa langt fra, at For-sig-væren umiddelbart skulde kunne fortrænge Andetheden, at den meget
mere maae, for at identificere den med sig, hæve
den til et for-sig-værende Andet. Derved faaer
For-sig-væren selv Tilværen. c) Idet For-sig-væren
gjennem Tilværen gaaer sammen med sig selv, og
saaledes opløfter sig over Tilværelsen, da skeer det
alene derved, at alle Tilværelsens Momenter selv paa
ethvert Punkt, i enhver Form opløftes til F -
sig-væren. Alt dette vil af det Efterfølgende blive mere
indlysende.

I For-sig-væren maae vi nu for det Første ikke
søge „\textbf{Noget}“; thi da det for-sig-værende Noget
% --- p. 16 ---
\setcounter{footnote}{0}%
\opage{16}er gaaet sammen med sit Andet, fortjener det ikke
længer dette Navn; ligesom ogsaa Andet, der har
identificeret Noget med sig, mister Andethedens Characteer.
--- Endnu mindre kunne vi her begynde med
den ubestemte Artikels \textbf{Et}; thi vel er det sandt, at
„Et“ er Fælledsbestemmelse for Noget og Andet, men
paa denne Maade er Væren fælleds for alle logiske
Categorier. „Et“ er ganske ligegyldigt imod den Forskjel,
der ligger i Noget og Andet; have vi ikke \textbf{et}
Noget, saa have vi \textbf{et} Andet; begges ophævede Forskjel
udtrykkes derimod bestemt i Categorien „\textbf{Eet}“,
som den egentlige Identitetscategorie. Endeligheden,
Dualiteten er forsvunden, i det Noget og Andet gaae
i \textbf{Eet}. Den Eenhed, hvori Noget og Andet her ere
ophævede, maa betragtes fra en dobbelt Side. Deels
er det Noget, der ophæver Andet, og gjør sig til „Eet“
dermed; deels er det omvendt Andet, der gjør sig til „Eet“
med Noget. Vi have derfor strax et dobbelt Eet, nemlig
Eet (a), der begyndte fra Noget, og Eet (b), som
begyndte fra Andet. Men da al Forskjel imellem Noget
og Andet er overvunden, saa er det fra dette Synspunkt ganske ligegyldigt, om vi begynde med Noget
eller med Andet; begge de nævnte Led ere da i Sandhed Eet. Saaledes maa Eet (c) udtrykke begges Eenhed.
Havde der nu været mindste Modsætning imellem
Eet (a) og Eet (b), saa vilde Eet (c) derved have
vindiceret sig en egen Betydning som den høiere Eenhed,
der ophævede begge de foregaaende Led; men da
disse ere \textbf{for sig} identiske, staaer det sidste Led ganske
\textbf{for} sig, som Repetition af det, der allerede er givet.
De trende Eetere ere altsaa i og for sig identiske, hvad
der kan udtrykkes ved et nyt Led, Eet (d); saaledes
producerer Eet sig selv i det Uendelige.
% --- p. 17 ---
\setcounter{footnote}{0}%
\opage{17}\textbf{Eet} sættes her imod \textbf{Eet}; men Modsætningen
forsvinder som en Skygge; thi Eet er - Eet. Det
er ved sin uudtømmelige Selvfrembringelse, at Eet overvinder
alt \textbf{Andet}.

Den eneste Form, under hvilken Modsætningen
endnu er bevaret, maa da søges i Frembringelsen selv.
Det producerende Eet i dets For-sig-væren kan
kun modsættes de producerede Eeteres hele Indbegreb,
saaledes som dette er for sig: \textbf{Eet} kan kun modsættes
de Flere. Eenhed og Fleerhed ere da den Dualitet,
som Eet her indeslutter; vi skulle nu see, hvorledes
denne udvikles og overvindes.

\textbf{Eet} er Eenhedens umiddelbare og bestemte Udtryk;
men just fordi Eet selv er den fulde Eenhed, kommer
Fleerheden, skjøndt et Product af Eet, til at ligge udenfor
Eenheden. Men udenfor Eenheden er Eet tillige
udenfor sig selv. Vel kan man sige, at ethvert produceret
Eet er sin egen Eenhed; men i denne Form
fremtræder det da ikke længere som produceret i Modsætning
til det producerende Eet. Som Led i Fleerheden
er ethvert Eet begrændset af andre Eetere; i
Fleerheden sønderbryder Eet sin egen Væren-for-sig.
Men denne Eetets Udgang af sig selv i Fleerheden
er netop dets Tilbagegang til sig selv og Eenheden. At Eet er gaaet ud over sig selv, vil med andre
Ord sige, at det er blevet sin egen Grændse.
Udenfor Eenheden, hvor vi skulde have ventet noget
Andet, træffe vi Eet, blot i omvendt Skikkelse. Eenhedens
Eet, der er Fleerhedens Princip, udgjør selv
det første Led i Fleerhedens Række. Som Princip
har det hele Fleerheden i sig, som Rækkens første Led
er det selv eet af de Flere, og coordineret med et hvilketsomhelst
Andet. Ethvert Led i Fleerhedens Række
er følgelig selv svangert med en Fleerhed. De Flere
% --- p. 18 ---
\setcounter{footnote}{0}%
\opage{18}ere selve Explicationen af Eet, og Eet Concentrationen
af de Flere. Saaledes fremstiller sig Eenhedens
Eet i dobbelt Form, nemlig som Fleerhedens Princip,
og som det, der selv resulterer af Fleerheden. Derved
tager Eenheden sig selv overalt tilbage af sin egen
Fleerhed og integrerer sin For-sig-væren. Det samme gjælder nu ogsaa om Fleerheden. Da
den sande Eenhed resulterer af Fleerheden, saa maae vi
her begynde med denne. De flere Eetere i deres simultane Existents gaae altsaa forud for Eenheden: deres
sporadiske Characteer vidner om Fleerhedens Væren-for-sig. Men denne For-sig-væren er endnu kun
abstrakt. Fleerheden støtter sig til enkelte Eetere,
der hvert \textbf{for} sig ere Eenheder. I ethvert isoleret
Eet gaaer Fleerheden under: det enkelte Eet er, som
saadant, Fleerhedens Grændse. Men denne Grændse
ophæver Fleerheden selv paa dobbelt Maade: a) Idet
Fleerheden forbyder det enkelte Eet at blive indenfor
dets egen Eenhed, saa retter den dets Bestemmelse
mod et andet Eet, og medens alle Eeterne udvortes
refereres til hverandre, kommer deres Eenhed derved udenfor
dem alle og udenfor sig selv; men denne Halvering
af Eenheden, der skeer saaledes, at dens ene Moment ligger
adspredt som Middelpunkt for de enkelte Eetere,
medens det andet Moment svæver over dem alle, er
netop Fleerheden. Fremdeles b) det Eenhedsmoment,
der overalt udgjør Centrum for de enkelte Eeteres Væren-i-sig, er, ifølge det Foregaaende, Principet for
Fleerhed. At Fleerheden saaledes fra alle Sider danner
Eenheden og resulterer af denne, udgjør dens fuldkomne
For - sig væren. Men \textbf{Eetet} selv synes herved givet til Priis for
den meest afgjørende Modsigelse. Det skulde være den
sande gedigne Identitet, og dog frembryde tvende mod%
% --- p. 19 ---
\setcounter{footnote}{0}%
\opage{19}satte Elementer i afsluttede Kredse af dets Midte. Dets Eenhed har sin Væren-for-sig lige over for
den for-sig-værende Fleerhed. Det er imidlertid saa
langt fra, at denne Dualitet skulde kunne sprænge
det adamantiske \textbf{Eet}, at dette just gjennem Modsætningens
Tilværelse hæver sig til den høieste Væren-for-sig. Den Fleerhed, gjennem hvilken Eenheden
i Eet har medieret sin For-sig-væren, er den
selvsamme, der igjen medieres igjennem selve Eenheden.
Eenhedens Selvmediation falder saaledes indenfor
Fleerhedens og Fleerhedens indenfor Eenhedens Grændser.
Eenheden opnaaer først gjennem Fleerhedens For-sig-væren
sin rette Væren-for-sig. Resultatet er
altsaa, at For-sig-væren udvikler sig i en Række
af for-sig-værende Momenter: den ene For-sig-væren
er kun for den anden; denne uendelige Selvfordobling er det Middel, hvorved Andetheden her beseires.

I dette Resultat culminerer Qvalitetens Udvikling.
Qvaliteten er For-sig-væren, sat i Umiddelbarhedens
Form. De Modsætningsmomenter, den rene
Væren absorberede, gives herved tilbage. Værens Eenhed
er Principet for uendelig Specification. Gjennemførelsen
af denne Specification danner et atomistisk
Element i Logiken, som Dialektiken aldrig kan opløse
uden selv at forsvinde. Enhver Categorie er et logisk
Atom, hvis Væren for de Andre uafbrudt medieres
gjennem dets Væren-for-sig.\footnote{Den her omhandlede Categorie spiller en stor Rolle i Theologien, hvis Opgave det netop er at mediere Guds Forsig-væren med Verdens Væren-for-sig. De hegelske Disciple synes imidlertid ikke at unde den, hvor mesterlig end Mesteren selv paa sin Maade har udviklet den. Grunden dertil er forklarlig. Man har fornøiet sig saa kostelig over at see, hvorledes det Ene ved Dialektikens Tryllemagt idelig kunde slaae over i det Andet, saa at Alt i den lovpriste Aandsphilosophie tilsidst er blevet en Væren for Andet. — Principet for den gjennemførte mythiske Behandling af Jesu Liv er Ideen om Selvbevidsthedens egen Christusfrembringelse. Christus, Christusbevidstheden, skal da være for Menneskebevidstheden (ɔ: for Andet) uden egentlig at være for sig. — Saaledes bliver jo Alt ganske flydende; Qvaliteterne i dette speculative Fluidum have ingen Afsmag; thi der ere ingen. Det er den pure Philosophie, luttret fra alle Forestillingselementer; dens rene Væren er renere end det rene Vand.}
% --- p. 20 ---
\setcounter{footnote}{0}%
\opage{20}Denne positive Form af For-sig-væren udtrykker den fuldkomne Umiddelbarhed. Ethvert Moment maa tages \textbf{for} sig; thi det er, som det er. Fra
en saa fuldstændig Positivitet synes derfor Intetheden,
der i Tilværen forstak sig under Andethedens Maske,
at være aldeles fortrængt. Hvor skulde den nu kunne
indtrænge sig, alle Pladser ere jo besatte med gedigne
Eetere og affirmative Qvaliteter? Men er Negationen ikke i den positive For-sigværen
selv, maa den have en Væren-for-sig. Denne
har den under vor Udvikling hidtil havt i den betragtende
Bevidsthed.

Naar der spurgtes om \textbf{Grændsen for} Eet, saa
spurgtes der egentlig om dets \textbf{Negation}; men istedetfor
Selvnegeringens negative Operation, fremstillede
sig strax Negationens positive Udtryk, ɔ: de Flere. Da
nu de positive Momenter saaledes overalt ere satte
Negationens Sted, saa ere de gjensidig hverandres 
Ne%
% --- p. 21 ---
\setcounter{footnote}{0}%
\opage{21}gation; For-sig-væren maa derfor udvikles gjennem
alle dens negative Forhold og Bevægelser.

Eet er, som vi have seet, Identiteten af Noget og
Andet. Denne Identitet har en aldeles negativ Characteer,
thi den fremkommer ved Forandringens Forandring.
I Forandringen virkede den værende Negation
derved, at Noget og Andet vare for hinanden Andet
og Intet; det ene Intet alternerede saaledes med det
andet. Denne Alternation bar et udvortes Præg; thi
uagtet Noget blev til Intet, hver Gang det gik over
i Andet, saa var det derfor endnu ligefuldt et positivt
Noget i sig. Negationen var saaledes blot relativ: Intet havde Tilværen, men endnu ingen Væren-forsig.
I den gjennemførte Forandring viser det sig derimod, at Noget ikke kunde blive til Andet, ɔ: til sit
Intet, hvis det ikke allerede i sig var et Andet. Men
idet Andet gjøres til Noget, bliver det sig selv Andet
og derved Intet i sig. Eenheden af Noget og Andet
er saaledes at sætte som det fordoblede Intet i
dettes fuldkomneste Liighed med sig selv: det positive
Eet er af negativ Herkomst. Vi kunne nu concentrere
alle disse Negationsmomenter paa følgende Maade:

Noget er sit Andet derved, at det selv er sit Intet;
thi 1) Noget er = Noget; heri ligger dets Væren (a):
det er sig selv og ikke Andet. 2) Noget er = Andet:
den Væren, Noget har, ligger ikke i det selv, men i
Andet. I sit Andet erholder det sin positive Væren (b).
Ifølge b maa a være Nogets første Negation, Intet (a): Noget \textbf{er ikke} hvad det er; det er i sig Intet,
derfor griber det ud over sig selv, for i b at \textbf{være}, hvad det ikke er. Da b udtrykker, hvad Noget ikke
er, bliver det igjen negeret af a og derved nedsat til
et Intet b. For at den dobbelte Væren i Noget ikke
længer skal være behæftet med Tilværelsens relative
% --- p. 22 ---
\setcounter{footnote}{0}%
\opage{22}Modsætning, men fuldkomnere gaae i \textbf{Eet}, maa hvert
af Leddene selv være sit Intet; uden denne vexelsidige
Selvnegation kunne de ei blive positivt identiske. Saaledes forvandler den i Noget optagne Andethed sig til
en gjennemtrængende Negation. Idet Eet bliver sat,
faaer det paa een Gang en positiv og negativ Side. Den positive Side fremkommer derved, at Væren (a)
identificeres med Væren (b); men Identificationen gaaer
gjennem Negationen; Intet (a) identificeres med Intet (b).

Heraf bliver det nu klart, at Eet er sin egen Negation. Den negative Side repræsenterer os her den
retorqverede Andethed. Udenfor Eet ligger altsaa ikke
noget Intet; thi den for-sig-værende Negation bevæger
sig i og igjennem Eet, endnu mindre nogetsomhelst
Andet; thi hele Andetheden er indsluttet i den negative
Bevægelse. Eet er gaaet i sig selv og har taget Alt
med sig: et \textbf{Udenfor} gjælder ikke længere. Men denne Identitet af den positive og negative
Side fremkalder Modsigelsens Uro i \textbf{Eetet} selv. De Momenter, der constituere den positive Side, ere de
selvsamme som de, der bestemme den negative. Her
er ikke en Position + en Negation, men en Position,
som selv er Negation, og omvendt. Kalde vi
f. Ex. den positive Side A, den negative B, da kan
Eet kun holde sig i Momentet A derved, at det
optager den negative Bevægelse af B, og gjør sig selv
til det, hvad B er, men B er netop gjennemførte
Negation. Saaledes bliver A det Negative af sig selv;
Positiviteten falder over paa B, der igjen absorberer de
negative Momenter i A. Under A's og B's fortsatte
Rotation vibrerer Eet uafbrudt imellem sig selv og sit
Intet.
% --- p. 23 ---
\setcounter{footnote}{0}%
\opage{23}Men dette Intet er selv et Eet, der gjør sit forudgangne Eet til Intet. Derfor bliver den for-sig-værende Negation at characterisere som \textbf{Exclusion}. I Exclusionen, den skarpeste Qvalitetsgrændse, ligger
netop den hele Haardhed, hvormed den afsluttede
Positivitet adskiller sig fra den afsluttede Negativitet.
Dette viser sig, naar vi forfølge den gjennem alle
dens Bevægelsesmomenter, som a) den begyndende,
b) den vexelsidige, c) den absolute Exclusion. Den begyndende Exclusion er Eetets egen Tilbli%
velse. Eet sætter sig i Momentet A, derved excluderer
det Positiviteten i B; B vindicerer nu sin Positivitet
og excluderer det positive A. Under dette dialectiske
Pulsslag maa Eet bevæge sig, endskjøndt det staaer
stille. Det bevæger sig; thi neppe har det fixeret sin
Positivitet, før denne strax nedbrydes af Negationen.
Men det staaer dog stille; thi der er og bliver kun
eet eneste Eet; naar det ene Eet kommer, er det andet
excluderet. Derved træder Eet i en Relation til
sig selv, som ingen Relation er, og dog en Relation;
thi det Eet, som opstaaer, forudsætter stedse det, som
forgaaer: Exclusionen bliver vexelsidig.

Ethvert Eet har altsaa den Egenskab at ponere sig
selv og excludere et andet Eet; men dette ponerer igjen
sig selv, og excluderer et andet o. s. f. Derfor existerer
Eet kun i Fleerheden. I Fleerheden dele alle de
enkelte Eetere sig imellem Selvvedligeholdelse og Exclusion
af de øvrige; disse tvende Virksomheder betinge
hinanden. Men Fleerhedens Ligevægt bliver idelig ophævet;
thi hvad det ene gjør, det gjøre de alle. Deducerer det ene Eet sin Existents af sig selv, da gjøre
de øvrige det samme: alle Eeterne i Fleerheden ere
ved eet Slag. Kan det ene ei existere uden ved at
excludere alle øvrige, saa kunne disse heller ikke kom%
% --- p. 24 ---
\setcounter{footnote}{0}%
\opage{24}me til Existents, før de have fuldbyrdet Exclusionen af
det ene og af hverandre indbyrdes. Dette er netop
det Desperate i den for-sig-værende Negation, at den
ikke kjender nogen Indskrænkning. I logisk Fortvivlelse
adsplitter Eet sig selv, for at redde sin Identitet
i Fleerheden; men Fleerheden hviler i Vedligeholdelsen
af Eet og gaaer under i dettes Opløsning. Herved
sprænges alle qvalitative Skranker.

Det Overraskende i denne Modsigelse hidrører fra
en resterende Umiddelbarhed. Man slutter ifølge principium exclusi medii: enten ere de enkelte Eetere excluderede,
eller de ikke ere det. Men et saadant Dilemma kan ikke holde sig, naar Exclusionen bliver absolut. Ifølge det Foregaaende ere nemlig Eeterne satte
i deres Enkelthed ved selve Exclusionen; naar denne
nu fuldbyrdes paa det enkelte Eet, da skeer dette derved
sin Ret. Alle Enkeltes Exclusion er Alles Position;
ethvert Eet sætter sig da først i og \textbf{for} sig selv,
naar det sætter sig i og \textbf{for} alle de øvrige: den absolute
Exclusion er en Assimilation.

\overgang

Naar det ovenfor i den positive Udvikling hed, at
Eenheden gjennem Fleerhedens For-sig-væren erholdt
sin fuldkomne Væren-for-sig, da var Hovedtanken
den, at den ene positive Modsætning skulde integrere
den anden. Derfor fremtraadte den afrundede Umiddelbarhed
allevegne, og Alt fremstillede sig i Qvalitetens
Form.

Nu derimod, da den gjennemførte Negation har ophævet
den hele qvalitative Umiddelbarhed, er den tilbagetrængte
og skjulte Andethed vedbørlig fremdragen,
og For-sig-væren \textbf{ud af sig selv} bleven en Væren-for-Andet. Fleerheden erkjendes igjen som Eenhedens
% --- p. 25 ---
\setcounter{footnote}{0}%
\opage{25}Andet; men Eet er ikke længere adsplittet i de Flere,
thi dets Væren-i-sig er tillige dets Væren udenfor
sig. Derved, at dette \textbf{Udenfor} er givet tilbage, forsætter
Eenheden sig selv ved en bøielig Overgang i
Andet, producerer sin Andethed allevegne og assimilerer
den med sig. Alle Specificationer blive nu at ansee
som Udvidelser og Indskrænkninger af Eet og det
Selvsamme. Dette Eet er ikke længere qvalitativt, thi
det har ingen Væren-for-sig udenfor den specificerende
Bevægelse. Selvudvidelse og Selvindskrænkning ere
dets egne immanente Bestemmelser. Saaledes er nu
den qvalitative For-sig-værens stive, faste Grændse
bleven elastisk og bevægelig; den bestemmer sig derfor
som

\afsnit{B. Qvantitet.}
\underafsnit{a) Den rene Qvantitet.}
\paragraf{14}{Continuitet, Discretion, Limitation.}

\emph{Qvantiteten}, Eenheden af Væren-for-sig
og Væren-for-Andet, er den egentlige Grændse.
I det Grændsen udvikler sin Identitet med det Tilgrændsende,
bliver den begrændset af sig selv, og er denne sin Selvudvidelse bestemt som \emph{Continuitet}. Men Grændsens Midte er den exclusive Negation
af alt Tilgrændsende og derved sin egen negative
Eenhed; som saadan falder Grændsen fra sig selv i
lutter exclusive Grændsemomenter, og er Discretion. \emph{Limitationen}, Continuitetens og Dis ---
% --- p. 26 ---
\setcounter{footnote}{0}%
\opage{26}cretionenS satte Eenhed, udgjør Grændsens egentlige
Tilværen eller dens fulde Selvbegrændsning.

Tilværen var kun en Væren for Andet, derfor gik
den igjennem Forandringen over i For-sig-væren. Ved at absorbere Andetheden optog For-sig-væren
sin exclusive Negation i sig, og maatte derfor gaae ud
over sig selv; Qvantiteten, der er begges høiere Eenhed,
skal nu udvikle Tilværens høiere Andet af For-sig-værens Eet, og maa derfor bestemmes som begges
Overgang i hinanden, d. e. som \textbf{Grændsen}. Den
umiddelbare eller qvalitative Grændse er en negativ
Midte imellem tvende positive Modsætninger: enhver
Qvalitet har sin Grændse. I Qvaliteten er Grændsen,
som den egentlige Midte imellem Noget og Andet, den rene i sig værende Negation. Den negative
Side af Grændsen derimod, der vender imod Noget,
bliver strax assimileret med det positive Noget selv,
ligesom ogsaa den imod Andet vendte Side gaaer op i
dettes Positivitet. Qvalitetens Grændse er derfor positiv
Bestemthed. Bestemtheden er Negation, men en
Negation, der opgaaer som svindende Moment i Positionen: Qvaliteten er en For-sig-væren, der maa tages positivt, som den er. Men er Grændsen, der adskiller
tvende Qvaliteter, som saadan hverken den ene
eller den anden af disse, da maae begge i selve Grændsen
betragtes som fraværende, og det ikke blot fra hinanden
indbyrdes, men ogsaa fra den egentlige Grændselinies Midte. De grændse da ikke umiddelbart op til
hinanden, men til den fælleds Grændse. Vi have altsaa:
Tilgrændsende, Grændse, Tilgrændsende\dots{}. i det
Uendelige. --- Men ligesom Væren (§ 11) gik over i Intet, fordi den, som saadan, havde Intet i sig, og omvendt,
saaledes gaaer Grændsen og det Tilgrændsende
% --- p. 27 ---
\setcounter{footnote}{0}%
\opage{27}her idelig over i hinanden. Kalde vi Grændsens Væren-for-sig a, dens Væren for det Tilgrændsende b, det Tilgrændsendes Væren for Grændsen α,
dets Væren for sig β, da viser det sig, at det er
Grændsen, der i det Tilgrændsende idelig grændser op
til sig selv. Da b nemlig ligger imellem Grændsen
og det Tilgrændsende, saa bliver det derved selv at betragte som den egentlige negative og exclusive Grændse,
for hvilken a og α ere tilgrændsende Momenter. Men
just fordi b nu er den for-sig-værende Grændse, bliver
α Grændsenegationen imellem b og β. Paa denne
Maade vedbliver Grændsen at forflytte sig i uendelig
Progres. Dette er da characteristisk for Grændsen,
at den paa exclusiv Maade stedse er baade i og udenfor
sig selv. Grændsemediationen gjør det Tilgrændsende
\textbf{til} Grændse i samme Nu, som den gjør Grændsen
til et Tilgrændsende. I denne Grændsens Selvudvidelse sættes Qvantiteten som \textbf{Continuitet}.

I Continuiteten viser det sig, at Grændsen ikke kan
være noget Stillestaaende. Hver Gang Grændsen sættes
med en vis umiddelbar Bestemthed, omgives den
af et Tilgrændsende; men dette er, som vi have seet,
selv Grændse, og følgelig Grændsen selv. Altsaa: kun
derved, at Grændsen gjennemgaaer alle sine Begrændsningsmomenter, sætter og ophæver alle sine Bestemmelser,
kan den fattes som sand Grændse. Men i
Grændsens Selvforflyttelse ligger tillige en uendelig
Gjentagelse af dens exclusive Negation. Hvad der
først bestemmes som den rene for-sig-værende Grændse,
bliver i næste Moment et Tilgrændsende; men ethvert
Tilgrændsende er ligeledes en for-sig-værende Grændse:
den ene Grændse excluderer derved den anden; men da
hver Grændse bestemmes paa lige Maade, saa er
Grændse stedse = Grændse; Grændsen excluderer alt%
% --- p. 28 ---
\setcounter{footnote}{0}%
\opage{28}saa sig selv, sætter sig derved som adskilt fra sig selv:
i denne fuldendte Selvadskillelse er Grændsen Discretion. Den rene Grændse opløser sig da i en Fleerhed
af Grændser, saaledes, at enhver igjen for sig indeholder
en Grændsefleerhed. I Grændsen er paa
een Gang saavel Continuitet som Discretion, eller, nøiagtigere udtrykt, Grændsen er selv Continuitetens og
Discretionens sande Midte. Det er denne skarpe Adskillelse
af begge Momenter, der ligger til Grund for
den eleatiske Fremstilling af Bevægelsens logiske Umuelighed. Man kan nemlig i Continuitetens Linie ponere
tvende hvilkesomhelst discrete Punkter, f. Ex.
A og B i Linien \textbf{A---c---B}; Linien selv er ganske ligegyldig
imod disse; thi den kan forlænges ud over
begge: Punkternes Discretion er paa udvortes Maade poneret i Liniens flydende Continuitet.

Spørgsmaalet om Bevægelsens Muelighed reduceres
nu til Spørgsmaalet om, hvad der betinger dens Resultat. Før Den, der bevæger sig, kan komme fra
A til B, maa han først være kommen den halve Vei,
nemlig fra A til c. Men Afstanden fra A til c kan
igjen halveres, og Halveringen af de halverede Afstande
fortsættes uden Ophør. Bevægelsens Absurditet
ligger da deri, at selv den korteste Vei ei kan tilbagelægges,
før en endnu kortere er gjennemløben: Bevægelsen
forudsætter sig selv i det Uendelige. Resultatet
af den allerførste Bevægelse skulde altsaa være: 1) en
\textbf{bestemt} Continuitet (ɔ: Continuitet med en endelig
lille Udstrækning); thi uden en vis Afstand imellem
Begyndelsespunkt og Endepunkt kan Bevægelsen
ikke tænkes; 2) en \textbf{udelelig} Continuitet (med en uendelig
lille Udstrækning); thi i modsat Fald vilde Bevægelsen
ei være den første. Modsigelsen ligger altsaa
deri, at man fordrer en Continuitet uden Discretion,
% --- p. 29 ---
\setcounter{footnote}{0}%
\opage{29}medens man dog gaaer ud fra den Betragtning, at
ethvert Continuum kan adskilles i discrete Momenter.
Derved kommer man til at tænke sig Bevægelsen udenfor
Bevægelsen selv. Forud for det allerførste Udgangspunkt
kan nemlig ingen Bevægelse være given, efterdi
det er det første; men naar Bevægelsen staaer stille i
det allerførste Punkt, saa kan den jo aldrig naae det
andet. Saaledes kommer man da ogsaa til at tænke
sig Hvilen udenfor Hvilen selv. Enhver Hvile er Resultat
af en foregaaende Bevægelse. Bevægelsen kan
nu ikke komme til Hvile i det andet Punkt, fordi den
bestandig hviler i det første; men den kan heller ikke
hvile i det første, fordi der forud for dette ingen Bevægelse er gaaet.

Sandheden er, at den fuldførte Bevægelse indeholder
Continuitet og Discretion i uopløselig Eenhed.
Eleaterne have blot indskrænket sig til at eftervise begge
Modsætninger, og istedetfor at erkjende den enes Overgang i den anden, for saaledes at begribe Bevægelsen
af Bevægelsen selv, have de urgeret Continuiteten og
Discretionen hver for sig, og reduceret Bevægelsen til
en discret og følgelig absurd Stillestaaen. --- Kun Bevægelsen kan ophæve den skarpe Modsætning imellem
Continuitet og Discretion. Ligesom Vorden var en
uafbrudt Overgang fra Væren til Intet, og derved
paa een Gang affirmerede og negerede begge Momenter, saaledes er nu Bevægelsen en bestandig Svæven
mellem Stillestaaen (Discretionens Moment) og Selvforflyttelse (Continuitetsmomentet), og følgelig det Medium,
hvori begge sættes.

Resultatet af Bevægelsen er altsaa en ved Discretionen
bestemt Continuitet, begrændset af en ved Continuiteten
bestemt Discretion eller, med andre Ord, Qvantitetens Tilværelse som Qvantum i continuerlig
% --- p. 30 ---
\setcounter{footnote}{0}%
\opage{30}og discret Størrelse. Den discrete Størrelse er, da
den har Continuiteten i sig, en Fleerhed af særskilte Continua, ligesom den continuerlige er en Eenhed af
flere discrete Momenter. Forsøger man nu at forvandle
den discrete Størrelse som saadan til continuerlig,
eller omvendt denne til hiin, saa kommer man til
den uendelige Deelbarheds Antinomie.

a) Den continuerlige Størrelse kan deles; thi den
har Discretionens Moment i sig, og hver af de discrete
Dele ere selv Continua. Ved Delingen opløses
altsaa den continuerlige Størrelse i en Fleerhed af continuerlige Størrelser, der igjen hver for sig kunne deles:
den continuerlige Størrelse er \textbf{uendelig deelbar}.

b) Da Resultatet af enhver Deling er en continuerlig Størrelse, saa er det saa langt fra, at Delingens
Fortsættelse skulde kunne ophæve Størrelsens Continuitet,
at den meget mere bliver en gjentagen Bekræftelse
af samme: den continuerlige Størrelse er følgelig absolut
\textbf{udelelig}.

Denne Modsigelse pleier man nu at tilhylle derved,
at man fremstiller enhver Størrelse som bestaaende af
uendelig mange, uendelig smaae Dele. En saadan Forklaring
viser en uendelig skjøn mathematisk Balance,
hvori den ene Tankeløshed opveier den anden. Man skal for det Første tænke sig, at en uendelig Deling fuldføres, det vil sige, faaer en Ende; dernæst, at de
Dele, der udbringes ved den uendelige Deling, ere
uendelig smaae (= 0), d. e. \textbf{ingen} Dele.

Det bliver heraf indlysende, at Deelbarheden og
Udeelbarheden gjensidig betinge hinanden. Delene ere
selv Størrelser og Størrelserne Dele. Derved, at den
ene Størrelse begrændses af den anden, er den ved
dens Dele begrændset i sig, og saaledes er Limitationen
eller den uendelige Begrændsning realiseret.
% --- p. 31 ---
\setcounter{footnote}{0}%
\opage{31}Limitationens Elementer ere lutter Størrelser og
lutter Dele. Hver Deel er som saadan en Enkelthed;
hver Størrelse derimod et Sammensat. Men da Delen
selv er Størrelse, og Størrelsen Deel, saa gaaer
det Sammensatte over i det Enkelte, og dette i hiint.
Hvert Enkelt maa som saadant sættes i dets Positivitet,
ligesom ethvert af Eeterne i For-sig-væren. Men
her sættes det enkelte. Eet med den udtrykkelige Bestemmelse, at det skal begrændses ved et andet enkelt % PRINTED AS IS: enkelte. (a stop in mid-sentence; S032)
Eet. Det gjælder da ikke her, som forhen, om Eenhedens
Tilbagegang til sig selv gjennem Fleerheden,
men om en \textbf{bestemt} Fleerhed og een ved Fleerheden begrændset Eenhed. Den limiterende Eenhed og Fleerhed
bestemmes ved en tredobbelt Act: 1) ethvert af
Fleerhedens Led adskilles og sættes for sig; 2) de ponerede
Led sammenfattes, hvorved Størrelsen, som omgrændsende
Eenhed, bliver sammensat; 3) den sammenfattende
Eenhed sættes igjen for sig som en ny
Enkelthed. I denne Adskillen og Sammenfatten er
Limitationen Numeration.

\underafsnit{b) Tallet.}
\paragraf{15}{Taleenhed, Antal, Utallighed.}

I Limitationen er Continuitetens Eenhed alsidig begrændset af Discretionens Fleerhed; begges
gjensidige Bestemmelse er at sætte som Tal. I Tallets Continuitet er den discrete Fleerhed paa een
Gang absorberet og excluderet; med denne Bestemmelse
er Continuiteten \emph{Taleenhed}, og Fleerhedens
Discretion \emph{Antal}. Da hvert af Antallets Eetere
% --- p. 32 ---
\setcounter{footnote}{0}%
\opage{32}er en Taleenhed, der indeholder sit Antal, saa bevæger
Tallet sig uendelig i sig selv, og gaaer dog
overalt udenfor sig selv: \emph{Utalligheden} er Tallets Selvopløsning.

Fra den logiske Udviklings Begyndelse af har der
foresvævet Tanken en Numerationsbestemmelse; thi uden
en saadan kan ingen Gjentagelse, ingen Modsætning
fastholdes. Tydelig yttrede sig Numerationens Moment
i Tilværens \textbf{Andet} (secundum), endnu tydeligere
i For-sig-værens \textbf{Eet} og \textbf{Flere}. Men med categorisk Accent kunde Numerationen ikke fremtræde i Qvaliteten.
Ved Bestemmelsen af Andet gjaldt det egentlig
om at vise, hvad det var i sig, som et Værende
(alterum, aliud), i Modsætning til Noget; dets Bestemmelse som et secundum var derimod ganske udvortes
og flygtig. Selv i For-sig-væren handledes der
egentlig kun om Eeternes qvalitative Genesis. For-sig-værens
Eet var derfor nærmest at betragte som et
Identitetens Eet; selv Fleerheden var qvalitativ, thi
den blev blot bestemt som den for-sig-værende Eenheds
ophævede Andet; derfor spurgtes der endnu ikke:
\textbf{hvormange}? Nu derimod, da Eenheden er bleven en
Continuitet, der har optaget Fleerheden som sin egen
Discretion i sig, sætter Qvantiteten sig som en gjennem
sine discrete Eetere fuldkommen bestemt Eenhed.
Herved kommer den logiske Tankebevægelse den ulogiske
Umiddelbarhed saa nær som mueligt; thi i Numerationen indskrænker den sig til en Efterviisning af de enkelte
Eetere, saaledes at hvert fastholdes for sig og
kun i udvortes Relation til de øvrige. Numerationens Resultat er Tallets Begreb. Rigtignok kan der i Virkeligheden
ingen Numeration finde Sted, uden at Talbestemmelsen foresvæver Bevidstheden, men under den
% --- p. 33 ---
\setcounter{footnote}{0}%
\opage{33}umiddelbare Tællen: 1+1+1 ꝛc., \textbf{er} ethvert af Eeterne
endnu nærmest opfattet som qvalitativ For-sig-væren; først ved Optællingens sammenfattende Operation fremkommer
det bestemte Begreb om en Eener, en Toer ꝛc.,
d. e. om \textbf{Tallet} som saadant. Den fuldkomneste
Værensbestemmelse, hvortil vi nu ere komne, maa da
sættes som en Talbestemmelse. Altilværelsen er \textbf{et}
Tal\footnote{Fra dette Synspunkt maa den pythagoræiske Philosophie bedømmes.}, thi den er en Eener, der i det Uendelige gjentager
sig selv, men ved hver Gjentagelse erholder en
ny og aldeles specifisk Betydning. Enhver Eener er
umiddelbart Eet med sig selv, men i denne Væren - forsig
har den sin særegne Plads i Eeternes Række. Det Characteristiske for Taleenheden er netop dette, at den
med al sin Varietet dog i ethvert Moment er saa absolut
begrændset. Forsaavidt Taleenheden betragtes som
saadan, f. Ex. en Femer modsat en Sexer, er den absolut
identisk med sig selv, og har excluderet al Fleerhed.
Men denne Exclusion er ligesaa meget en Ophævelse
og Optagelse af Fleerhedens Eetere; thi kun i
Relation til en begrændset Fleerhed kan Eenheden bestemmes.
Kun ved at sætte 1 + 1 + 1 + 1 + 1 erholde
vi Eenheden 5. Men i Fleerhedens Moment skal hvert
af Eeterne sættes for sig og følge en nødvendig bestemt
Orden; derved, at Fleerheden saaledes omgrændses
af Eenheden, bliver hvert af dens Eetere talt.
Disse Fleerhedens talte Eetere udgjøre nu Tallets anden
Hovedside, nemlig \textbf{Antal}. Antallet er netop
Grændsen imellem Eenhed og Fleerhed; thi medens
det paa den ene Side henpeger til de discrete Eetere,
saa udtrykker det paa den anden Side deres Sammen%
% --- p. 34 ---
\setcounter{footnote}{0}%
\opage{34}fattelse i Taleenheden. Antallet angiver stedse Eenhedens
og Fleerhedens vexelsidige Bestemmelse: saa stor
en Eenhed, saa stor en Fleerhed\footnote{Den eiendommelige Udvikling af Eenhed og Antal gives i Deductionen af de fire Regningsarter, men her gaae Talbestemmelserne saaledes over i Mathematikens egne Specificationer, at de ikke længere kunne blive Momenter i den speculative Logik. En concret Udvikling af Talforholdene i deres forskjellige Overgange vilde udkræve en gjennemført Mathematikens Philosophie. Medens nemlig Mathematiken umiddelbart forudsætter sine særegne Axiomata og ved stringente Forstandsdefinitioner characteriserer sine Problemer og retfærdiggjør Fremgangsmaaden ved Operationerne selv og deres Resultater, skulde Mathematikens Philosophie derimod deducere Operationernes og Resultaternes Eiendommelighed af selve Qvantitetens eget Begreb. Mathematiken erholder sin exacte Characteer ved den skarpe Begrændsning, hvormed den holder alle heterogene Elementer borte. Men netop den Omstændighed, at de mathematiske Resultater lade sig applicere paa den virkelige Tilværelse, viser, at Talbestemmelserne ere abstraherede af en concretere Verden. En Mathematikens Philosophie burde derfor eftervise den Eensidighed og Modsigelse, hvormed alle mathematiske Bestemmelser ere behæftede, og derved gjøre det indlysende, at denne Videnskab, saa selvstændig den end synes den umiddelbare Forstand at være, dog desuagtet dialectisk ophæver sig selv, for at blive Moment i en altomfattende Verdenserkjendelse.}.
I Antallet kommer derfor Tallets indre Modsigelse
klart tilsyne. Eenheden, Fleerheden og de talte Eetere
have nemlig i Antallet hver især deres abstracte Forsig-væren,
medens de dog kun ere for hverandre indbyrdes.
Taleenheden er i sin simpleste og almindeligste
Form given i Tallet 1. Men 1 er stedse = 1.
Denne Eenhedens tautologiske Liighed med sig selv
% --- p. 35 ---
\setcounter{footnote}{0}%
\opage{35}giør det nødvendigt, at man, for at \textbf{komme} til en Udvikling
af Fleerhedens Eetere, maa tage en anden Bestemmelse, der iovrigt, ifølge sin Natur, ligger ganske
udenfor Tallets Begreb, til Hjælp. Tanken om 1+1 kræver nemlig, at man skal adskille 1 fra 1. I Forestillingens
Sphære foregaaer denne Adskillelse særdeles
let, da man her i Almindelighed tæller selve de
tilværende Gjenstande. Men abstraherer man nu fra
alle specifiske Qvaliteter, saa kommer man til at sætte
1, 1, 1\dots{} Denne Sætten seer vel ved første Diekast % PRINTED AS IS: Diekast (D for Ø, Øiekast; S038)
ud som en Gjentagelse, men er det, ifølge det rene
Eenhedsbegreb, egentlig ikke. Tanken kan ikke adskille
det ene 1 fra det andet, hvis man ikke seer sig om
efter heterogene Criterier, f. Ex. Tiden og Successionen. Men Successionen ligger aldeles udenfor Tallets
Begreb. At 1 er et Tegn, der kan skrives paa flere for
skjellige Pladser, at det kan udtales ved at repetere
det samme Ord flere Gange, vedkommer ikke den abstracte
Eenhed som saadan. Naar altsaa en virkelig
Taladskillelse skal finde Sted, maa man have en vis
Qvalitet, ligegyldigt hvilken, in mente. Denne Qvalitet
er nu at ansee som Qvantitetens Andet, som den
Negation, gjennem hvilken den tager sig selv tilbage.
Hermed gaaer altsaa Tallet udenfor sig selv og reali:
serer saaledes Qvantitetens første Hovedfordring.

Men Qvaliteten selv erholder aldeles ingen qvalitativ
Bestemmelse ved selve Numerationen\footnote{Deraf kommer det ogsaa, at man ikke kan sammentælle forskjellige Qvaliteter, f. Ex. 4 Borde og 3 Stole. Ved en saadan Sammentællen vilde man nemlig sammenblande Qvantitetsbestemmelser med Qvalitetsbestemmelser, hvad der, ifølge Qvantitetens Natur, er absurd.}. Som det er
% --- p. 36 ---
\setcounter{footnote}{0}%
\opage{36}Tallet ganske ligegyldigt, enten vi tælle Stjernerne paa
Himlen eller Sandskornene ved Havets Bred, saa er
det ogsaa hver enkelt Stjerne ligegyldigt, hvad enten
den i Antallet sættes som den 1ste eller som den 10de. Paa denne Tallets Ligegyldighed imod alt Andet grunder sig Utalligheden af Tilværelsens Qvaliteter. Vel kunne Tingene i Verden tælles, men det er kun under
den Forudsætning, at de fremstille sig med udvortes,
tilfældig Indskrænkning og Begrændsning. Man kan
tælle Knapperne i en Kjole; men der ere i Verden aldrig
saamange Kjoler og aldrig saamange Knapper, der
kunde jo, hvad disse umiddelbare Qvaliteter selv angaaer, være mange flere. At Gud derimod har talt
Menneskenes Hovedhaar, er kun begribeligt derved, at
han, som sætter Tingenes Maal, tillige sætter den dybere
Eenhed af Qvalitet og Qvantitet. Skulde han
derimod først tage Haarene paa vore Hoveder som givne, og begynde at tælle 1, 2, 3, da vilde Numerationen
aldrig tilendebringes. Heraf indlyser det, at
Tallet negerer sin Andethed for at bestemme sig selv; Qvalitet og Qvantitet kunne ikke begrændse hinanden:
\textbf{Tallet tæller kun sig} \textbf{selv}. Saaledes sinde vi % PRINTED AS IS: sinde (long s for f, finde; S041)
Qvantitetens anden Hovedfordring, nemlig, at dens
Væren \textbf{udenfor} sig just skal være dens egen Væren i
sig, bestemt udtrykt.

Men i det Tallet excluderer sit Andet, giver det sig
selv sin Andethed. Det er denne Andethed, vi nu
nærmere skulle betragte.

Da hvert Tal er saa umiddelbart sig selv, at man
kun ved Indblanding af heterogene Bestemmelser kan
komme fra det ene Tal til det andet, saa holde de
enkelte Tal sig paa Grund heraf i abstract udvortes
Modsætning til hverandre. Tallets abstracte Væren i
sig er dets ligesaa abstracte Væren \textbf{udenfor} sig. Denne
% --- p. 37 ---
\setcounter{footnote}{0}%
\opage{37}Tallets Udvorteshed kjender ingen Grændse. Ethvert
Eet i Antallets Fleerhed kan sættes som Taleenhed for
et nyt Antal i det Uendelige: \textbf{Tallene ere utallige}.

Tallets \textbf{Utallighed} indeholder: a) Talbestemmelsens uendelige Reproduction; b) Talbestemmelsens
Forsvinden; c) Tallets logiske Ophævelse.

a) I de specifiske Taludtryk bliver Qvantiteten endeliggjort
og defineret ved Qvantum + Qvantum +\dots{}
De definitive Qvanta kunne tælles; men, hvorlænge
man end fortsætter Additionen og Multiplicationen,
faae de endelige Qvanta aldrig Ende. Qvantiteten,
der ophæver enhver Grændse, møder stedse nye Grændser,
og kommer saaledes aldrig til at fuldbyrde Grændsens
Ophævelse; den overskrider ethvert Antal og naaer
dog aldrig ud over Tallenes Endelighed. Utalligheden,
som det definitive Tals uophørlige Reproduction,
er herved sat som Grændsens Allestedsnærværelse i Qvantiteten, som den \textbf{uendelige Endelighed}.

b) Naar Qvantiteten udtrykkes i et Indbegreb af
endelige Qvanta, saa skulde man vente, at der i denne
fuldkomne Begrændsning var sat en bestemt Grændse
for et Maximum og Minimum; men det Største i Endeligheden har en endelig Storhed; der gives derfor
stedse et endnu Større; den mindste endelige Størrelse
har en Grændse for sin Lidenhed; udenfor denne Grændse
findes en endnu mindre Størrelse. Fremdeles ligger
der imellem tvende givne endelige Størrelser et uendeligt
Antal af andre endelige Størrelser. Imellem 1
og 3 ligger ikke blot 2×1, men 4×½, 8 ×¼ ꝛc.

Storhed og Lidenhed ere altsaa i den uendelige Endelighed aldeles relative Størrelser. Et Minimum og
Maximum kan ved udvortes Sammenligning her sættes
og ophæves overalt. Saaledes henflyder Tallenes Discretion i Continuitetens Strøm. Utalligheden, som
% --- p. 38 ---
\setcounter{footnote}{0}%
\opage{38}Tallets Opløsning, er en Ophævelse af alle ponerede
Grændser, \textbf{en af Endeligheden selv resulterende}
\textbf{Uendelighed}\footnote{Som et fra Mathematiken hentet Exempel paa den qvantitative Uendelighed i dens dobbelte Form, kan anføres Qvotienten af $\frac{1}{1-x}$, der (under Forudsætning: $x > 0$ og $< 1$) giver $1 + x + x^2 + \dots$. Qvotientens fortsatte Udvikling anskueliggør paa een Gang Forholdet imellem Antallet og Utalligheden og tillige den Modsigelse, der ligger i Tilnærmelsen til det uendeligt Store og uendeligt Lille. Ved hvilket Led man end afbryder Divisionen, viser denne Afbrydelse sig som vilkaarlig. Med hvert nyt Led erholde vi en forøget definitiv Talværdi: men den er dog stedse for liden; ved enhver Fortsættelse af Divisionen nærme vi os til et adæqvatere Udtryk: men den adæqvate Grændse selv flygter uophørlig længere bort. Den absolute Grændse kan derfor siges at være ligesaa langt borte fra $1 + x + x^2 + \dots x^{1000}$, som fra $1 + x$, hvorved hele Tilnærmelsen, betragtet fra det absolute Synspunkt, mister sin Betydning. Dette er, som Hegel ogsaa har udviklet det, Modsigelsen i den uendelige Progression: „Das Qvantum hat es in seinem Begriffe ein Jenseits zu haben. Dieß Jenseits ist erstlich das abstracte Moment des Nichtseyns des Qvantums: dieses löst sich an sich selbst auf; so bezieht es sich auf sein Jenseits als auf seine Unendlichkeit, nach dem qvalitativen Momente des Gegensatzes. Aber zweitens steht das Qvantum in Continuitet mit diesem Jenseits; das Qvantum besteht eben darin, das Andere seiner selbst, sich selbst äusserlich zu seyn, also ist dies Aeusserliche ebenso sehr nicht ein Anderes als das Qvantum: das Jenseits oder das Unendliche ist also selbst ein Qvantum. Das Jenseits ist auf diese Weise aus seiner Flucht zurückgerufen, und das Unendliche erreicht. Aber weil dieß zum Diesseits geworden wieder ein Qvantum ist, ist nun wieder eine neue Gränze gesetzt worden; diese, als Qvantum, ist auch wieder von sich selbst geflohen, ist als solches über sich hinaus, und hat sich in sein Nichtseyn, in sein Jenseits von sich selbst repellirt, das ebenso perennirend zum Qvantum wird, als dieses sich von sich selbst zum Jenseits abstößt“. (Wissensch. d. Logik. Werke 3 Band Pag. 265). Skal Qvotientudviklingen adæqvat afsluttes (som ved $x^n + \frac{x^{n+1}}{1-x}$), da er Progressionen herved reduceret til den i Forholdet mellem Dividendus og Divisor liggende Uendelighed, og Antallet sat i eet med Utalligheden. Paa denne Maade er den sande qvantitative Uendelighed givet i selve Udtrykket $\frac{1}{1-x}$. Da nemlig Dividendus og Divisor, hver for sig, ere bestemte endelige Tal, saa er deres indbyrdes Forhold her et bestemt, constant Forhold. Men dette i sig begrændsede Forhold er netop Kilden til den ubegrændsede Talstørrelse, der evolveres i Qvotienten. Udtrykket $\frac{1}{1-x}$ antager herved Characteren af en intensiv Størrelse, hvis Extensitet ligger i selve Qvotientens Evolution.}.
% --- p. 39 ---
\setcounter{footnote}{0}%
\opage{39}c) Under Mediationen af Antal og Utallighed bliver
Qvantiteten sin Grændse mægtig. Med Antallets
Utallighed synker den ned i lutter begrændsede Qvanta
og \textbf{sætter} derved \textbf{alle Grændser}. Det begrændsede Qvantum er endeligt, da det har en Grændse og stedse
viser ud over den, for at erholde sin Bestemthed i et
andet Qvantum. Men da denne Overgang i Andet
ligger i Qvantums egen Bestemmelse, da dette Andet
selv er Qvantum, saa er Begrændsningen tillige Grændsens Ophævelse. Den med sig selv medierede Endelighed
er da selve Uendeligheden. Gjennem denne Mediation
af endelige Qvanta tager Qvantiteten sig selv
% --- p. 40 ---
\setcounter{footnote}{0}%
\opage{40}tilbage, og \textbf{ophæver alle sine Grændser}. Forsaavidt
nu alle Begrændsninger ere satte og optagne i
Qvantiteten, er Tallenes Betydning bevaret; men derved,
at Grændsen allevegne bliver ophævet, er Tallet
nedsat til et svindende Moment. Istedetfor den umiddelbart begrændsede Taleenhed,
faae vi da en i sin Selvbegrændsning ubegrændset
Qvantitetseenhed ɔ: \textbf{en intensiv Størrelse}; og istedetfor
Antallet en Qvantitetsevolution, hvori utallige
Qvanta uafbrudt fixere og ophæve alle muelige Talbestemmelser,
ɔ: \textbf{en extensiv Størrelse}.

\underafsnit{c) Den absolute Qvantitet.}
\paragraf{16}{Intensitet, Grad, Approximation.}

Ved den uendelige Position og Ophævelse af Grændsen bevæger Qvantiteten sig i sig selv, og er
absolut Qvantitet. Naar Antallets Discretion sættes som ophævet Moment i Utallighedens Continuitet, bestemmes Qvantitetens Væren-i-sig, som
Intensitet. Ethvert af Intensitetens Momenter
har alle de øvrige ophævede til Eenhed med sig, og
er følgelig selv en Intensitet. Intensiteternes Relation
udtrykkes i \emph{Graden}. Graden er Qvantitetens qvalitative Bestemthed, men dog endnu abstract
qvantitativ. Den specifiske Begrændsning naaes kun
ved en \emph{Approximation}, hvis Utilstrækkelighed
nøder den i sig afsluttede Qvantitet til at gaae sammen
med sit qvalitative Andet i Modaliteten
% --- p. 41 ---
\setcounter{footnote}{0}%
\opage{41}. Categoriernes dialectiske Udvikling kunne vi adskille:
a) den successive Afvexling af de modsatte Led,
der danner en fortløbende Række, medens den, saa at
sige, skrider frem i en ret Linie; og b) den ene Modsætnings
nødvendige Fremgaaen af den anden, hvorved
begges speculative Eenhed erkjendes, og Bevægelsen
antager Characteren af en Circulation. Saaledes kan
Rækken: Væren, Intet, Væren\dots{} fortsættes uden Ophør; men i Vorden erkjendes det, at Intet fremgaaer
af Væren selv og Væren af Intet. Har man altsaa
giennemskuet det constante Forhold imellem Rækkens
tvende første Led, saa er dermed Betydningen af alle
de efterfølgende begrebet; thi, hvor ofte de end gjentages, saa ere de dog kun en tautologisk Repetition
af begge de første. Det Problem, den retliniede Bevægelse
fremkalder, bliver stedse løst i Cirkelbevægelsen.
Saalænge Qvantiteten blot opfattes under den udvortes
Modsætning af Grændse og Grændseløshed, af Antal
og Utallighed, indeholder denne Categorie en uopløselig
Modsigelse. Man fastholder da den Fordring, at Qvantiteten skal bestemmes ved discrete Qvanta, og at disse
maae angives i Tal. Men just fordi man er hildet i
Antallets Endelighed, bliver man ligesaa umiddelbart
overvunden af Utalligheden og Uendeligheden. Bestemt
udtalt lyder Modsigelsen saaledes:

: Tallet har Betydning; Tingene maae tælles, hvis
de skulle fattes; men nu kunne de ikke tælles: saa vist
som Rummet er uendeligt, saa vist ere Qvantitetens
\emph{Qvanta} utallige\footnote{Cfr. den indiske Mythe om Brima.}.

Naar man derimod ophæver Rækken og tydeliggjør
sig det constante Forhold imellem det bestemte Qvan%
% --- p. 42 ---
\setcounter{footnote}{0}%
\opage{42}tum og den ubestemte Qvantitet, imellem Antal og
Utallighed, da ville de abstracte Modsætninger ophæve
hinanden og Modsigelsen forsvinde. Derved, at det endelige
Qvantum begrændser sig selv, reduceres Qvantitetsbevægelsen
til en Cirkelbevægelse. Denne fuldender
nu Begrændsningen; thi det ene endelige Qvantum
fører idelig over i det andet Qvantum (a) over i
Qvantum (b), ligesom Væren over i Intet men enhver
Skranke er tillige hævet; thi det endelige Qvantum gjenfinder sig selv i sin Negation [Qvant. (b) fører igjen over i Qvant. (a)]. Qvantiteten afslutter
sig herved som absolut \textbf{Qvantitet}. Derfor er det
nu heller ikke fornødent at skride fra Antal til Antal,
for at tælle det Utallige. Det ene Antal begrændses
ved det andet; men Antal er = Antal; \textbf{Antallet er
sit eget} \textbf{Andet}, det gaaer sammen med sig selv: deri
ligger den medierede Utallighed, der følgelig ikke kan
forvexles med den blinde Ubestemmelighed. Ved denne
cirkelformige Mediation ophæves Qvantitetens ufordragelige
Længde og Brede. Det evige Grundforhold
imellem Antal og Utallighed, hvori dens disjecta membra, de sporadiske Qvanta, concentreres, kan nu sees
i eet Syn og fattes i een speculativ Tanke: \emph{Intensitet}en er Qvantitetens Sandhed.

I den absolute Intensitet maae nu for det Første
alle umiddelbare Discretioner forsvinde: dens Intension er, at de alle skulle være i hverandre og gaae sammen
i Eet; men ved deres fuldkomne Forsvinden vilde
Intensiteten selv forsvinde, derfor har den tillige en
afgjort Tendents til ogsaa at sætte de endelige Qvantas
Væren udenfor hverandre og give disse en Extension. Det er denne uophørlige Spænding imellem
Intension og Extension, der vedligeholder Grændsen og
den uendelige Alternation af Antal og Utallighed.
% --- p. 43 ---
\setcounter{footnote}{0}%
\opage{43}Under Extensionen begrændses det ene endelige Qvantum
ved det andet; men denne gjensidige Begrændsning, der danner Extensiteten, er kun relativ; i Intensiteten derimod have alle Qvanta deres fælleds absolute
Grændse: saamange ophævede Grændser i Intensiteten, saamange ponerede Grændser i Extensiteten.
Paa denne Maade bliver Antallet fastsat. Men da
Intensiteten har ophævet alle muelige Grændser, og
saaledes har dem i sig, maa ogsaa enhver muelig Begrændsning
kunne sættes i Extensiteten; derved er
Utalligheden bevaret.

Uagtet nu Intensiteten er den sande Grændse for alle
extensive Qvanta, saa kan den dog ikke have sin Væren
udenfor disse; thi i dette Tilfælde vilde den selv
nedsættes til et extensivt Moment; Alt gik da derved
op i Extensiteten, og Extensiteten selv maatte forsvinde,
da den ei kan bestaae uden Intensitet. Som Følge
heraf maa Intensiteten soutenere Discretionen og Continuiteten
derved, at den stedse sætter et \textbf{discret} Moment
som Udgangspunkt for \textbf{Ophævelsen} af alle øvrige. Tages der altsaa Hensyn til et enkelt bestemt Moment
for sig, da ligger det i Sagens Natur, at Intensiteten
er Extensiteten modsat. Det enkelte Moments intensive
Styrke bestaaer nemlig deri, at det har alle de
øvrige Momenter i sig. Men ifølge Extensiteten ere Momenterne udenfor hverandre. Det er følgelig en
ligefrem Absurditet, at et enkelt Moment skulde i een
og samme Betydning baade have og ikke have de
øvrige i sig. Fra dette Synspunkt slutter man med
Rette: \textbf{Jo større Extensitet}, \textbf{desto mindre} Intensitet. De Momenter, der i Extensiteten have deres Væren
udenfor hverandre, ere imidlertid de selvsamme, som
de, der i Intensiteten have deres Væren i hverandre;
% --- p. 44 ---
\setcounter{footnote}{0}%
\opage{44}Extensiteten er følgelig Intensitetens eget discrete og
distincte Udtryk: \textbf{Jo større Intensitet}, \textbf{desto større
Extensitet}. Idet nu Intensiteten skal acqviescere i et enkelt discret
Moment, medens den ophæver alle de øvrige Momenters Discretion, og hvert enkelt Moment i Intensiteten og Extensiteten her maa have samme Ret, saa
gjælder det om dem alle, at de hvert for sig skulle
sættes i deres Discretion ved de øvriges Ophævelse. Paa Grund heraf maa den absolute Intensitet extendere
sig i en \textbf{Række af} intensive \textbf{Momenter}.
Ethvert Moment er en Intensitet, men kun i Relation
til de øvrige ophævede Momenter. Ifølge Momenternes
indbyrdes Forskjel maa Ophævelsen foregaae
paa forskjellig Maade i hvert enkelt af dem.
Saaledes udvikler sig Intensiteternes Discretion. Men
denne Differents er endnu kun qvantitativ. Denne
qvantitative Forskjel kan ikke søges i Momenternes større
og mindre Udstrækning; thi deres umiddelbare Extension
er forsvunden. Det enkelte discrete Moments Extensitet
ligger alene i de øvriges Ophævelse, og er følgelig
selve dets Intensitet. Ligesaalidt kan Antallet
af de ophævede Momenter danne nogen Differents; thi
naar et hvilketsomhelst enkelt sættes i dets Discretion, sættes derved alle de øvrige som ophævede. De ophævede Momenters Antal er stedse lige stort. Forskjellen
kan da alene fremgaae af \textbf{de discrete} Intensiteters
\textbf{udviklede Relativitet}.

Ethvert intensivt Moment ophæver en Utallighed af
Momenter, hvori det selv utallige Gange skal ophæves.
Forsaavidt det har alle øvrige i sig, indeholder
det vel den hele Intensitet; men da det igjen hjemfalder
som ophævet til de øvriges intensive Discretion, er
det kun et relativt Udtryk for den absolute Intensitet.
% --- p. 45 ---
\setcounter{footnote}{0}%
\opage{45}Den absolute Intensitet sætter da enhver relativ Intensitet som ophævet Moment for en anden relativ Intensitet: jo flere ophævede Intensiteter, den relative
Intensitet indeholder, desto mere nærmer den sig den
absolute. Herved udfolder den absolute Intensitet sig i
en Række af \textbf{mindre} og \textbf{større} Intensiteter.

Paa Grund heraf maa man skjelne imellem en dobbelt
Ophævelse af de discrete Momenter. a) Den almindelige
Ophævelse reducerer Discretionen til dens Udspring.
Den sætter Intensitetens Continuitet og Eenhed som det
Primære, Extensitetens Discretion og Fleerhed derimod
som det Afledede og Secundære. I den universelle Ophævelse ere de discrete Qvanta som saadanne endnu ikke
satte. Men først naar Noget er sat, kan det virkelig
ophæves. Derfor følger nu b) den specifiske Ophævelse,
som forudsætter visse discrete Momenter allerede
ponerede, for ved deres Opgaaen i et efterfølgende Moment at danne en virkelig relativ Intensitet. Den relative
Intensitets eiendommelige Styrke beroer da ikke
paa, hvormange Momenter den indeholder ifølge den
almindelige Ophævelse (thi den indeholder alle), men
derimod paa, hvormange den ifølge den specifiske Ophævelse har i sig, det vil sige, hvormange relative Intensiteter den har under sig. Den første og svageste
Intensitet maatte nu opstaae derved, at det første discrete
Moment blev poneret ved en almindelig Ophævelse
af de øvrige; men da den relative Intensitet aldrig
dannes ved nogen blot almindelig, men først ved
en specifisk Ophævelse, saa vilde den svageste Intensitet slet ingen Intensitet være. Intensiteterne ere paa
Grund heraf utallige, Overgangen fra den laveste til
den høieste er uendelig; Bevægelsens Sandhed ligger
i Gradationen: \textbf{de relative Intensiteters Række}
\textbf{afrunder sig} i \textbf{Gradernes Cirkel}.
% --- p. 46 ---
\setcounter{footnote}{0}%
\opage{46}Graden sætter Intensiteten med dens bestemte Relativitet.
Enhver Grad er som saadan et qvantitativt
Noget for sig. Men den enkelte Grad, der i sig har
ophævet andre Grader, skal ifølge sin Bestemmelse igjen
begrændses og ophæves af en høiere Grad: dens Væren-for-sig henviser idelig til en anden For-sig-væren.
Gradernes Skala fremstiller os saaledes en Række af
for-sig-værende Momenter, der dog alle ere for hverandre indbyrdes. Denne For-sig-væren er nu Qvantitetens
egen Qvalitet, og denne Qvalitet den medierede
Relativitet. Gradationens Overgange og Nuancer ere
i og for sig utallige; men Relativiteten som saadan
er den blivende Bestemthed, der udgjør det Characteristiske
i alle muelige Gradationer, og derfor ophæver Gradernes uendelige Række. Gradationen udtrykker
netop den ufravigelige Fordring, at der være Modsætninger,
faste og urokkelige Differentspunkter i den
sande, absolute Qvantitet: \textbf{Graden adskiller sig fra}
\textbf{sig selv for at} \textbf{referere sig til sig} \textbf{selv}. Da Modsætningen hører med til Gradens egen Selvbestemmelse, maa der i den lavere Grad indeholdes et
Moment, som ikke lader sig ophæve af den høiere:
der maa være et Noget (et Qvale), der stabilerer Forskjellen
imellem de enkelte Grader. Under denne positive Specification har Graden sit Andet i sig selv; det
ligger i dens Tilværelse, at den maa forandre sig. Den
enkelte Grad, der selv er Noget for sig, er derved tillige
en Grad af Noget, som den ikke fuldkomment kan
udtrykke. Dette Noget griber da ud over den lavere
Grad, bevirker dennes Ophævelse og specificerer sig i
en høiere Grad, som derved erholder sin Væren-for-sig.
Men hvor ofte end Ophævelsen gjentages, stedse vil dog
det positive Noget forlægge sig udenfor de enkelte Grader. Derfor viser det sig ogsaa med en dobbelt Cha%
% --- p. 47 ---
\setcounter{footnote}{0}%
\opage{47}racteer. Det fastsætter for det Første Grændsen mellem
Grad og Grad, og danner dernæst den fælleds bestemte
Grændse for alle Grader. Det er dette Noget, hvorved
Graden bliver sig sit Andet, og derved tillige noget
Andet end sig selv. Da hver Grad er Noget for sig,
saa er Forandringen en Overgang fra den \textbf{ene} Grad
til den \textbf{anden}; men idet enhver Grad udvikler sig
som en Grad af de øvrige Grader, saa er Graden
stedse reduceret til den tautologiske Bestemmelse at være
en Grad af sig selv. Gradens Forandring ophæver
derved sig selv, og Gradationen mister sin Betydning.
Men naar Graden saaledes bestemmer sig som et Noget
uden Gradation, da er den ikke længere Grad; idet
Gradforandringen standser, har Graden absolut forandret
sig og er bleven til noget ganske Andet. Dog,
idet Graden taber sig i Andet, vinder den sig selv; thi
denne Andethed bliver den egentlige Basis, hvorved
Discretionen restaureres, og Gradationen kommer til
sin Ret. Hver enkelt Grad soutenerer nu sin specifiske Betydning, og Graderne blive hverandres indbyrdes
Andet, og sætte hverandre Grændser i Relation til det
qvalitative Andet, der ponerer alles absolute Grændse. Dette Andet kan imidlertid ikke saaledes fremtræde i
de enkelte Grader, at det selv specificeres; thi derved
vilde Graderne adskilles i lutter qvalitative Differentser,
og Discretionen blive saa prædominerende, at Continuiteten og med den Discretionen selv forsvandt. Heraf
følger, at \textbf{den Qvalitet}, \textbf{som ligger til Grund}
\textbf{for den virkelige Gradation}, \textbf{ikke selv er} Gradationen
\textbf{underkastet}\footnote{Fra dette Synspunkt have Stoikerne negtet Grader i Dyd og Forbrydelse.}.
% --- p. 48 ---
\setcounter{footnote}{0}%
\opage{48}Men naar Specificationen skyldes det qvantitative Noget,
og dette Noget, som Qvalitet, falder udenfor Gradationen, saa falder Gradationen udenfor sig selv; thi
den kan ikke undvære Specificationens Momenter.
Graden kan altsaa betragtes som den almindelige Negation
og Grændse, ved hvilken Qvaliteten fuldfører
sine egne Specificationer, efterdi ethvert Led i Rækken
determineres ved en Grad. Paa denne Maade kommer
Qvaliteten til at udvikle sig i en Gradation; men idet
Qvaliteten slaaer over i sit Andet og bliver Grad, slaaer
det, der forhen sattes som Grad, over i sit Andet, og
bliver Qvalitet. Resultatet heraf er da dette, at Gradernes Forandring eller den sande Gradation paa een
Gang er \textbf{en Afvexling af Grad} \textbf{med Grad og af}
\textbf{Grad med Gradnegation}.

Heri ligger en Modsigelse, der kommer til klar Udvikling
i \textbf{Approximationen}.

Approximationen sætter Graden af Noget (Qvantitet
med Qvalitet) som Udgangspunkt. Men ikke tilfreds
med den Maade, hvorpaa det qvalitative Noget i en
vis Grad er udtrykt, stræber man frem mod en høiere Grad, der skal give Qvaliteten en sandere Determination. Man vil imidlertid ikke have et andet Noget,
ikke gaae over til et andet Qvale; Andetheden søger
man alene i den forandrede Grad. Approximationen
\textbf{er} en fortsat \textbf{Gradation}. Det, der foranlediger
Gradationen, er alene den Fordring, at den paagjældende Qvalitet skal have sit blivende, adæqvate Udtryk.
Da Graden her er Qvalitetens eneste Determination,
saa maa det givne Qvale erholde forskjellige Bestemmelser i de forskjellige Grader. Ved Gradationen undergaaer
Qvaliteten altsaa selv en Specification. Sit
absolute Udtryk vilde det qvalitative Noget først erholde,
naar den høieste Grad var opnaaet; men Graden \textbf{er}
% --- p. 49 ---
\setcounter{footnote}{0}%
\opage{49}Qvantitet, Gradernes Antal utallige: Qvalitetens adæqvate Determination opnaaes følgelig aldrig. Er der
nu en relativ Forskjel imellem det qvalitative Noget,
sat i en lavere og høiere Grad, saa maa der være en
uendelig Forskjel imellem dets Determination, saaledes
som den kan gives i relative Grader, og saaledes som den
skulde gives i den absolute Grad. I Approximationen
ligger saaledes Gradationens Utilstrækkelighed; det qvalitative Noget man stræber at udtrykke
i de relative Grader, viser sig derved som et qvalitativt
Andet, der ligger udenfor alle Grader, d. e. som et
uopnaaeligt Maal\footnote{Den subjective Idealisme, der postulerer Udødeligheden, paa Grund af den endeløse Splid imellem Fornuften og Driften; Panlogismen, der urgerer Aandens uendelige Personificationsproces, men negter Guds absolute for sig værende Personlighed, befinde sig begge i Approximationens uopløste Modsigelse.}. Just fordi man paa Approximationens Standpunkt
saa umiddelbart fortroer sig til Gradationen og lægger
den qvalificerende Andethed i Graden selv, finder man
sig ligesaa umiddelbart nødsaget til at betragte det begrændsende
Andet som liggende udenfor enhver Grad. Naar det derimod bringes til Bevidsthed, at Gradens
Andet overalt baade er qvantitativt og qvalitativt, og
følgelig ligger baade i og udenfor Graden, da viser
Maalet sig paa een Gang baade som fraværende og
nærværende (quod petis hic est): den sande Approximation
er en uendelig Evolution.

\overgang

Qvaliteten udvikler sig som en Relation imellem for-sig-værende
- værende Momenter. Den qvalitative Grændse er
fast og ubevægelig; den kan ikke forandres, uden at
% --- p. 50 ---
\setcounter{footnote}{0}%
\opage{50}Qvaliteten selv forandres: Qvaliteten er som saadan
den fixerede Bestemthed.

Qvantiteten derimod er Grændsens Selvbegrændsning,
der har sin reneste For-sig-væren i Tallet. Men Grændsens Selvbegrændsning er Ophævelse af enhver
bestemt Grændse, og Tallet reproducerer sig selv i det
Uendelige. Qvantiteten er den bevægelige Grændse,
den ubestemte Formerelse og Formindskelse er dens Bestemthed. Qvaliteten og Qvantiteten ere altsaa, betragtede i
deres Umiddelbarhed, ligegyldige for hinanden og kunne
derfor ikke begrændse hinanden.

Men denne Ligegyldighed, der hidrører fra deres
umiddelbare, categoriske Forskjel, ophæver sig selv.

Idet nemlig Qvaliteten bliver sig selv sit Andet,
gjenfinder den sig selv paa hiin Side af den fixerede
Grændse, dens Selvbegrændsning er altsaa dens ubestemte
Selvudvidelse og Selvindskrænkning, den bliver
derved noget Andet end sig selv og bestemmes som
Qvantitet. Qvantitetens Grundbevægelse er Grændseophævelse. Men Grændsen kan ikke hæves, uden at
den først maa være poneret: Continuiteten kræver Discretion. I den ponerede Grændse og i Discretionen
som saadan gjenfinde vi altsaa de qvalitative Bestemmelser i Qvantiteten. Idet Qvantiteten qvalificerer sig og
sætter sit Andet, bliver den, som Approximationen viste,
det Negative af sig selv ɔ: Qvalitet. Resultatet er altsaa, at Qvaliteten medierer sin Selvbegrændsning gjennem Qvantiteten, og at Qvantiteten begrændser sig selv ved Qvaliteten.
% --- p. 51 ---
\setcounter{footnote}{0}%
\opage{51}\afsnit{C. Modalitet.}
\underafsnit{a) Den rene Modalitet.}
\paragraf{17}{Qvantitativ Qvalitet, qvalitativ Qvantitet, Modification.}

De qvalitative og qvantitative Momenters gjensidige
Relation udvikles i \emph{Modaliteten}. Ethvert
Qvale maa som saadant qvalificere sig. Dets Qvalification er en Selvbevægelse gjennem alle relative
Begrændsninger og Bestemmelser, og som Følge heraf
en Qvantification: det i sig begrændsede Qvale er en
\emph{qvantitativ Qvalitet}. Det bestemte Qvantum
resulterer af en Qvantificeren. Qvantificationen er
en Udvikling af en bestemt Relation imellem discrete
ɔ: qvalitativt forskjellige Qvanta: ethvert virkelig bestemt
Qvantum maa derfor sættes som \emph{qvalitativ Qvantitet}. Qvalificationens og Qvantificationens Overgang i hinanden er en \emph{Modification}, der
indleder Qvalitetens og Qvantitetens bestemte Tilværelse.

 Categorierne ere positive Udtryk for hypostaserede Tanker. Derved at Tanken i Categorien faaer sin faste Skikkelse, fremtræder den med en vis Indskrænkning, og, som Følge heraf, med en bestemt Eensidighed. --- Bevidstheden om denne Eensidighed fremkalder Modsigelsen, denne igjen den dialectiske Bevægelse; derfor gaaer Dialectiken stedse Categorierne imod. 
% --- p. 52 ---
\setcounter{footnote}{0}%
\opage{52}I Qvalitetscategorierne søger Tanken overalt sit
Holdningspunkt i den faste, ubevægelige Grændse: Væren er \textbf{ikke} Intet, Noget ikke Andet, Eet ikke
Flere. Men Dialectiken nedbryder stedse den ponerede
Skranke og udtrykker Grændseopløsningen i en særegen
Overgangscategorie (Vorden, Forandring o.s.f.). --- Overgangscategorierne
selv ere med Hensyn til den paagjældende
Categorieclasse deels \textbf{relative}, deels absolute. Den relative Overgangscategorie viser blot Utilstrækkeligheden
af enkelte underordnede Categorier og indleder
derved Tankebestemmelsen i høiere Categorier af samme
Classe; den absolute Overgangscategorie derimod fremstiller
den udviklede Categorieclasse i dens hele Eensidighed,
og fører derved over i et System af nye Categorier\footnote{Saaledes er Vorden en relativ Overgangscategorie; thi uagtet den ophæver den qvalitative Grændse imellem Væren og Intet, saa indleder den dog en ny qvalitativ Grændse imellem Tilværelsens Noget og Andet. Først naar den qvalitative Grændse har naaet sin høieste Bestemthed som exclusiv Grændse, og Exclusionen i sin absolute Gjennemførelse har viist sig som Assimilation, er enhver umiddelbar Qvalitetsgrændse at ansee som ophævet.}.

Idet den sidste og høieste Qvantitetsgrændse ophæves, saa bringes det herved til klar Erkjendelse, at
enhver poneret Grændse er variabel. Da nu den variable
Grændse som saadan er Qvantitetens Criterium,
saa følger heraf, at det, der kommer til bestemt Udtalelse
i Qvantitetscategorierne, ikke er Andet end det,
der allerede implicite har bevæget sig gjennem hele
Qvalitetsdialectiken. Men idet de qvalitative Bestemmelser forflygtiges i Qvantitetscategorierne, maa Dialectiken igjen hæve denne Eensidighed og idelig 
restau%
% --- p. 53 ---
\setcounter{footnote}{0}%
\opage{53}rere den forsvindende Grændse. Det var disse Restaurationsforsøg,
der i foregaaende Afsnit førte til stedse
høiere og fyldigere Categorier og derved omsider til den
Erkjendelse, at de qvantitative Bestemmelser maatte
integreres ved at gjentage den forsvundne Qvalitet. Altsaa: man begynder med de \textbf{værende}, umiddelbart
givne Qvaliteter (Qualia). Deres Umiddelbarhed ophæves,
idet de gjennemgaae deres \textbf{Vorden}. Denne
Vorden, hvori Qvaliteterne forsvinde, antager først
Characteren af en Qvantification. Men da Qvantificationen
er Qvaliteternes egen Genesis, maa den i sin
Udvikling vise sig som den egentlige Qvalification og
saaledes føre til qvalitative Resultater. Det er denne
Grundbevægelse, der bestemmer hele \textbf{Modaliteten}.

Enhver sand Qvalitet maa qvalificere sig. Saalænge Qvaliteten endnu ikke har gjennemgaaet sin Qvalification, repræsenterer den sig som et Værende, et
sandselig Givet, der blot tilhører Forestillingen. Den
hele sandselige Verden er i dens Umiddelbarhed et Indbegreb af saadanne uqvalificerede Qvaliteter. Ved at
reducere Sandseverdenen til sine simpleste Elementer,
kommer man til den crasse \textbf{Atomistik}\footnote{Selv de herbartske Realia ere uqvalificerede Qvaliteter.}. Atomerne
ere sandselige Umiddelbarheder; da de mangle Qvalificationen,
mangle de nødvendigviis ogsaa Qvantificationen,
og ere følgelig \textbf{uden Udstrækning}. Saaledes har da heller intet Atom Andethedens Moment i
sig selv; Atomet fremstiller kun det abstracte Noget,
hvis Andet er udvortes og ligegyldigt: Forskjel følger
paa Forskjel uden al Overgang. Det Medium, hvori
Atomerne befinde sig, er derfor ogsaa et udvortes Medium (Rummet, det Tomme) og den Magt, der forbinder
og adskiller dem, en udvortes ubegribelig Magt. Dette er da de sandselige Atomers logiske Modsigelse,
% --- p. 54 ---
\setcounter{footnote}{0}%
\opage{54}at de skulle have en Væren uden Vorden, en Tilværen
uden indre Forandring, en For - sig væren uden
gjensidig Exclusion og Assimilation.

Den rene Væren repræsenterer sig først som given
Qvalitet, som den \textbf{værende} Væren; Intet kan ligeledes fremstille sig som det \textbf{værende} Intet. Men idet
Væren resulterer af Intet, sætter den derved Intet
som en Grændse, den selv ophæver, medens Intet i
omvendt Retning tager sig selv tilbage gjennem Værens
ophævede Grændse. Begge \textbf{qvalificere} og extendere
sig, idet de \textbf{vorde} hver paa sin \textbf{Maade}: Modalitetens Qvalitet er en \textbf{qvantitativ Qvalitet}.

Den modale Væren sættes som den Urqvalitet, der
ved sin Qvalification frembringer alle værende Bestemmelser.
Men i den blotte Qvalificationsproces bliver
alt det Faste flydende og forsvindende, og Væren den
intensive Størrelse, der extenderer sig \textbf{nden Maal % PRINTED AS IS: nden (for uden; S066)
og Maade}. Den abstracte Qvalification er en blot
Qvantification; den sees her som en Tilbliven, der aldrig
kan blive til Noget, en rastløs Forandring, der aldrig
fører til For-sig-væren, en Eenhed uden Fleerhed, og
er følgelig en eensidig Abstraction. Da den abstracte
Qvantification aldrig bringer det til fast qvantitativ Bestemthed, kan den heller aldrig constituere bestemte
Qvanta. Qvalitetsmomenterne, der i det Uendelige
komme og svinde, ere utallige, og paa Urqvaliteten selv,
der gaaer op i den rene Qvantificationsproces, kan
Tallet slet ikke anvendes. Qvantificationen maa derfor
sættes som Moment for Qvalificationen, og denne
igjen føre til bestemte Qvaliteter. Som Begyndelsesmoment maa Urqvaliteten eller Almeenqvaliteten være den meest abstracte; jo længere
derimod Qvalificationen fortsættes, desto mere concrete
% --- p. 55 ---
\setcounter{footnote}{0}%
\opage{55}blive de deraf resulterende Qvaliteter. Det er ved
disse Grupper af universellere og speciellere Qvaliteter\footnote{Vi bruge her „Qvalitet“ i den meest udstrakte Betydning. Forskjellen imellem Qvaliteterne som selvstændige Existentser og uselvstændige Momenter kan først senere fremhæves.}
(genera og species), at Qvantiteten egentlig kommer
til sin Ret og bliver et Moment med i Tilværelsen.
Ligesom nemlig Qvantitetsbestemmelsen gaaer under,
naar alle Qvaliteter forsvinde i Almeenqvaliteten, saaledes
vilde den ogsaa tabe sin Betydning, hvis Qvalitetsbestemmelsen
dreves til sin yderste Spidse, og Specialiteten fastholdtes for sig. Kun derved, at de speciellere Qvantiteter reduceres til en universellere Qvalitet,
fremkommer en \textbf{qvalitativ} Taleenhed, og, idet
den universelle Qvalitet atter henpeger til sine Specificationer,
et \textbf{qvalificeret} Antal. Ved, f. Ex., at reducere
4 Borde og 3 Stole til den fælleds Qvalitet
„Meubel“, erholdes en qvalitativ Eenhed, hvormed Numerationen
kan begynde. Men bliver man nu staaende
ved den universellere Bestemmelse, saa er Meubel = Meubel; først ved at fremhæve Specificationen, bliver
det \textbf{ene} Meubel seet i dets Forskjel fra det \textbf{andet}, og
man kan nu tælle (bestemme \textbf{Antallet} af) de særskilte Meubler. Den modale \textbf{Qvantitet} er stedse
\textbf{qvalitativ}. Qvalificationens og Qvantificationens
Overgang i hinanden (quod modum efficit) lader sig
nu nærmere bestemme som \textbf{Modification}.

I Modaliteten blive Noget og Andet hvert for sig
at sætte som \textbf{qvalificerede} Qvaliteter. Disse Qvaliteter subsistere nu med relativ Selvstændighed i et
System af qvalitative Momenter. Bestemmelsen for
det qvalitative Noget er stedse den, at det maa gaae
% --- p. 56 ---
\setcounter{footnote}{0}%
\opage{56}under i Andet; Andet er dets \textbf{absolute} Grændse; hvor
Andet begynder, maa Noget ophøre. Men for at
Noget kan undergaae, maa det først opstaae. Det opstaaer
ikke med eet Slag; Qvaliteten maa qvalificere
sig; den har sin Væren i sin Vorden. Derfor blive
nu ogsaa de tvende Momenter, der i den abstracte
Vorden umiddelbart gik i Eet, hvert især udsondret
og specificeret i Modaliteten: den modale Vorden modificeres
ved en bestemt Modsætning mellem \textbf{Opstaaen}
og \textbf{Forgaaen}.

Idet Noget begynder sin Qvalification, sætter det
sit første qvalitative Moment, et Moment, hvori det
kan siges baade at være og ikke at være. Qvalificationen skrider derfor successivt frem, idet den stedse
ponerer nye Qvalitetsmomenter og bevarer de allerede
ponerede. Men Qvaliteten, som heraf resulterer, er
som saadan indskrænket; den er et Noget, der alsidig
begrændses af Andet. For at danne den fulde Begrændsning
af Noget, maa Andet altsaa baade gaae
foran og følge efter. Det forudgaaende Andet er Noget for sig, det efterfølgende ligesaa; Bevægelsen er
denne, at ligesom Noget opstaaer ved en Negation af
det forudskikkede Andet, saa fremgaaer det følgende Andet
igjen af det negerede Noget. Men denne Modsætning, der saaledes indtræder i Andetheden selv, skal
her blot sees som Moment for Tilværelsen af det mellemtrædende
Noget. Under sin Opstaaen maa Noget
idelig forandre sig. Forandringen bestaaer deri, at det
excluderer det forudsatte Andet ved at assimilere dets
Andethed og gjøre sig selv til sit Andet. I denne
Forandringens Modification har Noget altsaa sin Bestaaen,
sin \textbf{Bliven}\footnote{„Blive“ udtrykker baade \textit{manere} og \textit{fieri}.}. Noget bliver derved, at dets
% --- p. 57 ---
\setcounter{footnote}{0}%
\opage{57}Vorden (Qvalification) sættes som Moment for dets
Væren (Qvalitet). Men under Assimilationen af Andet
har Noget optaget sit eget Intet i sig. Andethedens
Momenter kom fra Andet, da dette skulde gaae
over til Noget, men nu er det igjen Nogets Bestemmelse
at gaae over i Andet; de Andethedsmomenter,
ved hvis Assimilation det har naaet sin Bestaaen, udvikle sig derfor til det, hvad de virkelig ere, nemlig
Negationen af selve Noget. Noget bestaaer i et Intet,
der igjen udvikler sig til positiv Væren under
Qvalificationen af det efterfølgende Andet\footnote{Den Grundtanke, at Noget bliver sig sit Andet gjennem noget Andet end sig selv, er af gjennemgribende Indflydelse paa den philosophiske Verdenserkjendelse. Atomistiken, der acqviescerer i den Betragtning, at Noget har Andet udenfor sig, fører til Materialisme; thi det hører Aanden til at blive sig selv sit Andet: Orientalerne, der lade Urqvaliteten være sig selv sit Andet uden egentlig at ponere noget Andet end sig selv, ere Pantheister. Den fichtiske Idealisme erkjender vel, at Intet skal assimilere sig sin Andethed gjennem sit Andet; men da dette Andet ikke i sig qvalificeres, bliver det staaende som umiddelbar Skranke. Christendommen erkjender Gud i sin Personlighed; thi her er han sig sit eget Andet (Sønnen) gjennem et i sig qvalificeret Andet (Verden).}. Nogets
Bestaaen indleder saaledes dets Forgaaen; med denne
Modification er Forandringen af Noget en Tilintetgjørelse.
Men hvis nu Noget forgik i samme Nu,
som det opstod, vilde der hverken være Tilbliven eller
Tilintetgjørelse: enhver af Forandringens Modificationer maa have sin Extension. For at soutenere sine
modale Modsætninger, modificerer Qvalificationen sig
til Qvantification.
% --- p. 58 ---
\setcounter{footnote}{0}%
\opage{58}Idet Noget under Qvalificationen sætter og ophæver
sine relative Grændser for at naae den absolute
Grændse, sætter det sig som et bestemt \textbf{Continuum},
og naar den absolute Grændse overskrides, begynder
Andet sin Continuation. Medens Discretionen altsaa
under den hele Tilbliven opgaaer i den qvalificerede
Continuitet, prædominerer Discretionen derimod i selve
Grændsen mellem Noget og Andet, og qvalificerer sig
derved i discrete Continua. Qvalitetseenheden af Noget
er saaledes at bestemme som en \textbf{intensiv} Størrelse,
der \textbf{extenderer} sig i Systemet af de qvalitative
Momenter. Paa denne Maade komme Qvaliteten
og Qvantiteten begge hver til sin Ret. Qvaliteten
kræver sin faste, bestemte Grændse; denne erholder den
nu ved at overføre det Variable paa Qvantiteten.
Men Qvantiteten varierer kun, saa længe den paagjældende Qvalification varer. Ved Qvalitetens Indskrænkning
erholder Qvantiteten sin faste, ubevægelige Grændse,
ja bliver endogsaa som den, der omfatter alle den særegne
Qvalitets qvalitative Momenter, det eneste bestemte
Udtryk for Qvaliteten i dens Heelhed: hvor
Qvalitetens absolute Grændse ligger, kan alene afgjøres
ved Qvantiteten. Det Constante i Qvaliteten betegner
det Foranderlige i Qvantiteten; men Qvantitetens
constante Grændse udtrykker den hele Qvalitets
egen Forandring. Saaledes fører Modificationen til

\underafsnit{b) Maalet.}
\paragraf{18}{Det umiddelbare Maal, det qvalificerede Maal, Maalets Modification.}

Den qvalificerede Qvalitet udtrykker sin Tilværelse
ved den tilværende Qvantitet, og bliver saale%
% --- p. 59 ---
\setcounter{footnote}{0}%
\opage{59}des bestemt i \emph{Maalet}. For Qvaliteten i dens
blotte Tilværen er den qvantitative Bestemmelse overfladisk
ɔ: et \emph{umiddelbart} Maal. Hvor Qvalitetens Væren derimod sættes i Eenhed med dens Tilbliven,
udvikler Qvantificationen sig af selve Qvalificationen, og Maalet qvalificeres. Qvalitetens \emph{qvalificerede},
immanente Maal bestemmes ved en Relation til andre Qvaliteters immanente Maal; ved denne Relation
bliver Maalet selv uendelig \emph{modificeret}. Qvalificationen maa føre til Qvalitet, ligesom Vorden til Tilværen, og Umiddelbarheden med den bestemte Grændse restaureres. Derved resulterer nu ogsaa
det bestemte Qvantum af Qvantificationen. Qvaliteten
er saaledes bleven til Qvantum; de have deres
Tilværelse i og med hinanden. Men paa samme Maade,
som Noget har qvantificeret sig og i sin Tilværen
viser sig som Størrelse, har Andet ogsaa erholdt sin
Qvantification og \textbf{er} ligeledes en Størrelse. Begge
ere da qvantitative Qvaliteter og i denne Henseende
hinanden fuldkommen lige. Men naar al Modsætning
imellem Noget og Andet ophæves, forsvinder selve Tilværen
i den rene Væren. Det gjælder altsaa nu om
at finde en ny og høiere Forskjel imellem Noget og
Andet, det vil sige, at bestemme Grændsen mellem de
qvalitative Størrelser indbyrdes.

Den \textbf{qvalitative} Grændse viste sig (§ 14, p. 26)
som den exclusive Negation; men da de tilgrændsende
Qvaliteter vare umiddelbare og som Følge heraf uden
Udstrækning, bleve de selv nedsatte til blotte Grændsemomenter.
Qvaliteten gik derfor op i Grændsens Selvbevægelse,
udviklede sig som \textbf{qvantitativ} Grændse
% --- p. 60 ---
\setcounter{footnote}{0}%
\opage{60}(p. 27), og frembragte derved en Udstrækning, der ophævede
al Begrændsning.

Den \textbf{modale} Grændse derimod er fast og bestemt
som den qvalitative, thi den adskiller \textbf{Noget} fra
\textbf{Andet}; men Noget og Andet ere \textbf{Størrelser}; den
modale Grændse er derfor bevægelig som den qvantitative,
thi den samler i sin Bestemthed Størrelsernes
relative Grændsemomenter. Den modale Grændse har
altsaa tvende Sider, hvis gjensidige Relation udvikles i
\textbf{Maalet}.

Da begge Maalets Hovedsider hver for sig ere modale,
saa indeholde de begge den umiddelbare Eenhed
af Qvalitetens constante og Qvantitetens variable
Grændse, men paa modsat Maade. Medens den ene
Side nemlig indeholder den bevægelige Grændse som
et i den constante ophævet Moment, og derved er sat
som en Qvalitet med \textbf{bestemt} Udstrækning, lader den
anden derimod ethvert constant Grændsemoment forsvinde i den varierende Grændse, og viser sig som et
Continuum med \textbf{ubestemt} Udstrækning. Udstrækningen
er da det for begge Sider fælleds Characteristicon;
dette er netop det Overfladiske i det umiddelbare \textbf{Maal}, at alle Qvalitetsbestemmelser her gaae
op i den blotte Udstrækning.

Men Qvaliteten tager idelig sig selv tilbage gjennem de qvantitative Momenter. Determinationens specifiske
Udvikling fremgaaer deraf, at den constante og
variable Grændse ere forenede i enhver af Siderne:
saaledes kunne de \textbf{maales} med hinanden. Under Maalingen
sættes det bestemte Continuum som \textbf{Maalestok}.
Maalestokken er en Qvantitet, der repræsenterer sig i
Forestillingen; \textbf{hvor} stor den er, maa umiddelbart
anskues: dens Udstrækning ɔ: dens specifiske Qvantitet
er Qvalitetens egentlige Bestemthed. Den constante
% --- p. 61 ---
\setcounter{footnote}{0}%
\opage{61}Grændse, der sættes i og med Maalestokkens Continuum,
bliver nu det Middel, hvorved det ubestemte
Continuum, ɔ: det som maales, reduceres til bestemte
Grændser. Den maalte Størrelse udvikler saaledes
begge sine Momenter: i og for sig har den en variabel
Grændse, der ikke blot mueliggjør, men endogsaa kræver constante Grændsebestemmelser; denne sidste Fordring
fyldestgjøres ved Maalestokkens Anvendelse. Men Maalestokkens
bestemte Continuum, der har den variable
Grændse i sig, kræver nu ogsaa, at Foranderlighedens
skjulte Moment udvikles og determineres. Med sin
umiddelbare Bestemthed kan Maalestokken udvides og
indskrænkes efter Behag; men i Relation til den
maalte Størrelse bliver det Vilkaarlige i dens Udstrækning
begrændset. Maalestokken er en given Eenhed,
der i den maalte Størrelse finder sit Antal.

Da det umiddelbare Maal blot udtrykker Qvantiteternes
overfladiske Forhold, følger heraf, at al dybere
Forskjel maa udelukkes, og den Qvalitetseiendommelighed,
der qvalificerer Størrelserne til at maales med
hinanden, umiddelbart gives\footnote{Derfor maales umiddelbart Linie med Linie, Flade med Flade.}.

Men den modale Tilværen udvikler sig som en Bliven
(§ 17), hvor Qvalificationens Vorden idelig afvexler
med Qvalitetens Væren. Da Qvaliteten \textbf{er}, idet
den \textbf{vorder}, maa den qvantitativt bestemmes som en
\textbf{intensiv} Størrelse, der extenderer sig i en Række af
intensive Momenter (§ 16. pg. 44). Men \textbf{Noget},
den bestemte Qvalitet, opstaaer og bliver kun ved Assimilation
og Exclusion af \textbf{Andet} (§ 17. pg. 56).
Den tilblivende Qvalitet kan derfor alene bestemmes
% --- p. 62 ---
\setcounter{footnote}{0}%
\opage{62}ɔ: \textbf{maales} ved en anden intensiv Qvalitet, der ligeledes
qvalificeres ved sin Extension.

Idet nu den ene af disse Qvaliteter sættes som ubestemt
Størrelse, hvis uvilkaarlige Formerelse og Formindskelse
blot angives ved en umiddelbar Maalestok,
vil den anden derimod vedligeholde sin qvalitative Eiendommelighed
derved, at den paa særegen Maade assimilerer
sig den variable Størrelse i forskjellig Mængde.
(Saaledes optager f. Ex. Vandet en forskjellig Varmemængde).

Den assimilerende Qvalitet har en begrændset For-sig-væren,
og, som Følge heraf, ligeledes kun en begrændset
Assimilation af dets variable Qvantum, der
her spiller Andethedens Rolle. (Vandet som Vand
optager kun en \textbf{vis} Varmemængde). Det er nu denne
Assimilation, der i dens hele Relativitet og Begrændsning
erholder sin Udvikling i det \textbf{qvalificerede Maal}.

Qvaliteten som saadan maa gives med værende Bestemthed,
følgelig ogsaa med en vis Assimilation af
dens Andet, ɔ: af den variable Størrelse. (Vandet
har ved selve Tøpunktet en \textbf{bestemt} Varme). Da
der nu forud for enhver qvalitativ Væren maa gaae
en qvalificerende Vorden, kan man kun gribe og fastholde
den constante Qvalitet i positiv Eenhed med den
variable Størrelse ved et Spring, der tilhører Umiddelbarheden som saadan. (Ved Thermometrets Nulpunkt
kan man f. Ex. baade have Vand og Jis). Men
da den værende Qvalitet ligesaameget maa \textbf{vorde} som
\textbf{være}, saa nedsætter den sin relative Væren til en
forsvindende \textbf{Tilstand}, og documenterer sin qvalitative
Bevægelse ved at lade de variable Tilstande afløse
hverandre.

Alt kommer saaledes an paa det dialectiske Forhold
imellem den constante Qvalitet og dens foranderlige
% --- p. 63 ---
\setcounter{footnote}{0}%
\opage{63}Tilstand. I den rene Umiddelbarhed fastholdes Qvaliteten
for sig og de varierende Tilstande for sig. Derfor
kan man da heller ikke begribe, hvorledes Qvaliteten
bevarer den uforandrede Liighed med sig selv, skjøndt
den dog stedse optager forskjellige Qvanta af den variable
Størrelse; endnu mindre indseer man, hvorfor
Qvaliteten, efterat den har viist sig ligegyldig imod
en Række af varierende Tilstande, omsider med eet
Slag altereres ved en blot Qvantitetsforhøielse.

Disse Modsigelser hæves ved Betragtningen af Maalets egen Qvalification: a) Da den constante Qvalitet
som saadan har et vist Qvantum af den variable
Størrelse i sig, saa ligger heri, at dens qvantitative
Varietet just udtrykker Qvalitetens egen Bestemthed,
og at denne, som Følge heraf, ikke kan være ligegyldig
mod nogensomhelst Tilstand. Det er kun igjennem
\textbf{Forandringerne}, at Qvaliteten kan udvikle sin
\textbf{constante} Characteer; disse høre altsaa med til dens
egen Qvalification. Naar det derfor hedder, at den
samme Qvalitet soutenerer sig i alle vexlende Tilstande,
saa er Sandheden denne, at Qvaliteten først
udtrykker sig fuldstændigt i alle de Tilstande, der
sættes og ophæves.

b) Da nu saaledes ingen enkelt Tilstand satisfacerer
Qvaliteten, vilde den aldrig erholde sit ædæqvate Udtryk, % PRINTED AS IS: ædæqvate (for adæqvate; S079)
hvis ikke Momenterne af en lavere Tilstand kunde
opbevares i en høiere. Men her er det, at Gradationen
udvikler sig i sin fulde Betydning. Den foranderlige
Størrelses ydre, umiddelbart bestemte Mængde
bliver i enhver given Tilstand assimileret til eiendommelig
Eenhed med Qvaliteten. Denne har altsaa
Størrelse i sig som bestemt intensiv Størrelse, det er,
ved en ligelig Fordeling af de Intensitetsmomenter, der
skulle characterisere de enkelte Tilstande, som Grad.
% --- p. 64 ---
\setcounter{footnote}{0}%
\opage{64}Graden er her sat i og ved en Qvalitet, derfor gives
her en \textbf{laveste} og \textbf{høieste} Grad. I den laveste Grad
er den variable Størrelses Mængde ophævet som Moment
for den Tilstand, hvori Qvaliteten begynder sin
Tilværen; i den høieste derimod fremtræder Qvaliteten
absolut identificeret med Tilstanden selv. (Ved den
100de Grad (C) har Vandet assimileret sig al den Varme,
det ifølge sin eiendommelige Qvalitet formaaer at
optage). c) Saalænge nu den relative Intensitet ikke
udtrykker Qvalitetens fuldendte Assimilationsact, saalænge
kan Qvaliteten overskride enhver Qvantitetsgrændse,
og forandre sin Tilstand uden selv at forandres;
naar Assimilationen derimod har naaet sin Grændse
og den adæqvate Tilstand, hvori den gjennem de qvantitative
Forandringer udviklede Qvalitet acqviescerer,
indtræder, da ligger det i Sagens Natur, at denne
Tilstand ikke ved en forhøiet Grad kan forandres, uden
at Qvaliteten med den altereres. Qvalitetens overraskende
Forandring er saaledes indledet ved Qvantitetsforandringen.
Herved er den constante Qvalitet begrebet
som den foranderlige Størrelses immanente Maal.

I det saaledes qvalificerede Maal have vi da: α) en
ved Qvaliteten positiv bestemt Qvantitetseenhed (fra
0° til 100° (C.) indeholdes altid et ligestort Varmeqvantum,
naar ikke andre Omstændigheder indvirke
paa Vandet); β) et ved denne Eenhed determineret
Antal, der poneres og umiddelbart observeres, idet den
variable Størrelse i Qvaliteten gjennemgaaer sine Gradationer;
og endelig γ) en sand Identitet af Eenhed
og Antal (Vandets Varme ved den 100de Grad indeholder 100 Varmegrader).
% --- p. 65 ---
\setcounter{footnote}{0}%
\opage{65}Det immanente Maal udvikler sig giennem utallige
Gradationer (imellem Frysepunkt og Tøpunkt ligger
ikke blot 80 eller 100 Grader, men, ifølge den uendelige
Deelbarhed, cfr. pg. 30, en vilkaarlig Mængde);
som Følge heraf bliver enhver Talbestemmelse ophævet.
Men hvor Tallet er forsvundet, er Maalets specifiske
Grændse udslettet; skal det qvalificerede Maals Gradationer
bestemmes, da maae Gradernes Marker vilkaarligt
afsættes, og Gradationen ved dem umiddelbart
observeres.

Men naar det qvalificerede Maal saaledes afhænger
af den umiddelbare Maalestok, opgaaer det derved i
uqvalificerede Bestemmelser, og bliver selv umiddelbart.
Denne Umiddelbarhed har en dobbelt Side: deels ligger
den nemlig i Maalestokkens vilkaarlige Inddeling,
deels i den constante Qvalitets egne Grændser, idet
denne ved sin Alteration umiddelbart angiver, naar
den variable Størrelse i den har overskredet sit Maal.

Den variable Størrelse selv altereres derimod ikke
(Varmen vedbliver at være Varme, efterat Vandet er
blevet til Damp), men søger sig indenfor en anden
constant Qvalitets Grændser et andet immanent Maal o. s. f. Herved udvikler det qvalificerede Maal sig i en
Række af immanente Maal, og bliver saaledes sig selv
udvortes. Af denne Udvorteshed fremgaaer nu Maalets
egne Specificationer. Da den variable Størrelse
determineres ved Totaliteten af de constante Qvaliteter,
hvori den finder sine immanente Maal, er dens Qvantum
herved specificeret; men ogsaa de constante Qvaliteter
ere qvantitativt specificerede, thi deres qvalitative
Eiendommelighed bestemmes alene ved den særegne
Mængde, de hver for sig optage af den variable Størrelse.
Sammenlignes nu disse Mængder med hverandre
indbyrdes, saa er Maalet reduceret til constante
% --- p. 66 ---
\setcounter{footnote}{0}%
\opage{66}Forhold. Qvantitetens varierende Umiddelbarhed er
herved ophævet, og den umiddelbare Maalestok, hvis
vilkaarlige Inddeling nu ikke længere har nogen Indflydelse paa Bestemmelsen af de constante Forhold selv,
nedsat til et udvortes Middel, hvorved det immanente
Maal maaler sig selv. (For at maale Temperaturen
er det ligegyldigt, enten Thermometret er inddeelt i
80° eller i 100°).

Reduceret til det constante Forhold har det immanente
Maal sit specifiske Udtryk i Forholdets Exponent.
Enhver umiddelbar Qvalitet er som saadan eensidig;
den søger sin Negation og Bestemmelse i en modsat
Qvalitet, ligesom denne i hiin; begge ere da at ansee
som hinandens Sider, og den fuldstændige Qvalitet som
begges Eenhed. Den udviklede Qvalitet er altid en
Dobbeltqvalitet, der i Eenhed med sig exponerer Forholdet
mellem trende qvalitative Sider. Men idet de
qvalitative Sider bestemme hinanden indbyrdes og negere
hinanden, forsvinder den qvalitative Umiddelbarhed
i deres gjensidige Relation; den ene er saaledes
optagen i den anden, den ene indeholdes i den anden,
den enes Formerelse medfører, hvis Forholdet skal bestaae,
en Formerelse af den anden: deres gjensidige Bestemmelse er aldeles qvantitativ. Da Forholdet, som
Exponenten viser, kan vedblive at være det samme,
naar kun Formerelsen eller Formindskelsen foretages paa
tilsvarende Maade med begge Sider (1:2 = 2:4),
saa har Qvantiteten i Exponenten erholdt en fast, specifisk
Characteer, medieret gjennem de qvantitativt variable
Sider. Denne specifisk bestemte Qvantitet udtrykker
netop den explicerede Qvalitet, idet de qvantitative
Siders Ligevægt udtrykker disses qvalitative Bestemthed. Saaledes bliver det immanente Maals fuldkomne
Udvikling at søge i Exponenten,
% --- p. 67 ---
\setcounter{footnote}{0}%
\opage{67}Men i det ovenfor udviklede System af immanente
Maal er denne Exposition endnu ikke gjennemført.
Vel indlyser det, at den specifiske Qvalitet overalt betinger
Qvantitetens Specification og omvendt; men
disse Specificationer ere endnu ikke deducerede af hinanden.
Den variable Størrelse, der kan betragtes som
de constante Qvaliteters fælleds Side, som den gjennem
dem alle gaaende Almeenqvalitet, tages som den
umiddelbart forefindes. Da ingen Formerelse eller Formindskelse
nogensinde bringer den til at opgive sin
Qvalitets Eiendommelighed, saa er ogsaa den hele Alteration,
den frembringer i de constante Sider, hvori
den finder sit Maal, en blot Modification. De constante
Qvaliteter bevare derfor alle Umiddelbarhedens
Præg, hvoraf da ogsaa følger, at Exponenterne ligeledes
vise sig som umiddelbare Talbestemmelser.

Skal den specifiske Qvalification opgaae i specifisk Qvantification, maae Qvaliteterne tilsætte deres Umiddelbarhed
i Qvantiteten, for i det qvantitative Forhold
at realisere den sande Qvalitet. Ethvert umiddelbart
Qvale maa da vise sig, som det er, nemlig som det
Neutrum, der hverken er sig selv eller sit Andet, men
venter paa et andet Neutrum, der vil gjøre det til sit
Andet, medens det selv gjør sig til et Andet for dette.
Denne Qvalification repræsenteres i den chemiske Neutralisation,
hvor Exponenten kan siges at være den
udviklede Qvalitet selv.

Lad A og B være de qvalitative Sider (Stoffer), af
hvis Foreening den qvalitative Exponent (Legemet) E
resulterer; deres umiddelbare Qvaliteter ere da ikke
længere at søge i E, thi i dette Tilfælde vilde E blot
være en Sammensætning af A og B, men ingen Qvalitet
for sig; men idet de umiddelbare Qvaliteter negeres
og ophæves i Qvantiteten, bliver deres Qvanti%
% --- p. 68 ---
\setcounter{footnote}{0}%
\opage{68}tet selv derimod opbevaret: i qvantitativ Henseende er E = A + B.

Her viser sig nu Qvalitetens og Qvantitetens gjensidige
Specification paa det Allerklareste. Hvilke Qvalitetsbestemmelser,
der ligge i A og B, kan alene
kjendes af det qvantitative Forhold, hvorunder de forbinde
sig med hinanden; men paa den anden Side kan
man heller ikke bestemme dette Forhold, uden netop ved
at prøve begges qvalitative Eiendommelighed, og adspørge
dem selv, under hvilke qvantitative Betingelser
de ville foreene sig. Heraf forklares det: α) at A og
B stedse maae foreene sig i et vist Forhold, for at
frembringe en vis qvalitativ Exponent; β) at man af
et saadant Forholds Natur ikke umiddelbart kan slutte,
at det er det eeneste; thi da det i sin Exponent kun
udtrykker een bestemt Qvalitet, kunne begge Sider under
en ny Forholdsexponent frembyde nye Qvalitetsmomenter;
endelig kunne A og B hver for sig være
Sider for flere andre umiddelbare Qvaliteter: A kan
f. Ex. foreene sig med C, D f.

Forholdet mellem Siderne selv og deres Exponenter
er nu dette, at hine opgaae som Momenter i qvantitative
Forhold, medens disse derimod repræsentere den
qvalitative For-sig-væren. I Rækken: E (A + B), E′ (A + C), E″ (B + C), E‴ (C + D) kan A bestemmes aldeles
qvantitativt; det viser sig derved som den variable
Størrelse, der i de constante Qvaliteter B og C
finder sine immanente Maal. Betragtes derimod A og B som constante Qvaliteter, da repræsenterer C den
variable Størrelse; det Samme gjælder om B, der som
Fælledsqvalitet for A og C ligeledes antager en qvantitativ
Characteer og af disse optages i forskjellig Mængde. Exponenterne E, E′, E″, derimod exclu%
% --- p. 69 ---
\setcounter{footnote}{0}%
\opage{69}dere hverandre som for-sig-værende Qvaliteter, og udtrykke
derved de qvantitative Siders specifiske Forhold.

Dog selv denne Modsætning imellem Exponenterne
og deres Sider maa, naar den nærmere udvikles, ophæve
sig selv. Exponenterne, som ved Negation af
Sidernes Qvalitet have gjennemgaaet deres Qvalification,
antage selv Umiddelbarhedens Præg, og kunne
saaledes igjen blive Sider for Exponenter af høiere
Orden. (E + E‴ = E⁗), og da enhver umiddelbar
Qvalitet som saadan forudsætter en Qvalification
(pg. 53), ere de underordnede Sider selv at betragte
som Exponenter Forsaavidt
nu alle Exponenter reduceres til Sider, ophæves alle
Qvalitetsbestemmelser i Qvantiteten; forsaavidt derimod
alle Sider sættes som Exponenter, tager den qvalitative
For-sig-væren sig selv tilbage paa alle Puncter.

Den Spænding, der herved indtræder imellem Qvaliteten
og Qvantiteten, hæves i \textbf{Affiniteten}.

Ifølge Affiniteten kunne Qvaliteterne ikke forene sig
i Flæng, men maae være disponerede for hinanden;
deres gjensidige Assimilation medfører en Exclusion,
hvorved Qvalitetsmomentet, som saadant, gjør sig
gjældende. Qvaliteternes større eller mindre Disposition
for hinanden staaer her ikke nærmest i Forhold til Sidernes
Qvantum; men Foreeningen selv er en Neutralisation,
der indrømmer Qvaliteten sin Ret. Da alle
Qvaliteter ere at betragte som Resultater af Urqvalitetens
Qvalificationer (pg. 54), saa ere de alle indbyrdes
beslægtede; men da Specificationen nødvendig hører
med for at danne Qvalitetens exclusive Grændsebestemmelser,
er deres Slægtskab ikke ligestort. Affiniteten
\textbf{har sine Grader}.

Som Moment i Qvaliteternes System kan ingen
qvalitativ Side betragtes isoleret. Enhver ny Foreening
% --- p. 70 ---
\setcounter{footnote}{0}%
\opage{70}er en Opløsning af tidligere indgangne Foreeninger.
Ved A's og B's Neutralisation constitueres et immanent
Maal. Men Foreening med B forudsætter
dets Adskillelse fra C; med C danner A ogsaa et immanent
Maal, dette maa da opløses, for at hiint kan
realiseres. \textbf{Valgforbindelserne}\footnote{At Affinitetens og Valgforbindelsens Categorier gjælde ikke blot i Chemien, men ogsaa i Psychologien, har Goethe viist i „die Wahlverwandtschaften“.} fremstille saaledes
den fuldkomneste Evolution af det qvalificerede Maal.
De immanente Maal, hvis Totalitet udtrykker Maalets
hele Selvudvikling, blive ikke her paa en blot udvortes
Maade sammenlignede med hverandre; men det
ene forudsætter det andet og gaaer frem af det andet;
det enes Tilblivelse er det andets Tilintetgørelse: de
immanente Maal modificere og negere herved hinanden
indbyrdes, og Maalet ophæves.

\underafsnit{c) Den absolute Modalitet.}
\paragraf{19}{Indifferents, Differents, den prægnante Væren.}

Den høieste Forskjel og Eenhed af Eenhed og
Forskjel, hvortil Væren, som saadan, kan bringe det,
sat med det immanente Maal, hvis Selvudvikling
afslutter sig i absolut Modalitet. Som Totaleenheden
af de forskjellige Momenter bevarer det
immanente Maal alle Forskjelle i sin Eenhed, saa
at Adskillelsen af Eenhed og Forskjel er ganske ligegyldig: Modaliteten bestemmes som Værens Indifferents. Men Maalets Momenter ere selv immanente
% --- p. 71 ---
\setcounter{footnote}{0}%
\opage{71}Maal, hvoraf ethvert concentrerer alle Værensbestemmelser
i sig; paa denne Maade er Foreeningen af
Eenhed og Forskjel indifferent, og Væren opgaaer i
uendelig \emph{Differents}. Da den hele Værenseenhed i enhver Forskjel, er Forskjellen og Eenheden
overalt bestemt paa selvsamme Maade; hermed forsvinde
alle Værensbestemmelser i den \emph{prægnante} Væren.
 Det er Philosophiens Opgave at erkjende den sande Eenhed i den sande Forskjel. Den Eenhed, der har Forskjellen udenfor sig, er selv forskjellig fra sin Forskjel og derved nedsat til et Forskjelsmoment. Paa samme Maade gaaer ogsaa den Forskjel, der har Eenheden udenfor sig, i Eet med sig selv, og udgjør derved en Side af Eenheden. Den sande Eenhed maa forene begge Eenhedens Sider i sig ɔ: sættes som en Eenhed af Forskjel og Eenhed. Men den Eenhed, der negerer og excluderer Forskjellen, er, idet den fremtræder som Forskjelsmoment, netop derved forskjellig fra sig selv: den er, under sin Adskillelse fra Forskjellen, i sig selv Forskjel. Ligeledes er den abstracte Forskjel, der negerer den umiddelbare Eenhed, for at opgaae som Eenhedsmoment i den udviklede Eenhed, forskjellig fra sig selv. 
 Adskillelsen af Eenheden og Forskjellen er saaledes ikke længere umiddelbar, men resulterer af begges Selvudvikling. Medens den høiere Eenhed gjennem de umiddelbare og derved underordnede Forskjelsmomenter gaaer sammen med sig selv og prædominerer, tager Forskjellen sig af de udviklede Eenhedsmomenter tilbage, som en høiere med den concrete Eenhed coordineret Forskjel. Men naar den udviklede Forskjel er koordineret med 
% --- p. 72 ---
\setcounter{footnote}{0}%
\opage{72}den udviklede Eenhed, kan denne ikke længere prædominere,
den nedsættes da til et Forskjelsmoment for en
høiere Forskjel mellem Eenhed og Forskjel.

Denne stigende Udvikling danner en dialectisk Skala,
der i den categoriske Umiddelbarhed afsætter sine bestemte
Trin. Saalænge altsaa en logisk Sphære har
specifiske Categorier for en høiere Eenhed, er den endnu
ikke dialectisk udtømt, naar derimod den Categori,
der udtrykker alle Sphærens specifiske Bestemmelser i
høieste Eenhed, udvikles, fremtræde ogsaa Bestemmelserne
selv i deres høieste Forskjel; viser det sig da, at
Eenheden, naar den fastholdes, taber sig i Forskjellen,
og at Forskjellen, naar den souteneres, forsvinder i Eenheden,
da bliver der hverken Forskjel eller Eenhed;
Categorierne ere da utilstrækkelige, og den paagjældende
Sphære ophæves.

Værens Eenhed og Forskjel med begges specifiske
Overgange i hinanden betragtes fra Qvalitetens, Qvantitetens
og Modalitetens Synspunkt.

Den umiddelbare qvalitative Eenhed i sin abstracteste
Skikkelse er den \textbf{rene Væren} selv; men idet den
reducerer alle Forskjelle til Ikke-Væren og gjør dem til
\textbf{Intet}, bliver den selv nedsat til umiddelbart Forskjelsmoment,
nemlig en Væren forskjellig fra Intet. Det
er denne umiddelbare Forskjel, der i \textbf{Vorden} higer
efter en høiere Eenhed, efterdi enhver af dens Sider viser sig forskjellig fra sig selv. \textbf{Tilværen} er nu den
concretere Eenhed, der har overvundet det tomme Intet,
og ikke ved den rene Negation kan reduceres til Forskjelsmoment;
den er ingen umiddelbar Forbindelse af
Væren og Intet, men en gjennem Negationen medieret
Forening af Væren og Væren. Men idet Væren i
Tilværen skal bestemmes ved Væren, maa den adskille
sig fra sig selv, og Forskjellen indtræder paa ny i \textbf{No}%
% --- p. 73 ---
\setcounter{footnote}{0}%
\opage{73}\textbf{get} og \textbf{Andet}. Denne Forskjel er høiere end Forskjellen
mellem Væren og Ikke-Væren; thi begge Forskjelsmomenter
have, hvert for sig, en Tilværen. Herved er da Tilværen selv bleven Forskjelsmoment, og
Eenheden brudt. Forskjellen ligger nemlig ikke blot
i en Adskillelse af Noget og Andet indbyrdes, men er
tillige en Forskjel imellem Tilværen som saadan, betragtet
fra Eenhedens Synspunkt, og Tilværen i Noget
og Andet, bestemt ved sine egne Forskjelsmomenter,
altsaa: en Forskjel mellem Eenhed og Forskjel. Forsaavidt
Tilværen fra alle Sider gives til Priis for
Forskjellen, udtømmer den sig i \textbf{Forandringen}.

Som Forandringens Forandring resulterer Qvalitetens
prædominerende Eenhed, nemlig For-sig-væren, der
lader Noget og Andet gaae i Eet. Men idet Eet
negerer alt Andet udenfor sig, kommer det derved udenfor
sig selv; Eet sønderbrydes i Flere, og Forskjellen
fortrænger atter Eenheden under Eeternes vexelsidige
Exclusion. Da den vexelsidige Exclusion imidlertid
giennem den absolute udvikler sig til en Assimilation,
hvorunder de adskilte Flere igjen flyde sammen i Eet,
kan Forskjellen heller ikke soutenere sin Overvægt. En
rolig værende Ligevægt imellem Eenhed og Forskjel kan
Qvaliteten aldrig tilveiebringe; thi fastholdes Eenheden,
gaaer Alt i Eet, og Forskjellen forsvinder, vindicerer
man derimod Forskjellen, saa er det i de Flere forbi
med Eenheden. Her er altsaa kun en Afvexling af
Eenhed og Forskjel, en urolig Higen efter Ligevægt;
men Ligevægten selv maa søges i \textbf{Qvantiteten}.

Den Categori, der staaer paa Grændsen af Qvalitet og Qvantitet, af Forskjel og Eenhed, er Grændsen
selv. I den qvalitative Grændse er Eenheden negativ,
Forskjellen positiv; her er altsaa specifisk Forskjel
mellem Eenhed og Forskjel: den qvantitative Grændse
% --- p. 74 ---
\setcounter{footnote}{0}%
\opage{74}derimod, der, idet den gjør det Tilgrændsende til
Grændse og Grændsen til det Tilgrændsende, og i Continuitetens
Bevægelse har Discretionens Hvilepunkter,
sætter Eenheden (Continuiteten) ligesaa positiv som Forskjellen
(Discretionen); i den qvantitative Grændse er
der altsaa Eenhed af Eenhed og Forskjel. Limitationen,
der viser den allestedsnærværende Grændse, indeholder
nu den Opgave overalt at finde og sætte Eenhedens
Forskjel (ethvert Continuum er som saadant discret),
Forskjellens Forskjel (ethvert Discretum er som saadant
Continuum og derved Discretion i sig) og Forskjellens
Eenhed. Denne Forskjel og Eenhed af Eenhed og
Forskjel er given i \textbf{Tallet}. Taleenhedens Eet er just
deri forskjelligt fra det qvalitative Eet, at dette negerer
og excluderer Fleerheden, medens hiint derimod ved Fleerhedens
Discretion vindicerer sin Eenhed. Den i sig
discrete Fleerhed bevarer sin Forskjel fra Eenheden derved,
at \textbf{Antallets} Eenere specifisk ordnes; da enhver
numereret Eener er anviist sin Plads, kan den qvantitative
Fleerhed ikke saaledes som den qvalitative atter
opløses i et enkelt staaende Eet. Men Discretionens
Forskjel afbryder dog heller ikke Continuitetens Eenhed,
thi ethvert af Antallets Eetere er selv en Eenhed for
et nyt Antal. Af \textbf{Utalligheden}, Taleenhedens Uendelighed
opløst i et uendeligt Antal, udvikler sig den
prædominerende Eenhed, for hvilken Taldiscretionens
Forskjel er sat som ophævet Moment. Denne elastiske,
med alle specifike Forskjelsmomenter svangre Eenhed er
\textbf{Intensiteten}, der i \textbf{Graden} ponerer den hele udviklede Forskjel ved at ponere sig selv. Eenheden er
her sin egen Forskjel, og kommer derved til at prædominere,
men hvert Forskjelsmoment har i sig den hele
Eenhed, og er altsaa selv sin Eenhed; dette tyder da
paa Forskjellens Overvægt. Den prædominerende For%
% --- p. 75 ---
\setcounter{footnote}{0}%
\opage{75}skjel er nu ikke længere en Forskjel imellem Qvantitet og
Qvantitet, thi enhver saadan er allerede udtømt, men
en exclusiv Forskjel derimod imellem Qvantitet og Ikke-Qvantitet,
hvori det Qvalitative atter optages. At nu
Qvantiteten ikke ved denne Exclusion nedsættes til et
umiddelbart qvalitativt Moment, hvorved Dialectiken
vilde synke tilbage paa Qvalitetens overvundne Standpunkt,
forebygges i \textbf{Approximationen}, hvis sande
Udvikling er en \textbf{Evolution}. Ifølge Approximationen
indledes det exclusive Qvalitetsmoment af den bøielige,
assimilerende Qvantitet selv, og ifølge Evolutionen er
den hele qvantitative Gradation selv et Udtryk for
Qvalitetens specifiske Positioner.

Den herved indtrædende Forskjel imellem det qvalitative
og qvantitative Moment ligger alene i den
\textbf{Maade}, hvorpaa Væren i enhver af dem bestemmes.
\textbf{Modaliteten} afgiver den høiere Eenhed. Qvantiteten
kommer nemlig ikke til udvortes fra, som en vilkaarlig
Hjælpebestemmelse for Qvaliteten, men er Qvaliteten
selv, bestemt paa en anden Maade. Ligeledes fremgaaer
det af Qvantitetens egen Dialectik, at den maa bestemme
sig paa en Maade, der netop er Qvantitetsbestemmelserne,
som saadanne, modsat. Da nu den qvalitative
Bestemmelsesmaade saaledes er den qvantitative modsat,
og „Alt kommer an paa Maaden“, saa er den alsidige Forskjel realiseret, som en Forskjel imellem tvende
Classer af Værensbestemmelser. Herved opløses Modalitetens
Eenhed i tvende modale Sider, \textbf{qvantitativ
Qvalitet} og \textbf{qvalitativ Qvantitet}, af hvilke enhver
for sig er den hele Modalitet. Men idet enhver
af Siderne adskiller sig fra sig selv, er Qvalitetens
Modalitet selv Qvantitetens, og Eenheden tilstræbes i
\textbf{Modificationen}. Modificationen udtrykker just den
flygtige Forskjel i den overveiende Eenhed. Qvalitet
% --- p. 76 ---
\setcounter{footnote}{0}%
\opage{76}og Qvantitet ere begge Sider af den \textbf{samme} Væren:
de specifiske Qvantitetscategorier maae i deres fulde
Udstrækning integreres ved tilsvarende Categorier af
Qvantiteten. Begge Categoriclassers uopløselige og
specifiske Eenhed er med eet Slag sat i \textbf{Maalet}. Men
ogsaa Maalets Eenhed bestemmer sig selv til Forskjel:
begge Siderne i det immanente Maal ere, hver for sig, selv
immanente Maal. I disse Maalenes umiddelbare og
uendelige Rækker adskilles det ene Maal fra det andet
derved, at enhver af Siderne, naar den betragtes som
Maal, sættes som Eenhed for ophævede Forskjelsmomenter,
medens omvendt ethvert Maal, forsaavidt det
selv er Side, afgiver et Forskjelsmoment for en endnu
høiere Eenhed. Paa denne Maade kommer det immanente
Maal overalt udenfor sig selv. a) Det høieste immanente
Maal vilde være det, hvori alle øvrige vare
ophævede som negerede Sider; men det ligger i ethvert
Maals Dialectik selv at sættes som Side. b) Det
høieste Maal skulde, som det, hvori alle andre Maal
sættes, være det allerconcreteste; men ethvert Maal er,
som Sidernes Negation, en reen Væren, der i sin
pure Umiddelbarhed kun fremstiller sig selv: ethvert immanent
Maal er altsaa, som saadant, lige abstract
og lige concret. Det absolute Maals Forskjel fra
de relative synes derfor at maatte reduceres til den
modale Genesis, hvorved det ene Maal vorder af det
andet. c) Det høieste Maal vilde da være det, hvoraf
alle de øvrige maatte deduceres. Men det høieste Maal
er som Negation og Ophævelse af alle lavere netop selv
deduceret af disse.

Det absolute Maal er herved bestemt som den modale
Intensitet, der under en graduel Approximation
søger sit adæqvate Udtryk i de relative Maal uden dog
% --- p. 77 ---
\setcounter{footnote}{0}%
\opage{77}nogensinde i dem at finde sig selv: \textbf{det høieste Maal
er} \textbf{uopnaaeligt}.

Som det, hvis Bestemmelser aldrig kunne udtømmes
i nogetsomhelst endeligt, opnaaeligt Maal, adskiller
sig det uopnaaelige Maal fra alle relative, og sættes
derved selv som Forskjelsmoment. Men adskilt fra
Relativiteten, bestemt alene ved dennes Negation, er
det i sin For-sig-væren selv relativt. Dets ufravigelige
Bestemmelse er at specificere sit Indhold i de relative
Maal, ligesom ogsaa disse, hvert for sig, stedse
i en vis Grad udtrykke det absolute.

Den Vorden, hvorunder det absolute Maal idelig
gaaer over i det relative, og dette i hiint, er \textbf{Maalets}
egen uendelige \textbf{Modification}; ved denne stræber
det gjennem alle sine Forskjelsmomenter mod sin høieste
Eenhed. Da enhver speciel og relativ Fremtræden
af det absolute Maal kun er en Modification af samme,
og Modificationens Forskjel kun en flygtig Forskjel,
maa Modificationens uendelige Række ende i den
\textbf{absolute Modalitet}.

Maalet giver sig i sin høieste Eenhed selv sine egne
Sider; ja Maalet er selv sin egen Side, og Siderne,
hver for sig, selve Maalet. I den absolute Modalitet
er saaledes enhver Forskjel forsvunden. Alle Forskjelsmomenter,
der af den rene Væren holdtes tilbagetrængte
i Ikke-væren, ere nu efterhaanden fremtraadte af deres
Intet som forbigaaende Værensbestemmelser, men have
derved tillige viist sig i deres sande Intethed, tilsat deres
Betydning i den dialectiske Proces, og saaledes forraadt
deres Ligegyldighed. Den Eenhed, der nu fremstiller
alle Forskjelle i denne deres gjensidige Ligegyldighed,
og ved sin Ligegyldighed mod dem alle viser
sig ligegyldig mod sig selv, kan alene bestemmes som
den \textbf{absolute Indifferents}.
% --- p. 78 ---
\setcounter{footnote}{0}%
\opage{78}Enhver Værenseenhed er stedse sat som de ophævede
Forskjelsmomenters Indifferents; men en saadan Indifferents
er endnu relativ, saalænge den selv kan blive
Forskjelsmoment. Saaledes er Væren i Momentet af
den rene Væren sat som Forskjellenes Indifferents, men
Forskjellen var latent i Intet; Tilværen blev derfor
Indifferentsen af Væren og Intet, og For-sig-væren ophævede
atter Noget og Andet i Eetets Indifferents.
Den hele Værensdialectik kan betragtes som en successiv
Reduceren til høiere og høiere Indifferentser. Den
absolute Indifferents er Værens \textbf{Alt}; thi alle Værensbestemmelser
culminere i den; men Bestemmelsernes
Væren i Indifferentsen er netop deres Ikke-væren; Indifferentsen
er den negative Eenhed, hvori alle Forskjelle
ophæve hinanden, som saadan er Indifferentsen Alværens
\textbf{Nul}. Nullet er det stillestaaende Intet\footnote{See Kants Versuch den Begriff der negativen Grössen in die Weltweisheit einzuführen. Werke (G. Hartenstein) Erster Band pg. 21.}, der ved
Forskjellens Indifferentieren er udrevet af Dialectikens
Strøm. Da den absolute Indifferents er ligegyldig
mod al Forskjel, er den følgelig ogsaa ligegyldig mod
den Forskjel, der er imellem den høieste Eenhed og
dens ophævede Forskjelsmomenter, og altsaa ligegyldig
mod sig selv: \textbf{Indifferentsen er Værens logiske}
\textbf{Beredvillighed til Alt}.

Det immanente Maal er Sidernes Neutralitet, ingen
af Siderne (momentorum neutrum) bevares som saadanne
i Neutralisationen; Maalets Væren er Sidernes
Ikke-væren. Idet nu Indifferentsen, som det immanente
Maal, selv udgjør sine egne Sider, sætter den derved
sig selv tilside: \textbf{Indifferentsens Væren} er eo ipso
\textbf{dens Ikke-væren}. For ret at vise sin Indifferents
mod sig selv, bestemmer den indifferente Eenhed sig til
Forskjel, og Væren oplader sig i \textbf{uendelig Differents}.
% --- p. 79 ---
\setcounter{footnote}{0}%
\opage{79}Hermed er da givet, at Væren concentrerer alle sine
Bestemmelser i en hvilkensomhelst enkelt Bestemmelse.
Den hele Væren er selv den rene Væren, den hele
Væren er et Intet, et Noget, et Andet, Qvantitet og
Maal; \textbf{den hele Væren er} \textbf{Alt og Nul}.

Enhver Differents af Væren er følgelig selv den
hele Værenseenhed. Med den absolute Differents er
Eenheden af Forskjel og Eenhed forsaavidt prædominerende,
som enhver Differents bliver sin egen Eenhed.
Men idet hver enkelt Differents i sin Eenhed med sig
er en Indifferents af alle øvrige Differentser, og nedsætter
deres Forskjelsmomenter i Ligegyldighedens Nullitet,
maa den selv forsvinde som indifferent Moment i
enhver anden Differents, og differerer saaledes fra sig
selv. Differentsens Ligegyldighed bestaaer netop deri,
at alle Differentser ere lige gyldige. Medens Væren
altsaa under den absolute Differentieren lægger sin Totalitets
Eenhed i enhver Differents, reducerer den med
det Samme enhver Differents til Indifferents, sætter
alle Forskjelsmomenter som forsvindende, Indifferentsens
Eenhed derimod som det Blivende; der indtræder da
atter en Forskjel imellem (den ligegyldige) Forskjel og
(den blivende) Eenhed. Men da den absolute Indifferents
er ligegyldig mod sig selv, og derved opgaaer i
uendelig Differents, saa er Forskjellen mellem Eenhed
og Forskjel ligesaa indifferent som Eenheden af Forskjel
og Eenhed.

Resultatet heraf er, at Væren ikke kan være uden at
give sig en Bestemmelse, og at Værensbestemmelserne
idelig negere hverandre. Eenheden fortrænger Forskjellen,
og Forskjellen absorberer Eenheden: den ene Bestemmelse
kan ikke være i Ro for den anden. Selv
den Uro, hvori Alværen herved befinder sig, kan ikke
engang udtrykkes i nogensomhelst Vorden; thi Vorden
% --- p. 80 ---
\setcounter{footnote}{0}%
\opage{80}forudsætter en længst forsvunden Forskjel imellem Væren
og Ikke-Væren. Væren er idet den ikke er, og
er ikke idet den er: \textbf{Alværen er prægnant}.

Ifølge denne Prægnants betyder enhver Værensbestemmelse
altid tillige det Modsatte af Det, den
udsiger. Heri ligger det høieste Criterium for, at Dialectiken
har udtømt og forbrugt alle Værens Bestemmelser.
Den absolute Væren vil under den successive
Negation af alle Forskjelsmomenter gaae sammen i
Eenhed med sig selv; men bliver ved selve Negationen
af alle Værensmomenter selv uendelig negeret og reduceret
til Ikke-Væren. Men Ikke-Væren \textbf{er}, og det
endog Alværen, da det er denne selv, der har sat sig i
Ikke-Værens Moment. I sin Affirmation er altsaa
Væren negeret, og i sin Negation tillige affirmeret.
Enhver specifisk Bestemmelse, der nu skal udtrykke
denne i sin Negation affirmerede Væren, kan da ikke
længere fastholdes fra Værens Side alene, thi derved
vilde den igjen negeres og overgaae i Ikke-Væren; ligesaalidt
kan den sættes som den rene Negation af Væren,
thi enhver Ikke-Væren ophæves igjen i Væren.
Enhver sand Bestemmelse maa altsaa paa een Gang
sættes i Værens og Ikke-Værens Moment, og categorisk
udtrykke Negationens og Affirmationens Prægnants.
Men i Værens Sphære er denne Prægnants en i sig
different Indifferents og en i sig indifferent Differents;
her ere altsaa ikke de høiere Bestemmelser selv, men
kun de dialectiske Veer, hvorunder den prøvede Væren
gjenfødes i nye Categorier.
% --- p. 81 ---
\setcounter{footnote}{0}%
\opage{81}
\capitel{Andet Capitel.}{Væsen.}
\underafsnit{Indledning.}
\paragraf{20}{Det Væsentlige, det Uvæsentlige, Relativitet.}

Den prægnante Væren er Væren og Ikke-Væren i begges Eenhed og Forskjel. Idet den sætter sin Ikke-Væren som Moment for en affirmativ, bestemt Væren, og viser sig som \emph{det Væsentlige}, maa den tillige sætte Væren som et i Ikke-Væren forsvundet Moment, og derved bestemmes som \emph{det Uvæsentlige}. Saaledes souteneres en værende Forskjel i den prægnante Væren (Criticismens Standpunkt). Men al Væren opgaaer i Vorden. Forsaavidt Ikke-Væren selv er Væren, kan det Uvæsentlige fra dets Umiddelbarhed idelig hæves til det Væsentlige (Sindrighedens Standpunkt); forsaavidt al Væren derimod selv er Ikke-Væren, gaaer det Væsentlige stedse under i det Uvæsentlige (Skepticismens Standpunkt). Det Væsentliges og Uvæsentliges Sandhed ligger i begges uendelige \emph{Relativitet}, der ophæver Syns%
% --- p. 82 ---
\setcounter{footnote}{0}%
\opage{82}punkternes Fluctuation, deducerer den ene Centralbevægelse
af den anden, og saaledes reducerer den
dialectiske Udvikling til væsentligt forskjellige Sphærer.

Den prægnante Væren er en Dobbeltværen i en Værenseenhed. Denne i sig fordoblede Væren er imidlertid, ifølge Eenhedens characteristiske Prægnants, nu ikke
længere at betragte som en Væren begrændset ved Ikke-Væren,
thi en saadan qvalitativ Grændse vilde ophæve
begges Væren i hinanden. Vi have altsaa ikke her
blot en Væren, der med Ikke-Væren som Grændse bliver
til „Noget“, medens Ikke-Væren, begrændset af Væren,
vorder et „Andet“, men en Væren, hvori Væren
er, i \textbf{Eenhed} med en Væren, hvori Ikke-Væren er, og
hvori Væren altsaa ikke er. Det er netop denne I-sig-væren
af den reduplicerede Væren selv, der adskiller
den prægnante Værens Selvfordobling fra Tilværen.

Enhver Værenscategori udtrykker som saadan kun den
rene abstracte Umiddelbarhed, en Umiddelbarhed, der
igjen ophæves af den enkelte Negation. I Categorierne: „Noget“, „Eet“, „det immanente Maal“, fremstilles
Væren vel med rigere og rigere Bestemmelser,
men i enhver af disse forefindes den dog stedse ligeover
for en anden Væren, hvori den igjen maa forsvinde.
Det Characteristiske i Determinationen kan derfor overoveralt % PRINTED AS IS: over-/overalt (the syllable set twice at the line end; S099)
reduceres til: \textbf{en umiddelbar Umiddelbarhed},
\textbf{begrændset ved en umiddelbar} \textbf{Negation}.

Det dialectiske Fremskridt i Væren skeer derved, at
det tilsvarende Intet (Ikke-Væren) bliver concretere,
eftersom Værens Indhold mere og mere udvikles, men
begges gjensidige Relation opløses stedse i en Alternation,
der idelig forhindrer Væren fra at komme til sig
selv. Er Væren i en Categori fastholdt med positiv
Umiddelbarhed (som „Noget“, „Eet“, „Maal“), saa er
% --- p. 83 ---
\setcounter{footnote}{0}%
\opage{83}netop den Væren, hvorved den skulde integrere sin Bestemmelse,
sat som umiddelbar Ikke-Væren, der er saaledes
i Væren selv ingen Ikke-Væren og i Ikke-Væren
aldeles ingen Væren. Da der ikke gives nogen categorisk
Opbevarelse af Umiddelbarheden i Negationen eller
af Negationen i Umiddelbarheden, kan Umiddelbarheden
saalidt som Negationen erholde nogen sand I-sig-væren.
Den prægnante Væren derimod, der sætter en categorisk
Eenhed af Væren og Ikke-Væren, søger sine høiere specifike Bestemmelser i en \textbf{prægnant} Umiddelbarhed
\textbf{med en prægnant Negation}.

Udvikles Umiddelbarhedens og Negationens Prægnants,
saa indtræder Modsætningen imellem det Væsentlige
og det \textbf{Uvæsentlige}.

Væren er i dens abstracte Umiddelbarhed bestemt
som den rene forskjelsløse Væren, ɔ: som Væren uden
al Negation. Men denne umiddelbare Umiddelbarhed
er dog vunden gjennem en Negation, thi kun ved at
negere Forskjellene kommer man til den rene Væren.
Fastholdes denne nu i umiddelbar Lighed med sig, faaer
den Intet til Grændse, bliver derved negeret og kommer
saaledes \textbf{udenfor} sig selv. Skal den blive i sig,
og affirmere sig selv som constant Væren, da maa den
vel aabne sig for Negationen og indrømme Forskjelsmomenterne
deres Væren, men istedetfor at lade denne
Forskjellenes Væren bestemme sig som dens egen Ikke-Væren,
maa den omvendt sætte deres Ikke-Væren som
sin Væren, kun saaledes negerer den Negationen, reduplicerer
sin Umiddelbarhed og fremtræder som en Væren
i Væren, ɔ: som en \textbf{væsentlig} Væren.

Men idet Ikke-Væren har maattet afgive sit positive
Værensmoment til den væsentlige Væren, beholder den
selv kun en negativ Væren tilbage. Vel maa det erkjendes,
at det Ikke-værende ogsaa \textbf{er}, men seet i sin
% --- p. 84 ---
\setcounter{footnote}{0}%
\opage{84}egen Negation kan det ikke længere antage den rene
Umiddelbarheds Characteer og optræde som det positive
Andet ligeover for det Væsentlige. Den Væren, som
Ikke-Væren har, er netop en i sig negeret Væren. Det
Ikke-værende maa saaledes gaae for hvad det er, nemlig
for den Side af Væren, der er forladt af den affirmative Væren selv. Da dets Ikke-Væren er dets
Væren, saa er denne Væren selv en Ikke-Væren. Vi
faae altsaa Ikke-Væren i dens egen Ikke-Væren, ɔ: en
prægnant Negation. Det Ikke-værende med Ikke-Væren
til sin Væren er saaledes bestemt som det \textbf{Uvæsentlige}.

Forskjellen imellem det Væsentlige og Uvæsentlige
er hermed udviklet som en Dignitetsforskjel, men Digniteten
selv afhænger af begges Relation. Det Væsentlige
er kun med \textbf{Hensyn} til det Uvæsentlige væsentligt,
ligesom omvendt det Uvæsentlige kun med
\textbf{Hensyn} til det Væsentlige er uvæsentligt. Ethvert
af de modsatte Led fordobler sig under Relationens
videre Udvikling.

Det Væsentliges umiddelbare Væren er, som saadan,
uvæsentlig, og antager først Væsentlighedens
Characteer ved en Privation af den Væren, der ellers
vilde tilkomme det Uvæsentlige. Det Væsentlige er
altsaa kun i det Uvæsentlige væsentligt, og har saaledes
Dignitetsforskjellen i sig selv.

Det Uvæsentlige tilhører, som saadant, ikke sig selv;
isoleret fra det Væsentlige er det kun en umiddelbar
Ikke-Væren, og følgelig hverken uvæsentligt eller væsentligt.
Heri ligger: a) at det Uvæsentliges negerede
Væren hviler i et positivt Andet, med Hensyn til hvilket
det bestemmes, at det følgelig ikke selv kan gjøre
sig til, hvad det er, men maa hente sin Uvæsentlighed udenfra; b) at det Væsentlige, der giver det Uvæsentlige
sin Uvæsentlighed, herved alene characteriserer sin
% --- p. 85 ---
\setcounter{footnote}{0}%
\opage{85}egen Væsentlighed, at denne Væsentlighed, i sig uvæsentlig,
som Følge heraf slet ikke er andet end selve
Uvæsentligheden i det Uvæsentlige; altsaa: Uvæsentligheden er det Uvæsentliges egen Væsentlighed.

Relationen maa da nærmere bestemmes saaledes, at
det Selvsamme, der i det ene Led er væsentligt, er i
det andet uvæsentligt, og da hvert af Leddene refererer
sig til sig selv som til sin Negation, saa er det selvsamme
Led i een Henseende væsentligt, i en anden
Henseende uvæsentligt. Dignitetsforskjellen er aldeles
afhængig af det forskjellige Hensyn (diversus respectus).
Men intet Hensyn kan finde Sted, uden at der
maa være Noget, hvortil man seer hen. Dette Noget
kan hverken søges i det Væsentlige eller i det Uvæsentlige,
da det netop er denne Dignitetsforskjel, der skal
bestemmes ved Hensynet selv. Det, hvortil man seer
hen, maa derfor repræsentere sig som et umiddelbart
Givet, som en abstract, for begge Led gjældende Væsentlighed.
Er f. Ex. A en saadan given Væsentlighed,
med Hensyn til hvilket a sættes som det Væsentlige,
b derimod som det Uvæsentlige, da er A, som Det,
hvorpaa Dignitetsforskjellen imellem de underordnede Led
beroer, det Eeneste, der i denne Relation kan kaldes
det sande Væsentlige. Men er A det Væsentlige, da
er det, i og for sig betragtet, tillige baade væsentligt
og uvæsentligt; dets hypothetiske Dignitet afhænger da
af et modsat Led B, der ligeledes i sig kan være baade
væsentligt og uvæsentligt. Dignitetsforskjellen imellem
A og B maa altsaa bestemmes ved en endnu høiere
Væsentlighed o. s. f.

\textbf{Forsaavidt} B altsaa er det Uvæsentlige, er A det
Væsentlige, og den deraf deducerede Dignitetsforskjel
imellem a og b staaer fast; \textbf{forsaavidt} B derimod
sættes som det Væsentlige, bliver A det Uvæsentlige,
% --- p. 86 ---
\setcounter{footnote}{0}%
\opage{86}og de underordnede Dignitetsmomenter altereres; forsaavidt
endelig baade A og B, ifølge høiere Hensyn,
hver paa sin Maade anerkjendes som væsentlige, blive
deres underordnede Momenter snart meer snart mindre
væsentlige, og Modificationen gaaer i det Uendelige.

\textbf{Criticismens} characteristiske Eensidighed bliver heraf
indlysende. Den critiske Methode forudsætter en væsentlig
Forskjel imellem det Væsentlige og Uvæsentlige;
den er saa langt fra at confundere disse Værensdigniteter,
at dens hele Bestræbelse meget mere gaaer ud
paa at skjelne og sondre dem indbyrdes. Men idet
Criticismen opfatter denne \textbf{væsentlige} Forskjel blot
som umiddelbart \textbf{værende}, og ved Forstandens Demarkationslinier
overalt søger at betegne den qvalitative
Grændse, maa den idelig forandre Standpunkt, og forvilder
sig saaledes i en Labyrinth af modsigende Hypotheser
og forskjellige Hensyn.\footnote{Disse „forskjellige Hensyn“ ere, saa at sige, den ydre Reflexions Kastevinde, der, hvis de faae Overhaand, maae bringe enhver philosophisk Idee til at kuldseile. Det forholder sig med den critiske Philosophi, som med den critiske Videnskab overhovedet, den siger altid: Forsaavidt — forsaavidt! Men Forsaavidtphilosophien er en Gammelmandsphilosophi, der trods alle \textit{reservationes mentales} maa \textit{crepere} af Skrupulositet. De forskjellige Inddelinger af \textit{Articuli fidei fundamentales} et \textit{non fundamentales} kunne, hvad det kirkelige Lærebegreb angaaer, afgive et Exempel.}
 Det næste Fremskridt er, at de forskjellige Hensyn
med alle deraf følgende Værensdigniteter opfattes i
deres hele gjensidige Relativitet, hvorved da enhver
umiddelbar Bestemmelse, hvad enten den i sin rene
Umiddelbarhed repræsenterer sig som væsentlig eller uvæsentlig, ophæves. Denne Ophævelse af al overfladisk
% --- p. 87 ---
\setcounter{footnote}{0}%
\opage{87}Umiddelbarhed characteriserer \textbf{Sindrighedens} og Aandrighedens
Standpunkt, hvor man i Alt søger en dybere
Betydning (ὑπόνοια).

Paa den ene Side har man altsaa her det hele Indbegreb
af Tilværelsesmomenter, der alle ere uvæsentlige,
fordi de alle ere umiddelbare; paa den anden Side
derimod den rene Væsentlighed, der, idet den giver alle
Bestemmelser deres væsentlige Betydning, derved tillige
bestemmer sig selv. Opgaven er da, at eftervise de
umiddelbare Momenters Relation til Væsentligheden;
thi i denne Relation ere de alle væsentlige. Paa
Sindrighedens Standpunkt, hvor selv det Ubetydeligste
erholder sin Betydning, ophæves altsaa den \textbf{værende}
Forskjel imellem det Væsentlige og Uvæsentlige. Methoden
gaaer ud paa at negere den abstracte Umiddelbarhed,
der som saadan er eet med selve Uvæsentligheden;
men denne Negation er en Omtydning, hvorved
ethvert Moment aldeles taber sin uvæsentlige Side, og
Væsentligheden selv identificeres med Indbegrebet af de
relative Væsensmomenter. Sindrigheden er endnu hildet
i det critiske Dilemma: enten maa A være væsentligt
eller uvæsentligt. Væsenets Væren i de relative
Momenters Totalitet er altsaa dog endnu en umiddelbar
Væren; man tilveiebringer kun en overfladisk Immanents,
uden at bringe det til Bevidsthed, at Væsentligheden
nødvendig maa have Uvæsentligheden, som saadan,
i sig. Eensidigheden forraader sig derfor ogsaa
paa en dobbelt Maade: deels kan man nemlig kun
udvikle Væsentligheden ved en successiv Omtydning af
de uvæsentlige Momenter, hvorved Udviklingen taber
sig i en uendelig Vorden og Approximation, der aldrig
naaer til den absolute Væsentlighed selv; deels bliver
den ved Omtydningen fremstillede Væsentlighed saaledes
neddragen i Momenternes Relativitet, at den i dem
% --- p. 88 ---
\setcounter{footnote}{0}%
\opage{88}mister al væsentlig Betydning, og Aandrigheden ender
i Aandløshed.\footnote{Exempler herpaa have vi i den overfladiske Teleologie, den smagløse Typologie og den vilkaarlige Allegoriseren.}

Er nemlig det Uvæsentlige selv blevet væsentligt og
sat som Væsentlighedens umiddelbare Udtryk, da er
hermed ogsaa det Væsentlige blevet uvæsentligt. Med
Uvæsentlighedens Forsvinden forsvinder den flygtige
Negation, som den, der alene vedkommer den udvortes
Betragtning, og kun den rene abstracte Umiddelbarhed
staaer tilbage; men denne er netop Uvæsentlighedens
sande Characteristicon.

\textbf{Skepticismen} viser os, hvorledes det umiddelbart
Væsentlige stedse gaaer over i det Uvæsentlige; enhver
affirmativ Væren opløser sig i de skeptiske „Troper“.
α) Den rene Umiddelbarhed, Tingen, som saadan, \textbf{er}
kun \textbf{for} en Bevidsthed; Tingens Væren affirmerer
da ikke sig selv, men Affirmationen tilhører Selvbevidstheden, der er Tingens Ikke-Væren. Da det
altsaa ikke er Tingen selv, men kun den subjective
Paastand om Tingen, der har Betydning, saa er
enhver Ting, som saadan, \textbf{uvæsentlig}: „Mennesket
er Tingenes Maal“.

Men Bevidsthedens affirmative Væren characteriserer
sig derved, at den kan negere en hvilkensomhelst
bestemt Væren: dens Selvaffirmation ligger
just i den negative Operation, hvorved den negerer
Alt, ogsaa sig selv. Selvbevidsthedens Væsentlighed
bestaaer altsaa deri, at den kan sætte sig selv som
\textbf{uvæsentlig}.
β) Bliver den Negation, hvorved Bevidstheden negerer
sin egen og alle Tingenes Væren, igjen negeret, da
antager den negative Operation saaledes vel et
% --- p. 89 ---
\setcounter{footnote}{0}%
\opage{89}positivt Anstrøg; men den væsentlige Umiddelbarhed,
hvori den af Negationen restaurerede Væren gaaer
sammen med sig selv, er dog kun den ubestemte og
Alt bestemmende Menen, den rene Postuleren. Men
enhver Postuleren maa postulere Noget. Forsaavidt
Noget postuleres, negeres Andet. Da nu Noget
alene postuleres, fordi det postuleres, og Andet
negeres blot, fordi det negeres, saa kan man omvendt
ligesaa godt negere Noget og postulere Andet.
Skal Grunden, hvorfor Noget affirmeres, ikke ligge
i Affirmationen selv, da maa den søges i et Andet,
som allerede er affirmeret; men skal denne Affirmation
nu heller ikke være grundløs, da forudsætter
den igjen et Andet, og saaledes i det Uendelige.
Dogmatismens Axiomer ere alle grundløse og som
Følge heraf \textbf{uvæsentlige}.
Reducerer man endelig Affirmationen af Noget og
Andet til deres vexelsidige Forhold, saaledes at Be%
vidstheden affirmeres, fordi den affirmerer Tingen,
og Tingen igjen affirmeres ved Bevidsthedens Affirmation, da indtræder en „Diallelus“, der gjør Bekræftelsen afhængig af Det, der selv afhænger af
Bekræftelsen.

Striden mellem disse forskjellige Standpunkter hidrører
fra en eensidig Opfattelse af den prægnante Værens uendelige \textbf{Relativitet}.
a) Relativiteten i al Tilværen forudsætter en For-sig-væren. At Noget er relativt, betyder, at det ligesaa
meget er for sig selv, som for Andet. Denne Side
bliver fastholdt i Criticismen.
Den sande Relativitet udtrykker de relative Momenters
Væren for hinanden indbyrdes. Heri ligger
deels en gjensidig Affirmation, deels en gjensidig Negation.
% --- p. 90 ---
\setcounter{footnote}{0}%
\opage{90}Sindrigheden opfatter Relationen fra det affirmative
Synspunkt; hvad her negeres er kun Momenternes
isolerede For-sig-væren. Den ene umiddelbare
Bestemmelse gaaer sammen med den anden,
og idet alle Værensmomenter knyttes og støttes til
hverandre indbyrdes, consoliderer Alværen sig selv
som den positive Væsentlighed. Sindrigheden erkjender
Værens prægnante Umiddelbarhed, men bemærker
derimod ikke, at det ene umiddelbare Værensmoment
just, som saadant, er det andets Negation,
og at den sande Væren idelig forsvinder, idet
man vil gribe og fastholde den.
Skepticismen tager Relativiteten fra dens negative
Side; for den er Indbegrebet af alle Værens Momenter
kun en Totalitet af gjensidige Negationer.
Umiddelbarheden forsvinder, hver Gang den dukker
op, og med den forsvinder ogsaa enhver Affirmation.
Skepticismen har kun Die for den prægnante Negation, % PRINTED AS IS: Die (D for Ø, Øie; S104)
hvorunder alle Bestemmelser gjensidig fortære
og opløse hverandre, men den føler ikke Væsentlighedens
uendelige Nærværelse i alle negative
Bevægelser, og kan derfor aldrig naae til den
sande Værens categoriske Udtryk.
Den absolute Relativitet fremstiller os derimod den
prægnante Umiddelbarhed i den prægnante Negation.
Medens den skeptiske Relativitet kun udtrykker
sig i en ydre Relation imellem negative Momenter, saa at dens Negation er successivt uendelig,
uden at være absolut, concentrerer den absolute Relativitet
derimod alle negative Momenter i en negativ
Eenhed, hvori enhver negeret Bestemmelse refererer
sig til alle øvrige, som til sig selv, og hæver
sig derved til affirmativ Umiddelbarhed. Ifølge
Relativiteten er den ene umiddelbare Værensbestem%
% --- p. 91 ---
\setcounter{footnote}{0}%
\opage{91}melse den andens Negation; heri ligger: α) at den umiddelbare Væren \textbf{er} Negation, og altsaa det Negative af sig selv; β) at dens Negativitet \textbf{er} selve
det Andet, hvori den negeres; γ) at dette Andet
ogsaa \textbf{er} det Negative af sig.

Alle negative Momenter \textbf{ere} altsaa, nemlig som
negerede, og det ene er, hvad det andet er, thi det er
en Fælledsbestemmelse for dem alle, at ethvert for sig
af dem er det Negative af sig selv. Enhver negativ
Bestemmelse er altsaa i sin Negativitet sig selv og kun
sig selv alene; denne dens Relation til sig selv udgjør
Umiddelbarhedens første Hovedside. Dernæst er det negative
Moment i sin Negativitet netop det Negative
af sig selv, det er selv alle andre negative Momenter. Det sætter sig derfor ikke istedetfor Andet og excluderer
Andetheden, saaledes som Eetet i For-sig-væren; Andetheden
\textbf{er} i det negative Moment som negeret; men den
negerede Andethed er Andetheden i sin egen Andethed,
altsaa Andetheden i sin Sandhed: det negerede Moment
er selv sit Andet, efterdi det er et Andet af sig selv.
Denne Relativitet udtrykker Umiddelbarhedens anden
Hovedside.

Men denne prægnante Umiddelbarhed, der er Negativitetens
\textbf{egen Umiddelbarhed}, er følgelig tillige
\textbf{Umiddelbarhedens egen Negativitet}.

Den negative Umiddelbarhed er den Værenseenhed,
hvori alle negerede Momenter synke sammen som i deres
Almeennegation. Almeennegationen er altsaa en Negation
af alle negative Momenter og, som saadan,
deres affirmative Eenhed. Som den Eenhed, der affirmerer
de negative Momenter, affirmerer den vel sit eget
Andet; men da dens Relation til sit Andet er dens absolute
Relation til sig selv, saa affirmerer den sig selv,
idet den affirmerer sin egen Negativitet. Enhver umid%
% --- p. 92 ---
\setcounter{footnote}{0}%
\opage{92}delbar Værensbestemmelse, der i den kan svinde, er
altsaa allerede svunden, og enhver, der kan komme, \textbf{er}
allerede kommen; som Negation af al Kommen og Svinden
er Eenheden Alværens positive Centrum, dens absolute
I-sig-bliven. Denne centrale, i Negationen luttrede
Væren er Alværens \textbf{Væsen}. Som Negativitetens Umiddelbarhed er Væsenet Negativitetens
Affirmation i Negativiteten selv, enhver anden
Væren, der viser sig med reen Umiddelbarhed, maa da
ligesaa umiddelbart vise sin Umiddelbarhed negeret; enhver
umiddelbar Væren vil da fremdeles blive at betragte
som Umiddelbarhedens egen Negativitet, som en
Væren, der umiddelbart udtrykker sin Ikke-Væren, ɔ:
som \textbf{Skin}\footnote{Man tænke sig her t. Ex. Tænkerens absolute Relation til sig selv igjennem sine Tanker.}.
 Det vil heraf være indlysende, at den prægnante
Umiddelbarhed og Selvaffirmation, som Væsenet vindicerer
gjennem den prægnante Negation, kun kan sættes
og udvikles, idet alle Tilværelsens Momenter, som
Skin og Værensreflexer, ophæves i Væsenets Dialectik.
a) Væsenets Selvbestemmelse fuldbyrder sig altsaa i en
Cirkelbevægelse, hvorunder det optager alle Værens Momenter: \textbf{Væsenhedens Dialectik maa have}
\textbf{sin særegne Sphære}. b) Det Samme gjælder om Skinnet, Væsenets væsentlige
Negation. Uagtet Skinnets Væren maa deduceres
af Væsenet, har det dog en Væren; just fordi
Væsenets Væren er den originale, og Skinnets den
deducerede, kan Skinnets egen Dialectik kun udvikle
sig ved at sætte Væsenheden som negeret Moment;
Skinnet ophæver saaledes overalt Væsenet, og 
be%
% --- p. 93 ---
\setcounter{footnote}{0}%
\opage{93}væger sig i sig selv som Phænomen. \textbf{Derved} erholder
\textbf{Phænomenet sin logiske Sphære}. c) Men idet Væsen og Phænomen, hvert i sin Sphære,
afrunde sig til dialectiske Totaliteter, viser det sig,
at disse kun ere relative. Deres sande Dignitetsforskjel
maa søges i deres absolute Relativitet.
Ifølge den absolute Relativitet refereres Totaliteten
til sig selv som til sit Andet, og til sit Andet som
til sig selv. \textbf{Totaliteternes gjennemgribende
Eenhed}, \textbf{der vil overvinde den} \textbf{væsentlige} Dualisme,
\textbf{udfolder sig} i \textbf{Virkelighedens Sphære}.

\afsnit{A. Væsenhedens Sphære.}
\underafsnit{a) Det rene Væsen.}
\paragraf{21}{Væsen, Skin, Reflexion.}

Det \emph{rene Væsen}, ɔ: den absolute Værenseenhed,
der ophæver al Forskjel, idet den negerer alle
umiddelbare og relative Værensmomenter, er, i Kraft
af denne Almeennegation, nærmest at bestemme som
det absolute Intet. Men Værens umiddelbare Bestemmelser tilintetgjøres kun i det absolute Intet
derved, at de tilintetgjøre hverandre indbyrdes, idet de
altsaa sættes i deres Intet, sættes de i deres Sandhed. Den endelige Væren bliver herved erkjendt for,
hvad den egentlig er, nemlig Negativitetens egen
flygtige Umiddelbarhed, det rene \emph{Skin}. Identiteten,
den væsentlige Eenhed af Væsen og Skin, er Forskjellens
Eenhed, og Modsætningen, begges væsentlige For%
% --- p. 94 ---
\setcounter{footnote}{0}%
\opage{94}skjel, er Eenhedens Forskjel: \emph{Reflexionen} udtryk- Væsenets Selvbevægelse og Gjenspeilen i alle % PRINTED AS IS: udtryk- (the word broken at the line end and not continued: udtrykker; S108)
Skinnets Momenter.

Uagtet Væren nu i foregaaende Capitel har adskilt
sine categoriske Bestemmelser, og derved efterhaanden er
bleven opløst i en Mangfoldighed af limiterede og modificerede
Momenter, saa er det dog med alt dette den
\textbf{ene og samme} Væren, der, ved at negere sin umiddelbare
Liighed med sig selv, har sat sine Forskjelsmomenter,
og, ved en successiv Negation af enhver specifisk
Bestemmelse, ophævet alle sine umiddelbare Categorier,
for at gaae sammen med sig selv i absolut Eenhed.

I den \textbf{rene Væren} supprimeres alle specifiske Værensbestemmelser,
og deres Plads indtages af det rene
Intet: i den \textbf{absolute Indifferents} udvikles de forskellige
Bestemmelser paa een Gang til Ligegyldighed
og til lige Gyldighed. Forsaavidt nu deres Ligegyldighed
er deres Ugyldighed, fremstiller Indifferentsen sig
som Værens blinde Eenhed; men, forsaavidt Momenterne
derimod i deres Ligegyldighed hvert for sig have
samme Gyldighed, gives Indifferentsen til Priis for
uendelig Differents, og Væren opnaaer aldrig nogen
sand Bestemmelse. I det rene \textbf{Væsen} derimod ere
alle Værensmomenter ikke blot udviklede og satte, men
tillige negerede og ophævede; deres umiddelbare Gyldighed
bliver herved bestemt som deres sande Ugyldighed,
og Væsenet sætter sig selv som det absolut Gyldige, ɔ:
som den sande Værenseenhed. Da Væsenet, som den absolute Væren selv, ikke kan
forandre sig, saa følger heraf, at det heller ikke umiddelbart
kan være „Noget“; thi ethvert „Noget“ staaer lige
imod et „Andet“, og er som saadant „Forandringen“
underkastet. Ligesaalidt kan Væsenet fastholdes som en
% --- p. 95 ---
\setcounter{footnote}{0}%
\opage{95}umiddelbar For-sig-væren, hvis qvalitative Eenhed skulde
excludere alle relative Værenselementer, thi derved maatte
disse, som for-sig-værende, igjen excludere hiint. Ifølge
en saadan vexelsidig Exclusion vilde da Væsenet kun
fremstille sig som „Eet“ iblandt „Flere“, og saaledes
blive inddraget med i Kredsen af de endelige Relationer.

Væsenets eneste Bestemmelse bliver da den, at negere
enhver Bestemmelse, hvorved det da atter synes
at repræsentere sig som den rene Væren. Men den
væsentlige Forskjel imellem Væren og Væsen er ikke
destomindre, naar den nærmere betragtes, iøinefaldende.
Medens nemlig den rene Væren i sin logiske
Uskyldighed endnu ikke veed af nogen Negation, og
derfor, om vi saa maae sige, overraskes af det begrændsende Intet, for i sin hele Ubestemthed at bestemmes
af dette, har Væsenet derimod optaget Negationen og
med den Principet for alle Bestemmelser i sig. Væsenet
negeres ikke af Noget, da det meget mere selv
negerer Alt; det \textbf{lader} sig ikke bestemme, men er selv
det, hvorfra al Bestemmelse udgaaer. Væsenets Bestemthed er altsaa at sætte som dets absolute Selvbestemmen.

Men uagtet Væsenet, som den i sig concentrerede
Væren, nu betragtes som det logiske Væld, hvoraf
alle Værens Momenter emanere, og følgelig selv maa
afgive Bestemmelser i det Uendelige, saa er denne dets
Bestemmelighed dog ingenlunde at forvexle med Indifferentsens Beredvillighed til at antage enhver Bestemmelse. Indifferentsen hjemfalder til Differentsen,
og tager, saa at sige, tiltakke med Værensmomenterne,
som det kan falde sig, just fordi den maa \textbf{lade} sig bestemme.
Den har vel Negationen under alle særegne
Skikkelser i sig, men den magter ikke selv denne Negation;
dens Negation altererer derfor stedse dens Vær%
% --- p. 96 ---
\setcounter{footnote}{0}%
\opage{96}en, og forhindrer den saaledes fra al Selvbestemmelse.
Væsenet derimod har Negationen selv i sin Magt, og
magter derfor ogsaa alle relative Værensbestemmelser.
Enhver umiddelbar Væren, der forsøger paa at undslippe
den centrale Eenhed og affirmere sin egen Umiddelbarhed,
bliver da strax indhentet og overmandet af
Væsenets uimodstaaelige Negation, der nøder den til
at bekjende sin Bevægelses Hvorfra? og Hvorhen? ? Den
værende Umiddelbarhed vorder af sit Intet; den kommer altsaa fra det Negative af sig selv, ɔ: fra Væsenet; under sin Vorden gaaer den derfor ogsaa tilbage
i sit Intet, ɔ: i Væsenets Negation.

Da Væsenet negerer enhver umiddelbar Væren, maa
det følgelig ogsaa negere sin egen Værens Umiddelbarhed.
Ved sin tveæggede Negation tilvinder det sig
vel en prægnant Umiddelbarhed, men dens tvende Sider udgjøre stedse een eneste, reent umiddelbar Eenhed.
Ifølge denne Eenheds rene Umiddelbarhed er Væsenet
under sin Selvbestemmelse idelig i Begreb med at blive
til „Noget“ og at repræsentere sig som et Værende. Er nu ethvert umiddelbart Noget hjemfaldet til Væsenets Negation, da maa denne Negation ogsaa idelig
ophæve Væsenets egen Væren. Den absolute Væren
bliver saaledes sat som det i sig absolute Intet, og
den positive Side bestemmer sig selv som den negative.
Er Væsenet i sig Negativitetens egen Umiddelbarhed,
da er det eo ipso ogsaa Umiddelbarhedens egen Negativitet:
\textbf{det rene Væsen er}, \textbf{som saadant}, \textbf{det
rene Skin}.

Det relative, endelige Intet, der kun negerer et vist
Noget, er i sig ligesaameget en Væren som et Intet; det
fremstiller sig derfor ogsaa som et tilværende Intet, og
antager Andethedens positive Præg; men det i sig absolute
Intet, der ikke fordrager nogensomhelst umiddelbar
Væren, kan derimod aldrig selv indtræde i Tilværelsen.
% --- p. 97 ---
\setcounter{footnote}{0}%
\opage{97}Resultatet synes da at være dette, at Væsenet, ifølge
sin Selvbestemmelse, maa frakjende sig selv enhver Værensbestemmelse,
at det altsaa overhovedet slet ikke er,
og at kun den endelige Væren, seet i sin egen Uendelighed
d. e. i Rigdommen af alle sine Momenter, er
den sande Væren.\footnote{I dette Resultat acqviescerer den atheistiske Logik; thi Atheismen er, logisk taget, Væsensfornegtelse.}
 Men er Væsenet uden Tilværelse, da maa Tilværelsen
ogsaa være væsenløs, og, som saadan, fremstille
sig i sin hele Uvæsentlighed. Tilværens dialektiske
Udvikling godtgjør vel, at Noget selv maa være sit
Andet, og at det derved bliver til noget Andet, end
det forhen var; men denne Ophævelse af Nogets Umiddelbarhed
er paa Værens Standpunkt kun en Afvexling
med en følgende Umiddelbarhed. Forandringens
Negation er altsaa blot en Fremskriden fra den ene
Umiddelbarhed til den anden. Hvad her negeres, er
kun det givne, bestemte Umiddelbarhedsmoment; Negationen
er derfor partiel og successiv. Det Noget, der
forgaaer, tilintetgjøres kun som \textbf{dette} eller hiint bestemte
Noget. Men ethvert Noget, der følger, maa
stedse bestemmes som \textbf{dette} eller hiint: det „\textbf{bestemte}
Noget“ bliver saaledes under hele Forandringen egentlig
uforandret. Hermed er da Forandringens Væren
ophævet og sat som et Skin. Hvorvidt Forandringen
skal gjælde, beroer alene paa, hvorvidt Noget umiddelbart
kan adskille sig fra Andet. Saalænge der er en
værende Forskjel imellem Noget og Andet, saalænge er
der ogsaa en Vorden og Overgang fra hiint til dette.
Men ifølge den prægnante Negation er Nogets umiddelbare
Væren egentlig dets Ikke-Væren, dets Væren
som Ikke-Noget. Noget er i sig et Skin; deraf 
kom%
% --- p. 98 ---
\setcounter{footnote}{0}%
\opage{98}mer det, at dets Overgang i Andet ingen Overgang
er. Nogets Andet \textbf{er} da, som Følge heraf, heller ikke
noget egentligt Andet, men kun det Gjenskin af Andet,
hvorved Noget viser ud over sig selv. Idet Andet
nu kaster sit Skin paa Noget uden selv at opgaae
deri, forbeholder det sin rene umiddelbare Væren for
sig selv. Andet er altsaa nu det egentlige Noget, det
Reale, der synes at have reddet sig ud af Skinnets Intethed. Men fastholdes nu Andet som Noget, da er det, ifølge det Foregaaende, atter et Skin i sig og et
Gjenskin af Andet. Saaledes bliver Tilværelsen med
al dens Mangfoldighed gjennemsigtig, og dens Umiddelbarheder
klare sig i utallige Gjenskin. Hvert Skin
i dette Endelighedens Rige er en flygtig Umiddelbarhed
, der i sin Negation viser ud over sig selv til et
Værende; men i det Værendes Sted indtræder idelig
et nyt Skin, og den affirmative Væren bortflygter i
det Uendelige. I Skinnet selv er der ikke Noget, og
ethvert Noget udenfor Skinnet forvandler sig til Skin;
saaledes staaer Tilværelsen som den absolute Tomhed:
„Alles Seyn ist Schein“.\footnote{Hermed er Akosmismens Standpunkt angivet. Akosmisten og Atheisten kunne forene sig i Erkjendelsen af Verdens Forgjængelighed. Men de opfatte denne paa modsat Maade. Atheisten betragter Alt under Værens og Forandringens Synspunkt, og kan derfor vel indrømme Tingenes Forkrænkelighed i det Enkelte; men da de partielle Umiddelbarheder under Forandringen idelig restaureres, urgerer han, hvad det Hele angaaer, Verdens evige Aprioritet og Uforkrænkelighed. Akosmisten gjennemskuer derimod den hele Tilværelses Intethed, idet han fra Væsenets Standpunkt fatter Alværens Indhold som Skin af Skin og Verden som Illusionernes Rige.}
 Den Intethed, hvori Alværen nu synes at forflygtiges,
maa imidlertid selv forsvinde som Skin, naar Væ%
% --- p. 99 ---
\setcounter{footnote}{0}%
\opage{99}sensnegationen gjennemføres. Det Forfærdende i Nihilismens Tankebevægelse ligger deri, at den forudsætter
en Dignitetsforskjel imellem den umiddelbare Væren paa den ene og det væsenløse, illusoriske Skin
paa den anden Side. Den søger uophørlig en Væren, der i sin Sandhed og Reenhed skulde være frigjort
fra ethvert Skin, og idet den da overalt finder, at en
saadan Væren kun kan bestemmes som Skinnets rene
Andet, og at ethvert saadant Andet, som Skinnets
Negation, er, naar man vil fastholde det, et Gjenskin
af Skinnet og i sig selv et Skin, fører den Bevidstheden
til det trøstesløse Resultat, at Alt \textbf{kun} er Skin.

Man vrager Skinnet, fordi man stoler paa Værens
rene Umiddelbarhed. Indlyser det derimod, at det netop
er den umiddelbare Væren, som, naar den nærmere prøves, selv er Skin, da har Væren sin Sandhed i Skinnet; Skinnet staaer da, som saadant, over den
blotte Væren, Dignitetsforholdet vender sig om, og
man maa tilstaae, at enhver Umiddelbarhed, der skal
ligge udenfor Skinnet, \textbf{kun} er en Væren. Skal nu den sande Væren søges i Skinnet, da er
hermed givet, at enhver anden Væren, der ikke er ophævet
i Skinnet, just derfor maa betragtes som en
egentlig Skinværen. Den væsentlige Adskillelse imellem
den sande og usande Væren vil derfor nu blive at
sætte som en \textbf{Adskillelse imellem det sande} \textbf{og
falske Skin}.

Det rene Skin er for det Første en værende Umiddelbarhed,
der kun repræsenterer sig selv; men som flygtig
Umiddelbarhed, ɔ: som Umiddelbarhedens egen Negativitet, maa det dernæst ligesaa umiddelbart slaae
over i sin Negation. Dette pludselige Omslag, der
overalt gjentager sig i Skinnets uendelige Række, medfører
for den naive, godtroende Bevidsthed altid en
% --- p. 100 ---
\setcounter{footnote}{0}%
\opage{100}Overraskelse. Den, der ikke aner, at Umiddelbarheden
selv er Negativitet, vil, naar han seer Skinnets rene
Umiddelbarhed, deri fængsle og fastholde den positive
Væren; men idet nu Skinumiddelbarheden strax forraader
sin Negations Hemmelighed, seer han sig skuffet,
og vil da reent overspringe Skinnet for at gribe Realiteten.
Det er denne uendelige Selvadskillelse, under
hvilken Skinnet idelig paa Ny gjentager sig, der frembringer
Illusionen: „\textbf{Skinnet bedrager}“.

Det enkelte Skin henviser i sin Selvnegation til et
andet Skin; Umiddelbarhedens og Positivitetens Bortflygten
viser sig som en uafbrudt Overgang fra Skin
til Skin; saaledes ligger Skinnets Affirmation i dets
Negation. Men idet Skinnet ved sin Negation afspeiler
og bekræfter sig selv, bliver det fuldstændig ophævet;
thi det strider imod Skinnets Bestemmelse kun
at affirmere og gienspeile sig selv alene. Det absolute
Skin, der gjennemskinner alle Vexelskin, er da tillige
Skinnets sande Andet, det Reale i Skinnet, hvis Umiddelbarhed
ikke er flygtig, men positiv og blivende; saaledes
er \textbf{Væsenet}, Skinnets \textbf{absolute Negation} i
\textbf{Skinnet selv}, \textbf{gjennemklaret som det rene} \textbf{Lys-} (τὸ

Lyset er Skinnets Eenhed; denne Eenhed er vunden
giennem Skinnets uendelige Selvadskillelse, den er altsaa
Forskjellens Eenhed og som saadan \textbf{Identitet}.
Som Skinnets Affirmation i dets Negation er Identiteten
Skineenhed og selv Skin; men som Skinnets
Negation i dets Affirmation, som dets Ophævelse i
Andet, er den Væsenseenhed, ja Væsenet selv, og Forskjellen
Skin. Da Identiteten saaledes baade er Skin
og Væsen, saa er den paa een Gang Forskjellens Negation
og tillige selv Forskjel. Begge disse Momen%
% --- p. 101 ---
\setcounter{footnote}{0}%
\opage{101}ter maae da poneres og negeres i det uendelige Vexelskin
imellem Væsen og Skin ɔ: i \textbf{Reflexionen}.

Reflexionens alsidige Udvikling characteriserer dens
trende Hovedskikkelser: 1) Som \textbf{Bevægelsens} Reflexion
udtrykker den Identitetens og Forskjellens uendelige
Overgang i hinanden; 2) som \textbf{Bestemmelsens
Reflexion} fixerer den Identitetens og Forskjellens Modsætning, og endelig 3) som den \textbf{absolute} Reflexion
viser den Bevægelsens og Bestemmelsens sande
Vexelskin.

1. Bevægelsens Reflexion er a) \textbf{sættende}, b) forudsættende og c) \textbf{modsættende}.

a) Umiddelbarheden, den rene abstracte Identitet, negerer sig selv, og forsvinder som Skin, idet den
repræsenterer sig som sin egen Forskjel; det Umiddelbare gaaer saaledes over i det Negative. Denne
Overgang er imidlertid ingen Vorden; hvad Umiddelbarheden
herved \textbf{bliver}, nemlig det Negative
af sig selv, det \textbf{var} den allerede inden Overgangen skete. Her begyndes da ikke med det Umiddelbare,
men Reflexionsumiddelbarheden \textbf{resulterer}
som Identitetsmedium imellem tvende negative
Extremer. Overgangen skeer altsaa fra det Negative
til det Negative, hvorved Negativiteten refererer
sig til sig selv, saa at Umiddelbarheden
kun er den Eenhed, hvori Relationen fastholdes.
Forskjellen, det Negative, der fremgaaer af den
umiddelbare Identitet, er netop Det, hvoraf Identiteten eller den værende Umiddelbarhed selv er
fremgaaet. Det er Umiddelbarhedens Negation,
der ophæver Identiteten; men da Umiddelbarheden,
som Skin, selv er denne Negation, saa ophæver
Identiteten igjen Forskjellen. Det Umiddelbares
Tilbagevenden til sig selv gjennem Negationen er
% --- p. 102 ---
\setcounter{footnote}{0}%
\opage{102}selve Negationens Tilbagevenden til sig gjennem
Umiddelbarheden. Umiddelbarheden sættes, idet
den resulterer af den ophævede Negation. Men
den sættes ikke af det abstract Negative, thi ogsaa
dette er resulterende, og bliver selv \textbf{sat} som det
Umiddelbares Ophævelse.

Hvad der i denne Overgang \textbf{sættes} og ophæves,
er altsaa et flygtigt Moment, et Skin,
der forsvinder, idet det kommer, og kommer igjen
i samme Nu, som det forsvinder; Væsenet derimod
er den uendelige \textbf{Sætten} og \textbf{Ophæven}, og som
saadant Bevægelsen og Reflexionen selv.

Som den rene Sætten er Væsenet en Bevægelse
fra Skin til Skin, og følgelig en Bevægelse, der
ikke bevæger sig ɔ: en Skinbevægelse, der selv
sættes af det Ubevægede. Idet Væsensidentiteten
gjennem Ophævelsen af Skinnets ponerede Væren
vil affirmere sig selv, opgaaer den i en væsentlig Forskjel; den bliver nemlig paa een Gang baade
det Ubevægede, der sætter, men ikke sættes, og
det Negative heraf, nemlig den sættende Bevægelse,
der selv sættes. Som det Ikke-satte, men Sættende, er Væsenet
Bevægelsens \textbf{Princip}, Bevægelsen derimod med
Alt, hvad den afsætter, Skin og \textbf{Resultat}. Men
som Bevægelsens Princip, maa Principet udtrykke
sin Væsenhed i Bevægelsen selv; det fremgaaer
saaledes af sin egen satte og ophævede Bevægelse,
og er ifølge heraf selv Resultat. Principet og
Resultatet ere nu identiske; men da Principet kun
er Princip ved at \textbf{sætte}, og Resultatet kun Resultat
ved at \textbf{ophæve} Bevægelsen, saa er Bevægelsen
selv væsentlig, og udtrykker den sande
Forskjel imellem Princip og Resultat.
% --- p. 103 ---
\setcounter{footnote}{0}%
\opage{103}Med Principet \textbf{begynder} den sættende Bevægelse,
med Resultatet \textbf{ender} den. Da Principet
selv er Resultat, maa Bevægelsen ende i samme
Nu, som den begynder, men da Resultatet selv er
Princip, maa den ogsaa begynde i samme Nu,
som den ender. Væsenets rene Sætten er saaledes
en uendelig Omsætten. Forud for ethvert
Resultat gaaer der en Sætten, der negeres i Resultatet selv; men da Resultatet herved omsættes
til Princip, saa fremgaaer den ophævede Sætten
igjen af dens egen Negation og fortsætter sig selv.

Som den Sætten, der allerede maa være sat,
for at kunne sættes, og desaarsag sætter sin egen
Ophævelse, for at sætte sig selv, er Reflexionen
\textbf{forudsættende}.
b) Den forudsættende Reflexion bestemmer Principets
og Resultatets væsentlige Identitet og Forskjel.
For at Principet kan sætte sit Resultat, maa det
allerede selv være. Men i denne abstracte Forsig-væren
er det kun et Skin; dets sande Væsen
ligger i den Sætten, hvorved det sætter Resultatet. Før det har sat sig selv, er det altsaa ikke
væsentlig Princip, og før Principet er blevet væsentligt,
kan det intet Resultat afsætte; men omvendt,
kan det heller ikke sætte sig selv, uden derved,
at det afsætter sit Resultat. I Principets og
Resultatets Væsen ligger saaledes en dobbelt Forudsætning.
α) For at sætte sig selv, sætter Principet sin egen Negation og Ophævelse, nemlig sit
Resultat, men det ophæver selv denne Negation;
thi det sætter kun sig selv ved at sætte sit Resultat,
som det, der endnu ikke er sat, men skal sættes.
Ifølge denne \textbf{primitive} Forudsætning sætter Principet sig selv \textbf{forud} ved at forudsætte
% --- p. 104 ---
\setcounter{footnote}{0}%
\opage{104}sit Resultat. Resultatet er forudsat som et
Principets Væsen forsvindende Skin, og derved
forudsat, som det, der skal sættes; med denne
Sætten indtræder Principets væsentlige Negation,
og Resultatet sætter sig selv. Men Principets
Negation er et væsentligt Moment i dets Selvposition; Resultatet sætter altsaa sig selv som det,
der ikke sætter men sættes; det forudsætter sit
forudsættende Princip; dets sande Væsen er dets
Selvophævelse, og dets Forudsætten \textbf{consecutiv}\footnote{Efterdi Gud forud har sat sin Verden (primitiv Forudsætten), saa forudsætter Verden nu ogsaa sin Gud (consecutiv Forudsætten).}. γ) Denne i sig fordoblede Forudsætten er imidlertid
ingenlunde at betragte som en Composition af
tvende heterogene Fragmenter. Den primitive og
consecutive Forudsætten ere identiske. Men deres
Identitet er heller ingen Tautologi; ethvert af
Reflexionens tvende Led er selv den hele Reflexion.
De ere begge lige selvstændige, og i denne deres
gjensidige Selvstændighed hinandens fuldstændige Negation og Forskjel. Den udviklede Forudsætten
er altsaa een eeneste i sig redupliceret Sætten,
der uafbrudt vender sig \textbf{imod} sig selv, og derved
characteriserer sig som en \textbf{Modsætten}. Idet den modsættende Reflexion fordobler sin egen
Selvstændighed, afsætter den tvende Centra, og
gaaer saaledes gjennem en dobbelt Bevægelse tilbage i sig selv. Principet tager sig gjennem Resultatet selv tilbage, og skriver, saa at sige, sit
Resultat som en Peripheri om sit eget Centrum. Det er saaledes paa een Gang sættende, sat, for%
% --- p. 105 ---
\setcounter{footnote}{0}%
\opage{105}udsættende og forudsat, og dets i sig fuldendte
Bevægelse udgjør dets Ubevægelighed. Omvendt
reflecterer Resultatet sig selv i sit Princip, dets
Selvophævelse er dets Selvposition. Men Resultatets
Selvposition er Principets væsentlige Negation:
Resultatet reflecterer altsaa sin Ubevægelighed
i sig, og hviler i det ophævede Princip som
i sig selv alene. Hermed er den absolute Identitet
sat som den absolute Forskjel. Principet er
nu ikke længere Princip, thi dets Overgang i Resultatet
er tilendebragt; ligesaa lidt kan Resultatets
Reflexion have nogen specifisk Betydning, da
Resultatet nu meget mere selv er Princip. Ethvert
af de modsatte Led er altsaa i sig afsluttet,
Bevægelsen forsvunden, og kun den rene Modsætning
staaer tilbage\footnote{Den væsentlige Forskjel imellem Pantheismen og Atheismen, Dualismen og Deismen paa den ene Side og Theismen paa den anden, viser sig som en Forskjel imellem den eensidige og alsidige Opfattelse af Reflexionsbevægelsen. Den rene Reflexion beskriver en Cirkel. Pantheismen adskilles nu fra Atheismen (Pankosmismen) derved, at hiin gjør Gud, denne derimod Verden til Cirkelens Centrum. Men den eensidige Reflexion uden væsentlig Modsætning er en abstract Uvæsentlighed, og den enkelte Cirkel et stort Nul. Dualismen (Deismen) holder sig derimod til Modsætningen. For Dualisten repræsenterer Tilværelsen sig under Billedet af tvende hinanden tangerende Cirkler. Men da man her forglemmer, at den forudsatte Urmodsætning forud er sat, altsaa ikke uoprunden, men udsprungen af en dybere Identitet, saa kan Dualismens dobbelte Centrum ogsaa siges at repræsentere den rene Ubegribelighed. Kun Theismen opfatter den sande Reflexionsbevægelse i dens hele Selvfordobling. Idet Guds og Verdens Centrum erkjendes som udsprunget af det guddommelige Væsens egen absolute Centralitet, ophæver den theistiske Reflexion Dualismens Cirkler: ⊙⊙ i Uendelighedens Symbol ∞.}.
% --- p. 106 ---
\setcounter{footnote}{0}%
\opage{106}2. I den forudsatte Modsætning er ethvert Led selvstændigt,
i sig \textbf{bestemt}, og maa tages, som det er
givet. Hermed forsvinder Bevægelsen som Skin i Bestemmelsens
Reflexion.

Bestemmelsesreflexionen er a) \textbf{distingverende}, b) antinomisk
og c) \textbf{fixerende}.

a) For den distingverende Reflexion repræsenterer
Modsætningen sig i sin rene Umiddelbarhed, og
fastholdes derfor som en Samværen af tvende hinanden tangerende Identiteter: A = A, B = B.
Idet nu Reflexionsbestemmelsen sættes som den i
sig acqviescerende Identitet, forsvinder Bevægelsen
som Skin; den sættes herved som Noget, der er
Identiteten selv aldeles uvedkommende, som en
Sammenstillen og Sammenlignen, der blot foregaaer
i Betragterens subjective Bevidsthed, altsaa
som en \textbf{ydre} Reflexion. Betragterens Die hviler % PRINTED AS IS: Die (D for Ø, Øie; S121)
snart paa A og snart paa B, og den subjective
Reflexion indprænter sig nu deres Forskjel: A er
ikke B; B er et Ikke-A; A er et Ikke-A; A er ikke Ikke-A.
I A's rene Identitet med sig selv ligger da aldeles
ingen Bestemmelse af nogensomhelst Ikke-A.
Ja selve Sætningen: A er A, udtrykker kun en
Reflexionsidentitet, der falder udenfor A selv. A
er nemlig her sat tvende Gange; men i sin
rene identiske For-sig-væren kan A ikke adskille sig
fra sig selv. Adskillelsen skyldes altsaa den ydre
Reflexion, der under sin Anstrængelse for at fastholde
Identiteten i dens fuldkomne Bestemthed,
og \textbf{distingvere} imellem A og Ikke-A (B), har
% --- p. 107 ---
\setcounter{footnote}{0}%
\opage{107}maattet sætte A imod, og ved Copula „er“
igjen gjøre denne Modsætning ugjort. Den rene
fra al Modsigelse frigjorte Bestemthed kan alene
udtrykkes som en Interjectionsidentitet: A!, og
den Substitution, hvorved den subjective Reflexion
indsætter A istedetfor A (A er? --- A!), er da herved
reduceret til en fuldkommen Tautologi\footnote{Den tautologiske Reflexion, hvis Criterium er, at den betragter Modsigelsen som en subjectiv Tankefeil, der kan hæves ved en Distingveren imellem de værende Identiteter, har fundet sin vidtløftigste Udvikling i Wolfianismen. Wolf definerer Contradictionen som en ”\textit{simultaneitas in affirmando et negando}”, uden Anelse om, at han herved definerer den sande Reflexions eget Princip. Idet han beskriver Identiteten: ”\textit{Eadem dicuntur, quæ sibi invicem substitui possint}”. (\textit{Ontolog, p. 148}) tilføier han: ”\textit{Facta nimirum substitutione perinde esse debet, ac si nulla facta fuisset}”, uden at observere, at enhver saadan Substitution just derved bliver aldeles betydningsløs.}.

Men selv denne Tautologi, der bestræber sig
for at undflye Modsigelsen, har sit egentlige Livsprincip
i Modsigelsen selv. Den subjective Reflexion
gjennemgaaer nemlig under sin Bevægelse følgende Sætninger: A er (maaskee) Ikke-A; A er
ikke Ikke-A og endelig A er A. At den sætter A istedetfor A har alene Betydning derved, at den
ellers maatte sætte Ikke-A istedetfor A. Men
saaledes sætter den jo netop A istedetfor (det befrygtede)
Ikke-A. Kun ved Modsigelsen selv
kan den subjective Reflexion distingvere sig fri
for Modsigelse. Modsigelsen er altsaa ikke længere
nogen Tankefeil, men en Tankenødvendighed,
% --- p. 108 ---
\setcounter{footnote}{0}%
\opage{108}og Reflexionen er hermed erkjendt som \textbf{antinomisk}. Ved den antinomiske Reflexion bringes det til
Bevidsthed, at Dualiteten, Modsætningen og Modsigelsen
høre med til Reflexionens eget Væsen;
men desuagtet bliver dog den Reflexionsbevægelse,
hvoraf Bestemmelserne resultere, sat som et Skinvæsen,
hvis Sætninger maae ophæves, saasnart de
ere satte; det sande Væsen holder sig derimod i
sin rene ubestemmelige Eenhed, og ethvert Forsøg
paa at bestemme det tilintetgjøres i Reflexionens "Antithetik". \emph{Thesis}: \textbf{Det sande} Væsen er, naar
\textbf{det nærmere bestemmes}, \textbf{den absolute} Identitet.

Den rene Identitet er \textbf{forskjellig} fra Forskjellen
og Modsætningen \textbf{modsat}. Herved har den
ophævet sin abstracte For-sig-væren og selv sat sig
som Forskjellens Forskjel og Modsætningens
Modsætning; den absolute Identitet er altsaa
i sig selv Modsætning. \emph{Antithesis}: \textbf{Det sande Væsen maa}
\textbf{bestemmes som den absolute} \textbf{Modsætning}.

I enhver Modsætning maa det ene Led nødvendigviis
forudsætte det andet; thi for at være
modsatte, maae de begge være \textbf{satte} mod hinanden:
Modsætningen er de modsatte Leds Inclusion. Men da ethvert af Leddene i Modsætningen
er det andet modsat, maae de gjensidig excludere
hinanden; Modsætningen er altsaa Leddenes gjensidige
Exclusion. Inclusionen og Exclusionen ere
selv hinanden modsatte. Saaledes er da Modsætningen
fra alle Sider i Modsætning med sig
selv, og derved, som absolut Modsætning sig selv
% --- p. 109 ---
\setcounter{footnote}{0}%
\opage{109}modsat. Men det Modsatte af Modsætningen er
Identitet\footnote{Den antinomiske Reflexion er characteristisk for Kantianismen.}

Det i sig blivende og i sin rene Bestemthed
ubestemmelige Væsen er imidlertid ikke at forvexle
med Værens absolute Indifferents. Den indifferente
Væren er idelig i Begreb med at antage en
hvilkensomhelst Bestemmelse, og befinder sig saaledes
i en uafbrudt Bevægelse og Uro. Det i
sig værende Væsen holder derimod alle Reflexionsbestemmelser
ude, sætter dem af sig selv som
Skin og Afglands, og acqviescerer saaledes i sin
rene \textbf{Transcendents}, opløftet over alle Modsætninger.

c) Naar denne Adskillelse af Identitet og Modsætning, Bestemmelse og Bevægelse, paa ethvert
Punkt gjennemføres, og derved drives til den
høieste Spidse, viser Reflexionen sig som \textbf{fixerende}. Den fixerende Reflexion fastholder paa det Strængeste
Væsenets rene Væren i dens Forskjel fra al
Negation. Selv Reflexionsbestemmelserne maae i
Conseqvents heraf fixeres som ubevægelige. Den
i sig fixerede Reflexionsbestemmelse\footnote{Den fixerende Reflexion er med en sjelden Energi gjennemført i den Herbartske Philosophi.} er nu at
sætte som en for-sig-værende Væsenhed, hvorved
Væsenets Uridentitet knuses, og adskilles i en Utallighed
af Identitetsfragmenter. Ethvert af disse
Fragmenter er selv et I-sig-værende, en reen Identitet, en Tautologi, altsaa en Modsigelse. Tautologiens Allestedsnærværelse er ogsaa Modsigelsens
Allestedsnærværelse; enhver specifisk Bestemmelse,
% --- p. 110 ---
\setcounter{footnote}{0}%
\opage{110}den subjective Reflexion tillægger det Bestemte,
maa den derfor idelig tage tilbage. Væsenets
ubevægelige Momenter vise sig for Reflexionen i
en uafbrudt Bevægelse, uagtet denne idelig forsikkrer,
at de ikke bevæge sig; de gaae under deres
gjensidige Bestemmelse idelig over i hverandre,
og Reflexionen maa selv anvende den enes Negation
for at bestemme den anden, uagtet det dog
paastaaes, at det ene Identitetsfragment („Reale“)
ikke vedkommer det andet. Reflexionen hilder sig
her i et uendeligt Selvbedrag; idet den saa eensidigt fixerer sine Bestemmelser, fixerer den sig selv.

3) Denne Illusion ophæves i den \textbf{absolute eller
alsidige} Reflexion, der, som a) \textbf{optisk-intuitiv} b) \textbf{dialektisk} og c) \textbf{dialektisk-speculativ}, lader Bestemmelsen
og Bevægelsen gjenspeile hinanden. Det transcendente Væsen, der svæver over Reflexionens Modsætning, bliver derved Modsætningen selv modsat, og saaledes i sig bestemt som
Modsætningens absolute Identitet. Væsenet afspeiler
nu sin Ubevægelighed i sin Bevægelse, det
skuer sig selv i alle sine Reflexionsbestemmelser,
og dets dunkle Transcendents forvandler sig ved
den optisk-intuitive Reflexion til den klare \textbf{Immanents}\footnote{Schelling.}.
Men den i Immanentsen anskuede
Bestemthed er, som den er, og det indlyser endnu
ikke, hvorfor den maa bevæge sig; Bevægelsen
skues som absolut, og det erkjendes endnu ikke,
hvorfor den maa afsætte det Bestemte. Identiteten
af Bevægelsen og Bestemmelsen er altsaa
endnu tautologisk umiddelbar, og Modsigelsen indlysende.
% --- p. 111 ---
\setcounter{footnote}{0}%
\opage{111}b) Denne i Bevægelsens og Bestemmelsens Identitet
immanente Modsigelse kommer først til sand Udvikling
i den \textbf{dialectiske} Reflexion.

Den absolute Identitet er, som saadan, selv
Identitet og Modsætning, og den udviklede Modsætning
ligeledes paa een Gang baade Modsætning
og Identitet. Ethvert Led i Modsætningen
er altsaa for sig den hele Identitet og den hele
Modsætning. Det soutenerer sin Integritet ved
at udelukke det andet Led; Identiteten excluderer
Modsætningen, just fordi den selv er Modsætning
og har denne i sig. Da ethvert af de modsatte
Led nu har sin Selvstændighed i sig ved at have sit Andet, og kun har dette Andet ved at excludere
det, saa excluderer det sin Selvstændighed i
samme Nu, som det vil vindicere den; det skrider
ud over sig selv og kommer i Bevægelse, just
som det vil blive i sig og fixere sin Bestemthed.
Ifølge denne immanente Modsigelse har ethvert
af Leddene dets Ulighed ophævet i dets egen Lighed,
og er saaledes at sætte som det \textbf{Positive},
men omvendt har det ogsaa en i Ligheden reflecteret
Ulighed, og maa derfor ligesaavel sættes som
det \textbf{Negative}\footnote{Dieß ist also derselbe Widerspruch, der das Positive ist, nemlich Gesetztseyn oder Negation, als Beziehung auf sich. Aber das Positive ist nur an sich dieser Widerspruch; das Negative dagegen der gesetzte Widerspruch; den in seiner Reflexion in sich, an und für sich Negatives oder als Negatives identisch mit sich zu seyn, hat es die Bestimmung daß es Nichtidentisches Ausschliessen der Identität sey. Es ist dies gegen die Identität identisch mit sich zu seyn, hiermit durch seine ausschließende Reflexion sich selbst von sich auszuschliessen. Das Negativeist also die ganze, als Entgegensetzung auf sich beruhende Entgegensetzung, der absolute, sich nicht auf Anderes beziehende Unterschied; er schließt als Entgegensetzung die Identität von sich aus; aber somit sich selbst, den als Beziehung auf sich bestimmt er sich als die Identität selbst, die er ausschließt. Hegel.}.
% --- p. 112 ---
\setcounter{footnote}{0}%
\opage{112}c) Men naar det Positive i dets Udvikling selv er det Negative, saa er Identiteten af den udviklede Positivitet
og Negativitet, nærmere betragtet, den udviklede
Tautologi selv. Modsigelsen er da endnu ikke
løst, og kan umuelig forsvinde, saalænge Bestemmelsen
selv endnu ikke er nærmere bestemt. Det Positive maa da sættes imod sig selv som et Dette! imod
Dette! Men denne specifiske Antithese kan Reflexionen,
som saadan, ikke fatte; thi for den rene Reflexion
er Dette = Dette, som A - A. Den fixerede
Bestemthed culminerer i Umiddelbarheden, d. e. i
Reflexionsbevægelsens Negation. Det demonstrable
Positivitetsmoment maa skues, som det er sat af
Væsenet selv. Det sandses imidlertid ikke, som et
umiddelbart værende Noget, dets specifiske Positivitet
er tillige dets specifiske Negativitet, det er, som
det Ubevægelige, et Moment med i den dialectiske
Bevægelse selv. Integreret ved dette Bestemmelsesmoment
er Reflexionen \textbf{dialectisk-speculativ}. Idet Væsenet, som uendelig Negativitet, afsætter
de positive, specifiske Bestemmelser, og derved
tillige sætter dem som Reflexionsbevægelsens negative,
ubevægelige Momenter, blive de thetiske
Forskjelle opbevarede og afspeilede i Reflexionen
selv. De thetiske Momenters dialectiske Udtræden
af Væsenets pleromatiske Centrum constituerer saaledes
den reale Forudsætning, og reflecterer den
absolute Væsenheds egen speculative Grundtypus.
% --- p. 113 ---
\setcounter{footnote}{0}%
\opage{113}
\underafsnit{b) Væsenets Dualitet.}
\paragraf{22}{Det primitive Væsen, det consecutive Væsen, Subordination.}

I den absolute Identitet af alle positive Reflexionsbestemmelser har Væsenet sin uendelige \emph{Fylde}. Som Fyldens negative Eenhed (transcendent Hypostase)
maa det afsætte Bestemmelsens specifiske Momenter; men da disse svinde som Skin i dets egen
Selvposition (transcendental Hypostase), forbliver Fylden esoterisk, og den deri reflecterede Modsætning
ophæves i det \emph{primitive} Væsens egen, for-sigværende Identitet (absolut Hypostase). Det primitive
Væsens uendelige Positivitet er tillige dets
uendelige Negativitet; den esoteriske Fylde bliver derfor
negeret i den exoteriske, og Forskjellen indtræder
i alsidig Specification. Den exoteriske Fyldes selvstændige,
i sig reflecterede Eenhed er herved sat som
det primitive Væsens egen centrale Modsætning ɔ: som \emph{consecutivt} Væsen. Men idet den esotoriske % PRINTED AS IS: esotoriske (for esoteriske; two readers)
Fylde under den i sig reduplicerede Forudsætten subordinerer
sig den exotoriske, bliver det primitive Væsen % PRINTED AS IS: exotoriske (for exoteriske; S137)
selv det consecutive subordineret, og med denne vexelsidige \emph{Subordination} gaaer den centrale Modsætning
tilgrunde.

Tilværelsen har opløst sig i sit Intet; den svandt
som Skin og Skygge i det rene Væsen. Væsenet i
denne dets Abstraction var selv en negativ Væren;
% --- p. 114 ---
\setcounter{footnote}{0}%
\opage{114}men af dets uendelige Negativitet udviklede sig tillige
Skinnets Negation og Ophævelse; Væsenets Selvreflexion
constituerede sig som Princip for alle Reflexionsbestemmelser,
og den bestemte Bestemthed culminerede i
den værende Umiddelbarhed. Hermed er altsaa Tilværelsen
reproduceret som den af Væsenet \textbf{satte} Væren.\footnote{In der Sphäre des Seyns war das Daseyn das Seyn, das die Negation an ihm hatte, und das Seyn der unmittelbare Boden und Element dieser Negation, die daher selbst die unmittelbare war. Dem Daseyn entspricht in der Sphäre des Wesens das Gesetztseyn. Es ist gleichfalls ein Daseyn, aber sein Boden ist das Seyn, als Wesen oder als reine Negativität; es ist eine Bestimmtheit oder Negation nicht als seyend, sondern unmittelbar als aufgehoben. Daß Daseyn ist nur Gesetztseyn dieß ist der Satz des Wesens, vom Daseyn. Hegel. Logik IItes Buch. Pag. 23.}
Tilværelsens hele Rigdom af specifiske Bestemmelser
\textbf{er} altsaa, og maa tages, som den \textbf{er} givet; men den er
ikke af sig selv; som umiddelbart værende, er den uendelig
Fragilitet og Selvopløsning; dens Vorden er en
uafbrudt Resulteren, og den selv et Resultat, der forudsætter
sit forudsættende Princip; ifølge denne consecutive Forudsætten ophæves alle Tilværelsesmomenter i
Urvæsenets negative Eenhed. Men som negativ Eenhed
er Væsenet ikke et abstract Medium eller udvortes
Receptaculum, hvori Momenterne kunne tænkes opbevarede
som værende. Deres Væren i Væsenet er netop
deres Ikke-Væren; de ere i Urvæsenets Selvposition
kun satte, som de, der skulle sættes.\footnote{Denn man kann nicht sagen, daß in Gott sey Feuer, Bitter oder Herbe, vielweniger Luft, Wasser oder Erde, allein man siehet, daß es daraus worden ist. Jacob Böhme.} Forsaavidt
% --- p. 115 ---
\setcounter{footnote}{0}%
\opage{115}Væsenet bestemmes som det, \textbf{der} sætter sig selv ved at
sætte Tilværelsens uendelige Mangfoldighed, er det
med denne prægnante Identitet reflecteret i sig som i
den absolute Fylde
(πλήρωμα). Væsenet er her consecutivt forudsat uden dog
egentlig at være sat; det er sat af den satte Væren,
altsaa af Det, der Intet kan sætte, efterdi det selv er
berøvet al Selvstændighed. Den Reflexionsbevægelse,
som Væsenet altsaa selv har at gjennemgaae, for at
kunne afsætte Fyldens specifiske, i deres For-sig-væren
reflecterede Momenter, skal da nu nærmere betragtes: a) Da alle Reflexionsbestemmelser forsvinde i Væsenets Ur-Eenhed, saa er denne Eenhed reflecteret i
sig som i Bestemmelsernes Negation; Væsenet affirmerer sig selv ved at negere alle Bestemmelser,
dets Ubestemmelighed er dets i sig reflecterede Bestemthed.
Men i denne bestemte Ubestemmelighed
er dets Væren uendelig affirmativ, og Væsenet er
herved forudsat som den \textbf{transcendente} Hypostase.
Idet Væsenet hypostaseres, ophæver det
den consecutive Forudsætten, og bestemmer sig selv
til at være det Altsættende. Herved ere da Reflectionsbestemmelserne
forudsatte, som de, der endnu
ikke ere satte. Men de skulle sættes; thi det
transcendente Væsens Ubestemmelighed er alene
bestemt som Principet for alle Bestemmelser.
b) Forsaavidt den transcendente Hypostase nu altsaa
sætter sig selv, maa den bestemme sine Bestemmelser,
og derved afsætte Differentsen og Mangfoldigheden,
som Det, hvori den selv negeres; men
forsaavidt den skal sætte sig som det Altbestemmende,
maae Bestemmelsernes Reflexer opgaae som
forsvindende Momenter i dens egen Selvposition;
kun saaledes kan Væsenet tilegne  sig alle Re%
% --- p. 116 ---
\setcounter{footnote}{0}%
\opage{116}flexionsbestemmelser, assimilere dem med sig, og
gjøre dem til sine egne. Men herved transcenderer
det transcendente Væsen tillige sin rene abstracte For-sig-væren, og oplader sin forborgne Fylde,
for at fastholde sig selv som Reflexernes bestemte
Identitet. Ved denne Reflexionsbevægelse
er Væsenet sat af sig selv, og har saaledes bestemt
sig som den i sig bestemmelige ɔ: transcendentale
\textbf{Hypostase}. Transcendentaliteten
ligger nu, umiddelbart taget, ikke Mere, end hvad
der alt maatte være givet i Transcendentsen selv;
men, ifølge Reflexionsbevægelsen, dog noget Mere.
I den rene Transcendents vare nemlig alle Bestemmelser
forsvundne, og Væsenet consecutivt forudsat; den transcendentale Hypostase ere de derimod
satte som Reflexer i Væsenets egen Selvposition,
og Væsenet har fuldbyrdet den Bevægelse, hvorved
det constituerer sig selv som det Altsættende.
Hvad der altsaa laae forborgent i Transcendentsens Involution, bliver klaret og sat i Transcendentalitetens
Evolution. Ifølge den Selvbevægelse,
der indtræder imellem det Transcendente og
Transcendentale, forholder Væsenet sig nu til sig
selv, som Princip til Resultat. Den transcendentale
Hypostase forudsætter den transcendente, som
Resultatet forudsætter Principet. Men Principet
er kun Resultatets Princip. Den transcendente
Hypostase med sin forborgne Væren maa derfor
sætte sine Reflexionsbestemmelser som den Selvnegation
og Modsætning, i hvilken den først kommer
til at besidde sig selv og sin Fylde. Begge Hypostaser
ere altsaa, hver i sig og hver for sig, det
hele Væsen med dets hele Fylde, og ved denne
deres For-sig-væren reflecterede i hinanden som
% --- p. 117 ---
\setcounter{footnote}{0}%
\opage{117}gjensidige Negationer. Hermed \textbf{er} Urvæsenets prægnante Identitet i sig redupliceret som Urmodsætning
og \textbf{Dualitet}.
c) Denne Dualitet i Urvæsenet er imidlertid ikke at
fastholde som en blot værende Forskjel imellem
Væsen og Væsen. Selvbevægelsen er ingen Vorden,
ingen successiv Overgang. For at Væsenet
kan være et Transcendent, maa det allerede have
sat sig som et Transcendentalt: Principet er selv
Resultat, og Resultatet Princip. Den Negation
og Forskjel, der fjerner Tautologien, ligger i Fyldens
Differentser og Modsætninger; men da alle
Momenter endnu kun ere forsvindende Skin i Væsenets
Selvudvikling, bevæger Væsenet sig fra sig
selv til sig selv, og sætter Bestemmelsens Reflexer
som det Medium, hvori det afspeiler sig selv alene.
Det i sig udviklede Væsen er saaledes paa een
Gang baade sættende og sat, og fastholder i sin
\textbf{absolute Hypostase} den hele Fylde som esoterisk.

Den absolute Hypostase forudsætter den udviklede
Forskjel imellem Transcendentsen og Transcendentaliteten;
den er saaledes hverken den ene
eller den anden af disse; men holder sig i sin
For-sig-væren som Dualitetens og Modsætningens
Identitet. Hvis den kun var Transcendentsens og
Transcendentalitetens udviklede Tautologie, forsvandt
igjen al Bestemmelse; Væsenet blev da
atter forvandlet til \textbf{et} transcendent Abstractum, og
Bevægelsen maatte begynde forfra; men hvis den
absolute Hypostase omvendt satte den transcendente
og transcendentale udenfor sig, og, saa at sige, \textbf{ved}
Siden af sig, da vilde Urvæsenets Selvaffirmation
opløses i lutter Relativiteter; det vilde derved re%
% --- p. 118 ---
\setcounter{footnote}{0}%
\opage{118}præsentere sig som en sat Tilværelse, der atter
maatte ophæves og negeres. Væsenets Selvposition
er da at sætte som den absolute Identitet i
den absolute Modsætning. Hvert af Modsætningens
Led er selv den hele Identitet, derved er
Identiteten sig selv modsat, men denne Modsætning
er Identitetens Negation i Identiteten selv, altsaa
dens egen Selvbestemmelse, hvorved Modsætningen
ophæves i den for-sig-værende Identitet. Saaledes
er Urvæsenets absolute Eenhed reflecteret i
dets udviklede Triplicitet (τρεις ὑποστασεις ἐν μια οὐσια),
og Væsenet har affirmeret sig selv
som \textbf{det primitive} Væsen.

Som transcendent Hypostase er det primitive
Væsen et blot \textbf{værende} Princip, et væsentligt
Skin, der reflecterer sin egen Forskjel: som transcendental
Hypostase er det derimod \textbf{sat} som Princip,
og er dermed sit eget Resultat, og endelig som
absolut Hypostase constituerer det sig som det \textbf{sig
selv sættende}, i sig fuldendte og væsentlige
Princip.

Men det \textbf{væsentlige} Princip er kun det væsentlige
Resultats Princip; det reflecteres derfor og
fuldendes i sig, som det, der endnu ikke er fuldendt.
Som det sig selv sættende Princip vindicerer
det primitive Væsen sin egen centrale Identitet
med sig; denne Centralitet ligger deri, at enhver
Hypostase for sig kan sættes som Eenhed af de
tvende andre; idet den transcendente saaledes sættes
som absolut, reflecterer den i sig den transcendentale
som transcendent og den absolute som transcendental.
Den Selvbevægelse, hvorved det primitive Væsen sætter sin egen Identitet og Forskjel,
er altsaa, for sig tagen, en Skinbevægelse,
% --- p. 119 ---
\setcounter{footnote}{0}%
\opage{119}og kun som primitiv Forudsætten absolut væsentlig. Reflexionen, hvorved det væsentlige Princip
reflecteres i sig, maa derfor fra alle Sider
fordobles, og derved vise sig som en Reflexion,
hvorved det reflecteres i Andet. Urmodsætningen
mellem den transcendente og transcendentale Hypostase,
hvoraf den centrale Identitet resulterer, maa
altsaa selv reflecteres i en central Modsætning,
hvorved det primitive Væsen negeres i det ikke-primitive. Urvæsenets Selvnegation er saaledes i sig redupliceret:
deels er den nemlig en Negation, hvorved
Væsenet sætter sig selv som sit eget Negativum, hvad der, som ovenfor er viist, frembringer
Urmodsætningen i det primitive Væsen selv; deels
er den en Negation, hvorved det primitive Væsen,
som concret for-sig-værende Identitet, sætter sit
Ikke-Væsen ɔ: sit væsentlige Andet. Uden en saadan
central Modsætning er Urmodsætningen et
tomt Skin, og den centrale Identitet et transcendent
Abstractum; men, reflecteret i sin Centralmodsætning,
er Urmodsætningen selv den primitive Forudsætning,
og det primitive Væsens centrale Identitet det sig selv sættende Princip, der i den centrale
Modsætning afsætter sit væsentlige Resultat.

Denne det primitive Væsens centrale Modsætning
kan nu hverken søges i Tilværelsen, da denne
selv er forsvunden som Skin, ei heller i det rene
Skin, da dette kun er Negation af det rene Væsen,
som selv er Skin. Det primitive Væsens
væsentlige Negation maa derfor sættes som det i
sig reflecterende Skin ɔ: som det uvæsentlige eller
 consecutive Væsen. 
% --- p. 120 ---
\setcounter{footnote}{0}%
\opage{120}Det rene Skin er uendelig Forskjel; det affirmerer
sig selv som Skin af Skin i uafbrudt Selvadskillelse.
Det i sig reflecterede Skin gjenspeiler
altsaa en Utallighed af Skin, hvoraf hvert enkelt
igjen er et i sig reflecteret Skin, en Skinhypostase.
Forsaavidt nu Skinbevægelsen overalt gaaer ud
paa en Hypostaseren, antager ethvert Skin et Væsenheds-Præg,
og vindicerer sig en relativ Selvstændighed;
men forsaavidt det væsentlige Skin
kun er et Skinvæsen, opgaaer dets hypostatiske
Affirmation som Moment i dets væsentlige Negation:
det consecutive Væsens Reflexion-i-sig hensvinder
overalt som uendelig Reflexion-i-Andet;
idet det herved negerer den centrale Identitet, og
forvandler sig til en Skinidentitet for de ydre Reflexionsbestemmelsers
hypostaserede Forskjelle, sætter
det sig som det udadkastede Væsen med en exoterisk
Fylde.

Den esoteriske og den exoteriske Fylde negere og
reflectere hinanden indbyrdes. I den esoteriske kan
nu ethvert Forskjelsmoment forsvinde som Skin,
just fordi Forskjellen selv poneres og hypostaseres
i den exoteriske, og omvendt kan Identiteten i den
exoteriske sættes som Skinidentitet, efterdi den esoteriske
Fyldes Identitet er affirmeret og centraliseret
i det primitive Væsen selv.

Centralmodsætningen mellem det primitive og
consecutive Væsen fremstiller sig nu ogsaa som en
sand Modsætning imellem det \textbf{Positive} og Negative.
Det primitive Væsen, der i sin abstracte
For-sig-væren vilde nedsynke til et flygtigt Skin
uden al Selvstændighed og Reflexion-i-sig, befrier
sig selv fra denne Abstraction ved at sætte
det consecutive som sin sande Negation, og viser
% --- p. 121 ---
\setcounter{footnote}{0}%
\opage{121}saaledes, at det har den uendelige Negativitet i
sig; men det sætter tillige det consecutive Væsen
som et væsentligt Resultat, der maa forudsætte
sit forudsættende Princip, og gaaer derved uafbrudt
tilbage i sig selv, som det absolut fuldendte. Al
Negativitet og Uselvstændighed er hermed sat som
tjenende Moment for dets Positivitet og Selvstændighed;
hvad der udtræder i den exoteriske Fylde,
er allerede givet som negativ Væren i den esoteriske;
den Reflexionsbevægelse, hvorved det primitive
Væsen sætter Foranderligheden, og gaaer over
i den exoteriske Fylde, som i sin egen concrete
Transcendentalitet, er altsaa den selvsamme, hvori
det reflecterer sin Uforanderlighed, og vindicerer sin
evige Transcendents: \textbf{det consecutive} Væsens
\textbf{hypostaserede} Skin \textbf{ere} alle \textbf{Udstraalinger
af det primitive} \textbf{Lys}.

Det consecutive Væsen viser sig derved som
Negativitetens Selvudvikling. Det fremgaaer i
alle Henseender af sit Andet. Reflecteret i det
primitive, \textbf{er} det, før det bliver \textbf{sat}; det \textbf{er} forudsat,
som det, der skal sættes. Men reflecteret i
sig selv, er det sat, før det kommer til at være;
dets egentlige Væren fremgaaer af den Selvposition,
hvorved det negerer og ophæver sin Forudsætning.
Men heri ligger netop dets væsentlige
Modsigelse: det er, for at sættes af sit Andet;
og det sættes af sit Andet, for at sætte sig selv
og at være. Denne Modsigelse finder sin Opløsning
derved, at det consecutive Væsens eiendommelige
Reflexionsbevægelse lader dets hele Positivitet
opgaae som et i den satte og udviklede Negativitet
forsvindende Moment. Ifølge denne eiendommelige
Reflexionsbevægelse \textbf{er} det consecutive
% --- p. 122 ---
\setcounter{footnote}{0}%
\opage{122}Væsen et selvstændigt, for-sig-værende, med sig
identisk Væsen. Men dets Identitet er negativ.
Det sætter sig selv, som det, der sættes af Andet;
dets Selvstændighed udtrykker dets Uselvstændighed:
Positionen er Negation. Idet Identiteten
vil hypostasere sig, forsvinder den som Skin i en
Utallighed af relative Hypostaser; og idet hver
særskilt Væsenhed vil hypostasere og soutenere sin
For-sig-væren, sætter det sin egen Selvmodsigelse;
som Resultat af en svunden Selvstændighed, forsvinder
enhver Skinhypostase i en anden selvmodsigende Selvstændighed. Forskjelsmomenterne i den
exoteriske Fylde produceres altsaa, for at destrueres,
og destrueres, for at reproduceres: det væsentlige
Resultat er en sig selv opløsende Identitet.\footnote{Guds Væsen er ifølge dets Triplicitet uendelig affirmativt; Verdensvæsenet søger derimod uafbrudt Selvstændighed, uden at kunne naae den i sig selv. Den Selvstændighed, det har i dets eiendommelige Selvbevægelse, er kun Medium for Aabenbarelsen af Guds absolute Selvstændighed. Kunde det derimod i Virkeligheden tænkes løsrevet fra Identiteten med Gud, da vilde det derved blive diabolisk; thi det diaboliske Væsen er just den sig uophørlig hypostaserende Modsigelse, der dog aldrig kan naae en sand Hypostase.}

Det er det primitive Væsen, der uafbrudt sætter
og ophæver det consecutives Modsigelser; og
det er paa den anden Side Negativiteten i det
consecutive, der udtrykker Positiviteten i det primitive.
Idet de saaledes forudsætte hinanden gjensidig,
er Dignitetsforskjellen tillige sat i og med
den reduplicerede Forudsætning: deres vexelsidige
% --- p. 123 ---
\setcounter{footnote}{0}%
\opage{123}Relation er en gjensidig \textbf{Subordination}; thi
Subordinationen er det Uselvstændiges Reflexion i
det Selvstændige.

Det consecutive Væsen kan ikke vindicere sig
udelukkende Selvstændighed, thi derved vilde det
ophøre at være consecutivt, det kan heller ikke opgive
sin Selvstændighed og forsvinde, thi derved
vilde det primitive selv tilintetgjøres. Men den
Modsigelse, at det consecutive Væsens Selvstændighed
ligger i dets Uselvstændighed, hæves i Subordinationen.
At det har Selvstændighed, følger netop
deraf, at det i sig har det primitive som sit væsentlige
Princip; den Selvstændighed, det har i sig,
er altsaa ikke dets egen; det opoffres, idet det subordineres.
Men da det, som et i sig reflecteret
Skin, reflecterer, har og er i sig sin egen væsentlige
Forskjel, saa voldtages det ikke under sin Opoffrelse
af et fremmed Væsen: dets Underordnen er
dets egen fuldstændige Selvposition, og dets Selvstændighed resulterer uafbrudt af dets Uselvstændighed.

Det primitive Væsen kan da alene underordne det
consecutive under sig derved, at dette underordner
sig selv. Det Consecutives Selvposition er altsaa
den Basis, paa hvilken det Primitives Selvstændighed
hviler, den reale Forudsætning, af hvilken
dets hele Primitivitet fremgaaer. Men saaledes
bliver det primitive Væsen selv consecutivt, og det
consecutive primitivt. Subordinationens Forudsætten
forvandler sig da under sin Udvikling til en
blot Omsætten af det Selvstændige og Uselvstændige,
altsaa til en Coordination af begge. Det
væsentlige Fremskridt ligger deri, at alle de Reflexionsbevægelser
og thetiske Momenter, den cen%
% --- p. 124 ---
\setcounter{footnote}{0}%
\opage{124}trale Modsætning indeholder, nu ere satte og bestemte;
men da ethvert af Leddene i denne Modsætning
selv er paa een Gang baade det primitive
og consecutive Væsen, og det ene altsaa i Grunden
det selvsamme, som det andet, saa er den
hele Modsætning selv en central Modsigelse, og
maa følgelig gaae tilgrunde.\footnote{Naar den Centralmodsætning, som Logiken herved har antydet mellem Guds og Verdens Væsen, ikke kan holde sig, men Dialektiken atter maa lade dem forsvinde i hinanden, da er Grunden denne, at de Categorier, hvori Gud og Verden her opfattes, endnu ere for abstracte. Men den theistiske Type vindicerer ligefuldt sin Betydning; thi enhver Centralmodsætning, der forsvinder, danner derved en Overgang til fyldigere Categorier, i hvilke den vender tilbage med større Klarhed. Den typiske Udviklingsgang er da denne: Jo dybere Immanents, desto høiere Transcendents, jo concretere Universalitet, desto fastere Centralitet. Den speculative Logik kan derfor ikke standse, før den theistiske Centralmodsætning er sat med absolut Frihed, saa at det primitive Væsen ikke længere kan forsvinde i det consecutive.}

\underafsnit{c) Det absolute Væsen.}
\paragraf{23}{Urgrund, Indhold, Mediation.}

I det absolute Væsen ophæves enhver tilværende
Selvstændighed; Alt er i Grunden Eet, saaledes er
\emph{Urgrund}. Men Forskjellen er ligesaa
grundig som Identiteten. Alt er i Grunden forskjelligt; som Grunden til Alt oplader Urgrunden sig i et \emph{Indhold} af relative Grunde. Denne dobbelte
% --- p. 125 ---
\setcounter{footnote}{0}%
\opage{125}Grundbevægelse, hvormed Urgrunden deels ophæver Grundelementer, for selv at begrunde sin egen
Transcendents; deels sætter sig selv til, for i
Tilværelsens bestemte relative Grunde at udtrykke sin væsentlige Immanents, fuldbyrdes i \emph{Mediationen}. I „det rene Væsen" opgaae alle Værensbestemmelser
som Reflexer, det ene Reflexionsmoment gjenspeiler det
andet; i „Væsenets Dualitet„ bliver Reflexionen i % PRINTED AS IS: Dualitet„ (the closing quote set low; S145)
Andet alsidig bestemt som en Reflexion-i-sig. Væsenet
ophæver Skinnet, og reflecteres i sig som Væsen i Væsen,
ligesom ogsaa Skinnet omvendt optager Væsenheden
i sin Selvreflexion. Denne centrale Modsætning
imellem det primitive og consecutive Væsen reflecteres
igjen ved Subordinationen i een eneste Grundidentitet,
hvorved den immanente, alsidige Reflexion gjennemføres,
og Væsenet sætter sig som det \textbf{absolute Væsen}.
I det absolute Væsen fremtræder nu ethvert Væsenhedsmoment
med absolut Reflexion-i-sig. Det primitive
Væsen er absolut, thi det ophæver det consecutive som
et Moment i sin egen Selvreflexion, som Substrat for
sin Bestaaen; enhver af dets Hypostaser er følgelig
ogsaa absolut, thi den reflecterer hele det primitive Væsen
i sig. Paa den anden Side er det consecutive
Væsen, der i sin Selvbevægelse ophæver det primitive,
og, naar dets Reflexion er fuldendt, subordinerer det
under sig, ligeledes absolut; og da nu enhver af de
consecutive Hypostaser igjen reflecterer alle de øvrige,
saa er det Absolute selv sat i ethvert Moment af det
absolute Væsen. Da ingen Væren, intet Skin kan
falde udenfor det Absolute, kan dette igjen heller ikke
falde udenfor nogetsomhelst af dets Momenter; ethvert
Moment af det Absolute bliver følgelig selv sat
som et absolut Selvstændigt, ja som det Absolute selv.
% --- p. 126 ---
\setcounter{footnote}{0}%
\opage{126}Men det absolut Selvstændige taaler ingen Andethed
udenfor sig. Hvis det nemlig lod Andetheden tilbage,
da vilde denne atter, i Kraft af det Absolute, selv blive
et absolut Selvstændigt, og vi erholdt saaledes tvende
Absolute, af hvilke intet var absolut. Idet et Moment
af det Absolute bliver absolut selvstændigt, subordinerer
det derved alle de øvrige, og forvandler dem til Substrater.
Som nu alle Momenterne have lige Krav paa
absolut Selvstændighed, saa maae de ogsaa alle, uden
Undtagelse, forvandles til Substrater, og ere følgelig,
hvert for sig, absolut uselvstændige. Her er altsaa en
Utallighed af Selvstændige, der ophæves og forsvinde
i den absolute Selvstændighed.

Naar det Selvstændige ophæves, gaaer det tilgrunde.
Den selvstændige Væren er først den sande Tilværen;
at tabe sin Værens Selvstændighed, er at tabe sin egentlige
Væren; Overgangen fra absolut Selvstændighed til
absolut Uselvstændighed er derfor en Overgang saa
umiddelbar og brat, som fra Væren til Ikke-Væren:
at gaae \textbf{tilgrunde er at tilintetgjøres}. Men det
absolut Selvstændige ophæves kun, fordi det i sin Selvstændighed
selv har sin Uselvstændighed, fordi det Andet,
der udgjør dets Fundament og Substrat, selv vil
sættes som absolut selvstændigt; det gaaer altsaa kun
tilgrunde, forsaavidt det søger sit selvstændige Grundlag;
i denne Henseende er det saa langt fra at tilintetgjøres,
at det ved at \textbf{gaae under} meget mere selv bliver
det gedigne \textbf{Underlag} og Grundlag, hvori det
selvstændige Andet maa sætte sin sande Selvstændighed;
\textbf{der er altsaa ingen} \textbf{Tilintetgjørelse}. Men reduceres
nu den hele Mangfoldighed af Selvstændige til Substrat
for det Absolutes Selvstændighed, da bliver denne en
Selvstændighed, hvori Intet er selvstændigt, følgelig en
absolut Uselvstændighed; gaaer den absolute Selvstæn%
% --- p. 127 ---
\setcounter{footnote}{0}%
\opage{127}dighed omvendt op i de Selvstændiges Mangfoldighed
og Forskjel, da gaaer den netop derved under i lutter
Relativiteter, og vi have heller intet absolut Selvstændigt.
I Grunden selv ligger da Selvstændighedens og
Uselvstændighedens umiddelbare Identitet og Forskjel.
For at Selvstændigheden kan begrundes, maae de Selvstændige
gaae under, men for at begrunde de Selvstændige,
gaaer Selvstændigheden selv tilgrunde. Idet
Alt saaledes skal begrundes, bliver Intet begrundet.

Denne Grundmodsigelse i det absolute Væsen kan
først hæves ved den sande Grundbevægelse. Det Absolute
begrunder sig selv, for at begrunde sit Andet, og
det begrunder sit Andet, idet det begrunder sig selv:
det Absolute er kun absolut derved, at det selv er Mediet
imellem sig og Andet. Grundbevægelsen er da
som Følge heraf en dobbelt: deels en primitiv, hvorved
det Absolute begrunder sig selv som absolut; deels
en \textbf{consecutiv}, hvorved det lægger sig til Grund for det
Ikke-Absolute. Den ene af disse Grundbevægelser begrunder
igjen den anden, er altsaa den andens Medium: Begrundelsen er \textbf{Mediation}.

Som den absolute Selvstændighed, hvori alt Selvstændigt
har sin Grund, er det absolute Væsen den
primitive Grund ɔ: \textbf{Urgrund}. Da den rene absolute
Selvstændighed ophæver enhver selvstændig Væren, saa
er den i denne sin abstracte For-sig-væren forudsat som
Negationen af al Selvstændighed, som Det, hvori
endnu Intet er begrundet, men hvoraf al Væren, selv
den absolute, dog skal fremgaae og begrundes. I denne
sin Transcendents er Urgrunden altsaa bestemt som det
reale Intet, \textbf{den negative Identitet af} Alt. Det
Ikke-Absolute fremgaaer af sit Andet, og kommer just
derfor af det Absolute som af sit Intet; det Absolute
selv er sit eget Andet, det oprinder af sig selv, og er
% --- p. 128 ---
\setcounter{footnote}{0}%
\opage{128}er just  derfor det Uoprundne ɔ: det af Intet Oprundne.
Intet er altsaa Det, hvoraf Alt oprinder, men kun det
reale Intet begrunder en sand, real Væren, I „Vorden“
oprinder vel den rene Væren af det rene Intet,
men her er ingen Begrundelse; thi da Væren, som saadan,
kun er en abstract Umiddelbarhed, er ogsaa det
værende Intet en tom, abstract Negation. I Reflexionen
er Væsenet Skinnets og Skinnet Væsenets Intet,
men de vexle som flygtige Reflexer i det klare Lys, og
mangle Grundens Væren. I Mediationen derimod,
hvor enhver Væren er Reflex, og enhver Reflex en Væren,
hvor Vorden er Reflexion, og Reflexionen Vorden,
gaaer den reale Væren frem af det reale Intet som af
sin Grund.

Den \textbf{transcendente} Urgrund, der concentrerer Alværen
i negativ Identitet, er følgelig ingen umiddelbar
Værensindifferents, der selv umiddelbart og uden al
Grund slaaer over i Differents, thi dens satte Reflexion
er dens gedigne I-sig-Væren, hvorved den fastholder
Alt i absolut Eenhed; men dens I-sig-Væren er heller
ingen abstract Væsensreflexion, hvori den ene Hypostase
gjennemklarer den anden; i det reale Intet synker al
Væren \textbf{umiddelbart} sammen, og Reflexionens Lys
udslukkes i Urgrundens Nat. Idet Reflexionen gaaer
tilgrunde, ophæves den i Vorden; som det reale Intet
er derfor Urgrunden den umiddelbare Identitet af Alværens Princip og Begyndelse. Forsaavidt den nemlig
er den reale Værens negative Forudsætning, det absolute
Andet, hvoraf det Absolute selv skal fremgaae, er
den Princip; forsaavidt den derimod endnu ingen Væren
har begrundet, efterdi den selv er Intet, befinder den
sig selv i en Vorden, og maa derfor sættes som den
absolute Begyndelse. Med den absolute Begyndelse
begynder Alt; enhver Forskjel imellem Princip og Re%
% --- p. 129 ---
\setcounter{footnote}{0}%
\opage{129}sultat som imellem det Primitive og Consecutive, enhver
Adskillelse imellem den successive Udviklings forskellige
Led maa ogsaa have sin Grund i Urgrunden;
derfor vil Alt vorde med eet Slag. Men idet Alt paa
een Gang skal vorde og begrundes, og det Enes Begrundelse
forudsætter det Andets Undergang, efterdi det
forudsætter sit Grundlag, er den transcendente Urgrund,
som negativ Identitet, ɔ: som Eenhed af stridige Tendentser,
i en uendelig Spænding, Uro og Gjæring.\footnote{Denne Urgrundens dialektiske Selvmodsigelse har Jacob Böhme skildret med mystiske Farver i sit Skrift: \textit{De tribus principiis}: „Man findet im Urkund die allerschreklichste und strengeste Geburt, alles Herbe, Bittere und Feuer: da kann man nicht sagen, daß es Gott sey, und ist doch der innerlich erste Quell, der in Gott dem Vater ist, nach welchem er sich einen zornigen, eifrigen Gott nennet, und derselbe Quell ist das erste \textit{principium}, und ist Gott der Vater in seinem Urkund, daraus diese Welt sich urkundet. — In diesem \textit{principio} stehet nichts als nur die allerschrecklichste Gebährung, die größte Aengstlichkeit, feindliche Wonne, gleich einem Schwefelgeist, und ist eben der Höllen Porten und Abgrund, darinnen Fürst Lucifer in Verlöschung seines Lichts geblieben“.}

Denne Strid imellem Intet og Væren, der i Grunden
er en Strid imellem Grundlag og Begrundelse, opløser
Urgrundens abstracte Transcendents. Grunden
underlægger sig Alt, for at være Alt, dens Uselvstændighed
gaaer under i dens absolute Selvstændighed, den
giør sig selv til Grundlag for sin egen Begrundelse.
Da Grundens Vorden altsaa bestaaer deri, at den, ifølge
den primitive Grundbevægelse, ophæver det reale Intet
i den reale Væren, og sætter sig selv som sin egen
Følge, bliver den herved bestemt som den \textbf{positive}
% --- p. 130 ---
\setcounter{footnote}{0}%
\opage{130}\textbf{Identitet af Alt} ɔ: som den \textbf{transcendentale
Urgrund}.

Forsaavidt den transcendentale Urgrund forudsætter
sit reale Intet, idet den forudsætter en Grundbevægelse, hvoraf den resulterer, er der i Grunden Intet, hvis
umiddelbare Væren jo er ophævet, Intet, som jo er
gjort til Grundlag; men forsaavidt det derimod er den
transcendente Urgrund, der selv er bleven transcendental,
er der heller Intet, hvis reale Væren jo er sat i
og med Grunden, Intet, som jo er blevet begrundet.
Hvis Eenheden af den transcendente og transcendentale
Urgrund var umiddelbar, da var der ingen Grundbevægelse,
og følgelig hverken Grund eller Begrundelse.
Urgrunden selv forudsætter altsaa en Grundforskjel, og
hviler saaledes i tvende hinanden modsatte Grundlag.
Men hvis omvendt Urgrundens negative og positive
Identitet ikke var een og den selvsamme Grundidentitet,
da ophørte al Grundrelation, og der blev heller
Intet begrundet. Da den sande Urgrund først kan resultere
af den forudsatte Modsætning, er der ingen
Grundrelation, der kan forbinde de tvende primitive
Grundlag; de forholde sig derfor ligegyldige til hinanden,
og ere, hvert for sig, ubegrundede, efterdi de selv
udgjøre Grundbevægelsens Forudsætning. I Principets
Dualitet ligger altsaa den absolute Begyndelse.
Men da den primitive Grund selv resulterer af det
dualistiske Princip, indeholde de tvende Grundlag alene
Urgrundens Grund, og have følgelig Identiteten og den
gjensidige Grundrelation i sig; ethvert Grundlag er altsaa
selv Mediet imellem Grunden og Begrundelsen.
Ifølge denne Udvikling er Identiteten af den negative
og positive Identitet medieret, og den primitive Grund
bestemt som den \textbf{absolute} Urgrund.
% --- p. 131 ---
\setcounter{footnote}{0}%
\opage{131}Den negative Grundidentitet, det reale Intet, udgiør
da nu den hele absolute Urgrund, men kun under
den Forudsætning, at den positive Identitet, den reale
Væren, paa een Gang sættes baade som dens satte
Andet og tillige, ifølge Begrundelsen, som dens ophævede
Andet ɔ: som dens Grundlag; ligeledes vil nu
ogsaa den positive Grundidentitet, den reale Væren,
udgjøre den absolute Urgrund selv, der har overvundet
og ophævet det reale Intet. Den primitive Grund er
saaledes i sig selv hverken en abstract Væsenhed eller
en abstract Tilværen, hverken et realt Intet eller en
blot real Væren, men begges Medium. Den satte
Væren kommer altsaa, ifølge den gjennemførte Mediation,
ligesaalidt af Intet, som Intet af Væren; men
\textbf{Alt maa have sin} \textbf{tilstrækkelige Grund}.

Ifølge den primitive Grundbevægelse er enhver Forskjel
og Modsætning forsvunden som Grundelement i
Urgrundens Identitet; Alt er da i Grunden Eet; den
primitive, absolute Grund er den eneste Grund til Alt.\footnote{Andagten, saaledes som den beskrives af \textit{Dr.} W. H. Rothe (Læren om Treenighed og Forsoning): „Han (Mennesket) er borte, Sjælen svæver i et Lyshav, milde Følelser gjennemstrømme Hiertet, klare Livsskikkelser fare forbi Sjælen; i Syner, flygtige som Morgenskyer, i Hjertets Dyb bevæger hans Aand sig, men ikke som hans ꝛc.“ udtrykker den subjective Bevidstheds Nedsynken i Livets Urgrund.} Heraf udvikler sig nu en dobbelt Grundmodsætning.
For det Første er Urgrunden sig selv modsat; den er sit
eget Andet og som saadan selv baade Grund og Følge; dernæst er den sit Andet modsat, thi ifølge dens Selvbegrundelse
er den endnu blot bleven Grund; men en
Grund, der Intet begrunder, \textbf{er} ingen Grund. Den
% --- p. 132 ---
\setcounter{footnote}{0}%
\opage{132}primitive Grundmodsætning er saaledes vel en grundig
Forudsætning for den consecutive, men uden den aldeles
ugrundet. Kun det Ubegrundede kan begrundes;
thi hvad der er umiddelbart identisk med Grunden selv,
kan ikke følge af nogensomhelst Grundbevægelse, og
følgelig heller ikke begrundes; Grundens Andet maa
altsaa forefindes som et umiddelbart Værende, der slet
ikke ligger i Grunden, men er aldeles forskjelligt fra
den. Men kun Det, der fremgaaer af Grunden selv,
kan begrundes, og kun Det, der ligger i Grunden, kan
fremgaae af Grunden: det Ubegrundede maa dog ligge
i den absolute Grund. Den primitive Grund har altsaa
i sig Det, som er forskjelligt fra den, og som den
ikke selv er: den har et Indhold.

Indholdet er for det Første Grundens eget ifølge
den primitive Grundbevægelse, identisk med Urgrunden
selv; dets Momenter falde da idelig sammen med Grunden
og derved med hverandre indbyrdes. Urgrunden
ruger over sit \textbf{esoteriske} Indhold som over et evigt
Mysterium. Men paa den anden Side er Indholdet
Grundens Andet, der emanerer af Grunden, Det,
hvori den absolute Grund gaaer ud over sig selv og
sætter sin Relativitet; idet Indholdets Grundmomenter
saaledes adskille sig fra Grunden selv, skille de sig derved
fra hverandre indbyrdes, og fremtræde, hvert for
sig, med en relativ, begrundet Selvstændighed: Grundens
Mysterium oplader sig i det \textbf{exoteriske} Indhold.
Det esoteriske og det exoteriske Indhold ere i Grunden
Eet, og netop derfor grundforskjellige. Af den Grundsætning,
at Alt er Eet, udledes nu den, \textbf{at Alt} i
Grunden \textbf{er} forskjelligt; herfra begynder den consecutive
Grundbevægelse.\footnote{Det er Professor Martensens Fortjeneste, i sin „Mester Eckart“ at have fremdraget et Cardinalpunkt i Mystiken, hvorom Man saavel i Philosophien som i Theologien hidtil kun har havt en meget uklar Forestilling, nemlig den mystiske Bevidstheds Forhold til Mysterium og Aabenbaring. Logisk taget kunne vi nu sætte den speculative Mystiks Eensidighed deri, at den i sin dybsindige Skuen uafbrudt forsøger at reducere alle Tilværelsens Modsætninger til Urgrundens esoteriske Enhed, men ogsaa derfor ligesaa uafbrudt ved den dialektiske Modsigelse nødes til at vende tilbage til det exoteriske Indhold.}
% --- p. 133 ---
\setcounter{footnote}{0}%
\opage{133}Forsaavidt det exoteriske Indhold i dets Forskjel  fra
den absolute Grund dog er identisk med den, er det
herved bestemt som Urgrundens egen Alteration, ɔ: som
en Utallighed af \textbf{tilværende Grunde}; forsaavidt det
derimod sættes som Grundens Andet, repræsenterer det
sig som en \textbf{begrundet Tilværelse}. Den absolute
Grunds I-sig-væren er Forudsætning for dens Væren
i Andet; dens Transcendents begrunder dens Immanents;
den er saaledes selv Mediet imellem de tilværende
Grunde og den begrundede Tilværelse, og Mediationen
ligger deri, at det Absolute bliver relativt,
for at det Relative kan blive absolut. Da det er Urgrunden
selv, der er altereret, er det Absolute derved
selv blevet relativt; ethvert Tilværende har sit Maal
af det Absolute i sig selv, og udtrykker saaledes i sin
eiendommelige Begrændsning det hele Absolute.\footnote{Paa dette Standpunkt befinder Man sig, naar Man søger Guds Væsens Aabenbaring i et Græsstraa.} Men skal det Relative udtrykke det Absolute, da maa
det selv sættes med absolut Relativitet. Som et isoleret,
for-sig-værende Noget, har det Tilværende blot en
relativ Grund; kun under den Forudsætning, at det
selv er Moment i al den øvrige Tilværelse, i hvilken
den altererede Urgrund udtømmer hele sin 
Mangfoldig%
% --- p. 134 ---
\setcounter{footnote}{0}%
\opage{134}hed af relative Grunde, kan det siges at have den absolute
Grund i sig. Altsaa: det Tilværende har sin
Grund i sig, og er umiddelbart identisk med den, thi det
er, fordi det er; men ligesaa umiddelbart er det nu
ogsaa forskjelligt fra Grunden; thi det er kun til under
Forudsætning af en anden Tilværelse, og har følgelig sin Grund \textbf{udenfor} sig. Som Tilværelsens
Grund er den altererede Urgrund immanent i Tilværelsen
selv derved, at det ene Tilværende begrunder
det andet. Ethvert Tilværende er altsaa paa een
Gang baade Grund og Følge. Som Følge opstaaer
det ved et Andets Undergang, som Grund gaaer det
derimod selv under, for at et Andet kan opstaae. Men
ogsaa kun forsaavidt Tilværelsen udtrykker Urgrundens
egen Alteration, kan det ene Tilværende begrunde det
andet. Hvis ethvert Tilværende blot havde Grunden
umiddelbart i sig, da vilde Tilværelsens Mangfoldighed
udtrykke Mangfoldigheden af den altererede Urgrunds
tilværende Grunde, den hele Tilværelse vilde da være
begrundet med eet Slag, og Grundbevægelsen ophøre;
Grundens Forskjel fra sig selv blev da saa umiddelbar,
at Grundidentiteten og med den Grundrelationen aldeles
forsvandt. Kun derved, at det ene Tilværende
uafbrudt gaaer over i det andet, som Grunden i det
Begrundede, kan Grundbevægelsen vindiceres, og Identiteten
uafbrudt resultere af Forskjellen. Men hvis omvendt
det ene Tilværende blot havde sin Grund i et
andet Tilværende, da maatte det Ene forudsætte det
Andet i en uendelig Række; Forskjellen vilde da være
sat saa umiddelbart i Tilværelsen selv, at Grundbevægelsen
tabte sig i en uafbrudt Forandring uden Begrundelse.
Den Forudsætning, at det ene Tilværende
begrunder det andet, slaaer da ligesaa umiddelbart over
i den Forudsætning, at ethvert Tilværende umiddelbart
% --- p. 135 ---
\setcounter{footnote}{0}%
\opage{135}begrundes af Urgrunden\footnote{Traducianismen, der i Forældrenes Sjæle alene søger Grunden til Børnenes, er, til Ex., ligesaa eensidig som Creatianismen, der i Guds Skabervillie umiddelbart søger Grunden til enhver Sjæl.} Denne abstracte Modsigelse 
medieres i den consecutive Grundbevægelse derved, at
\textbf{Tilværelsens Grundrelationer specificeres}.

Fra den umiddelbare Værens Standpunkt viser Tilværelsens
Forskjel sig som en ydre overfladisk Forskjellighed;
her forefindes kun den ubestemte Vrimmel af
„Nogle“ og „Andre,“ der komme og svinde i ligegyldig
Vorden og Forandring, og under deres Adskillelse og
Forbindelse uophørlig søge et høieste Maal. Men
som Væren først finder sit Maal i Væsenet, saa er det
igjen Væsenets Bestemmelse, at gjøre sig til Værens
Grund og begrunde Tilværelsen. Betragtede fra det
absolute Væsens Standpunkt fremstille Tilværelsens
Forskjelle sig derfor som væsentlige og begrundede;
Grunden sætter det bestemte Maal, og Modaliteten lader sig
nu ikke mere opløse i reen Indifferents. Hvis
et hvilketsomhelst Noget kunde lægges til Grund for et
hvilketsomhelst Andet, vilde Eenheden selv være grundløs,
og enhver grundig Forskjel forsvinde. Det er da,
som Følge heraf, kun visse særskilte Tilværende, der
kunne sættes som Grundelementer for visse andre, og
blive Forudsætninger for disses Begrundelse; med det
væsentlige Maal vender saaledes ogsaa Valgforbindelsen
tilbage, og Tilværelsen grupperes, ifølge sit Grundindhold,
i afsluttede Kredse. Idet nu den ene og samme
Grundrelation, der gaaer gjennem Altilværelsen, maa
specificere sig til særskilte Relationer, for at blive begrundende,
medieres Urgrunden med enhver tilværende
% --- p. 136 ---
\setcounter{footnote}{0}%
\opage{136}Grund, og Modsigelsen ophæves i Grundens Selvbestemmelse. I det ene Tilværende ligger den bestemte Grund
til et andet Tilværende under den Forudsætning,
at ethvert „Noget“ ligesaa umiddelbart refererer
sig til en Kreds af „Andre“ som til sig selv.
Forsaavidt det Tilværende har Grunden i sig, refererer
det sig til sig selv alene, og kan ifølge
denne Grundrelation ikke umiddelbart indeholde
Grunden til Andet; men forsaavidt det har sin
Grund udenfor sig, har det derimod et overgribende
Indhold, et Indbegreb af Bestemmelser og
Relationer, hvorved det henviser til visse andre
Tilværende, som et Noget, der selv er sat \textbf{paa
Grund af Omstændighederne}. Det Indhold,
ethvert Tilværende saaledes har i sig, og det Indhold,
det har udenfor sig, ere da Indholdsmomenter
af een og samme Grundrelation, og i Mediationen
af disse Momenter ligger den specifiske
Grund. α) Grunden er den med sig medierede
Identitet af det Begrundende og Begrundede; det
Tilværende, der skal begrunde, og Det, der skal
begrundes, have, som Følge heraf, samme Indhold;
thi Andetheden selv ligger i det Tilværende,
der indeholder Grunden, efterdi dets Relation til
Andet kun er et Moment i dets Relation til sig
selv. β) Men den reale Grund ligger paa den
anden Side i den grundige Forskjel imellem det
Begrundende og Begrundede; det Tilværende, der
skal begrunde, maa derfor omvendt sætte sin Relation
til sig selv, som Moment for Relationen
til Andet; Grundens rene Mediation falder da
udenfor de Grundelementer, gjennem hvilke den
fuldbyrdes: enhver specifik Begrundelse \textbf{for}%
% --- p. 137 ---
\setcounter{footnote}{0}%
\opage{137}\textbf{udsætter sine særskilte Betingelser}. γ) Det ene Tilværende indeholder altsaa Grunden til det
andet derved, at der ligger et heelt Indbegreb af
Betingelser imellem Grunden og Det, der skal begrundes.
Betingelserne maae forudsættes som umiddelbart
givne, de udgjøre kun de discrete Punkter
i den Tilværelseskreds, indenfor hvis Peripheri
Begrundelsen søges; som \textbf{et} saadant Indbegreb af
Tilværende forholde de sig ligegyldige til den hele
Grundrelation, og give sig som \textbf{ubetingede}. Men
i deres ubetingede Væren ere de blot tilværende,
og kunne følgelig Intet betinge, efterdi deres Relation
til sig selv indbyrdes forsvinder i abstract
Umiddelbarhed: kun i Forholdet til Grundrelationen
slaaer deres umiddelbare For-sig-væren
over i en Væren for Andet; for en anden Grundrelation
vilde de selvsamme Tilværelsesmomenter
afgive andre Betingelser, de ere altsaa, som Grundens
Betingelser, selv \textbf{betingede}. Afhænge nu
saaledes alle Betingelser af den bestemte Grundrelation,
da maa denne derved fremtræde som ubetinget;
ved en blot Addition af Betingelser tilveiebringer
man ingen Begrundelse; her maa da begyndes
med et Tilværende, der kan siges at indeholde
Grunden. Men kun forsaavidt Betingelsernes
Indhold optages som Momenter i den tilværende
Grunds Indhold, kan der sættes en real
Forskjel imellem det Begrundende og Det, der skal begrundes; Grundrelationens Eiendommelighed er
altsaa dog betinget ved Betingelsernes umiddelbare
Tilværelse. Ved den immanente Eenhed af Grundens rene
Mediation og Betingelsernes Umiddelbarhed bliver
Grunden specificeret og bestemt. Men den saa%
% --- p. 138 ---
\setcounter{footnote}{0}%
\opage{138}ledes bestemte Grund er i sig utilstrækkelig. De
særskilte Betingelsers afsluttede Kreds er nemlig
kun et Led i den uendelige Række af alle øvrige
Tilværelseskredse, og den specifiske Grundrelation,
der gjennem sine Betingelser ubetinget fremhæves
som en for-sig-værende Bestemthed, kun et
Moment i den universelle Grundbevægelse. For
nu at frigjøres fra denne ubestemte Uendelighed,
maa den bestemte Grund da tillige sættes ved en
ubetinget Position af Urgrunden selv. Denne specifiske Position, hvorved Urgrunden fastsætter den
relative Grund ligeover for de særskilte Betingelser,
for at ophæve begge som Momenter i Mediationen
af den tilstrækkelig bestemte Grund, bærer
nu selv Umiddelbarhedens Præg. Ifølge denne
Umiddelbarhed er Urgrundens ubetingede Position
da ogsaa betinget; dens Specification har nemlig
den relative Grund og dens Betingelser til sin
ydre Forudsætning, hvorved dens Relation til sig
selv forsvinder i en Relation til Andet. Men de
relative Grunde og deres Betingelser, der udgjøre
den hele begrundede Tilværelse, ere kun et umiddelbart
Udtryk for Urgrundens egen Alteration;
den Forudsætning, hvorved dens specifiske Position
betinges, har den selv forudsat, den betinges altsaa
kun ved de Betingelser, den selv har givet sig,
og fremtræder saaledes i den tilstrækkelig bestemte
Grund som \textbf{det absolut Ubetingede}. I det absolut Ubetingede ere nu alle Grundbestemmelser
satte, og derved tillige ophævede. Tilværelsen
med dens relative Grunde er kun Forudsætning
og Betingelse for Urgrundens egen Selvbegrundelse;
derfor ophæver den primitive Grund de
Forudsætuinger, den selv har sat. I Tilværelsens % PRINTED AS IS: Forudsætuinger (for Forudsætninger; read so by two readers)
% --- p. 139 ---
\setcounter{footnote}{0}%
\opage{139}Umiddelbarhed er dens Indhold altereret, og denne
Alteration er et Ferment, hvorved den fuldbyrder
sin egen reale Mediation, og lader alle Betingelser
svinde som Momenter i det esoteriske Indhold.
Som Tilværen opstiger af Urgrunden, saaledes
maa den ogsaa synke tilbage i den; derfor gaaer
det Tilværende uophørlig tilgrunde; de relative
Grunde ere i sig selv kun Skingrunde. Men
omvendt ophæver Urgrunden nu ogsaa sig selv, idet
den ophæver sit exoteriske Indhold; dens Grundidentitet
er nemlig selv Betingelse og Forudsætning for Udviklingen af den grundige Forskjel. I de ydre Betingelser og de tilværende Grundes
Umiddelbarhed ligger Urgrundens Alteration og
Selvophævelse, men idet den saaledes igjen maa
ophæve sin egen Selvophævelse for at sætte sig
selv, viser sig Betingelsernes Negation som
en ufravigelig Betingelse for Urgrundens ubetingede
I-sig-væren, for dens esoteriske Indhold. Da
det exoteriske Indhold er sat af det esoteriske, saa
ligger det i Grunden dog stedse udenfor Grunden,
det har en fremmed Grund i sig, er et blot Begrundet,
et Substrat for den sande Grunds Mediation
med sig selv, og maa derfor ophæves.
Men idet den ydre Mangfoldighed ophæves som
betinget og sat, bliver den just derved selv et
Sættende, og forvandler sig til en real Grund for
det esoteriske Indhold, der nu bliver et Sat og
Begrundet. Som sat og begrundet maa det esoteriske
Indhold igjen ophæves, og i dets Selvophævelse
ligger Begrundelsen af det exoteriske. Tilværelsens
Skingrunde ere alle \textbf{begrundede Skin},
 de ere \textbf{Phænomener}. 
% --- p. 140 ---
\setcounter{footnote}{0}%
\opage{140}Uagtet nu denne Mediation, hvorved Urgrundens
Transcendents og For-sig-væren uafbrudt sætter sig
som Forudsætning for sin Immanents og Væren
for Andet, i Grunden aldrig kan ophøre, saa er
det dog betydningsløst fremdeles at ville fastholde
Adskillelsen af de modsatte Grundbestemmelser, efterdi
de alle have viist sig som eensidige Abstractioner.
Den ved Grundens Ophævelse gjenvundne Umiddelbarhed træder da frem med ubetinget Selvstændighed.
Da Mediationen saaledes er udtømt,
slaaer Identiteten atter umiddelbart over i Forskjellen, og denne i hiin; Grunden selv træder derved i Baggrunden og en ny Sphære aabner sig.

\overgang

Væsenets Identitet kan ikke være uden Forskjel. For
det rene Væsen er Forskjellen Skin, for Væsenets og
Grundens Primitivitet er den et consecutivt Væsen, en
Tilværelse af relative Grunde. Men, som et i Identiteten
forsvindende Moment, er Forskjellen endnu ikke kommen
til sin Ret; derfor maa nu Identiteten omvendt
lægges til Grund for en selvstændig Udvikling af phænomenale
Modsætninger, og Reflexionen selv forandre
sin Characteer. Medens nemlig den immanente, dialectisk-speculative
Reflexion prædominerer i Væsenheden,
bliver den udvortes distingverende og fixerende Reflexion
derimod fremherskende i
B. Phænomenets Sphære.

\afsnit{B. Phænomenets Sphære.}
\underafsnit{a) Grundphænomenet.}
\paragraf{24}{Materie, Form, Reflexivitet.}

Idet Phænomenet, den umiddelbare Eenhed af
Væsen og Skin, sættes som Resultat af det absolute
% --- p. 141 ---
\setcounter{footnote}{0}%
\opage{141}Væsens Mediation, falder det Umiddelbare udenfor
det Middelbare. Det i Phænomenet umiddelbart
Værende repræsenterer sig som \emph{Materie}, den udenfor Materien faldende Reflexion derimod som Form. Denne Forskjel bliver igjen ligesaa umiddelbart ophævet: Ingen Materie uden Form --- som umiddelbart
fixeret: Den samme Materie kan have de forskjelligste
Former, og den samme Form indeholde de forskjelligste
Materier. Dette ydre Reflexionsforhold udtrykkes i
Materiens og Formens \emph{Reflexivitet}.

Den dialektiske Overgang fra Væren til Væsen er
en tilsyneladende Fremgang, men en væsentlig Tilbagegang. Efter sin rene Umiddelbarhed repræsenterer Tilværelsen
sig, som om den kun var grundløs, fordi den
ikke behøvede nogen Grund, som om den selv var den
ubetingede Forudsætning for al Bevægelse og Udvikling,
altsaa som det Primitive. Det Væsenløse i Væren
selv er altsaa, fra denne Side betragtet, saa langt
fra at være en Mangel, at det meget mere udtrykker
den Aprioritet, hvorved Tilværelsen bliver istand til
at sætte Væsenet. Bevægelsen fra Væren til Væsen
er altsaa, ifølge denne Synsmaade, progressiv: Væsenet
\textbf{er en Følge af} \textbf{Væren}. Men idet Værens
Selvudvikling ender i Væsenet, indlyser det, at den
Uro og Bevægelse, der har drevet Væren ud over sig
selv, just hidrører fra dens egen Mangel og Eensidighed;
dens Overgang i Væsenet vidner om dens Oprindelse
af Væsenet. Da Væren saaledes viser sig som
den Umiddelbarhed, Væsenet selv har forudsat, for ved
Ophævelsen af denne Forudsætning at sætte sig selv,
saa er Bevægelsen fra Væren til Væsen \textbf{regressiv}.
% --- p. 142 ---
\setcounter{footnote}{0}%
\opage{142}Væsenet, der fremkommer som det Sidste, \textbf{er} da det
egentlige Første: \textbf{Væren er en Følge} \textbf{af Væsenet}.

Den Modsigelse, at Væsenet ikke kan fremkomme uden
som Følge af Væren, uagtet Væren selv først bliver
til som en Følge af Væsenet, finder sin Løsning i
Grundens Mediation.

Forsaavidt Grundens esoteriske Indhold forudsætter
et exoterisk som Det, der uophørlig sættes og skal sættes,
følger Grunden selv af den umiddelbare Væren;
men idet den centraliseres til Urgrund, ophæver den
Alværens Umiddelbarhed i sin egen absolute Selvstændighed,
følger derved af sin egen Følge, og nedsætter
Tilværelsen til en Følge af Grunden. Værens Fremgang
til Væsenet er da hermed sat som en Tilbagegang
til dens egentlige Udspring, og Grunden har ophævet
den Illusion, at Væsenet skulde være en Følge
af den umiddelbare Væren. Da der i enhver sand
Illusion imidlertid dog ligger et Skin af Sandhed, saa
kommer dette Sandhedsskin i Grunden nu ogsaa til
sin Ret. Det har nemlig viist sig, at det esoteriske
Grundindhold i sig selv kun vilde være et Skin, hvis
det esoteriske med dets relative Grunde og Betingelser
kun vare Skin; forsvinder den ene af disse Hovedsider,
da forsvinder hermed tillige den anden. Som det altsaa
er en Illusion, at ville fastholde en for-sig-værende
Urgrund, og lade den umiddelbare Tilværelse komme
og svinde som overfladisk Skin, eller at ville paa den
anden Side opløse den i sig centraliserede Urgrund i
de relative Grunde og deres Betingelser, saa er det
ogsaa en logisk Illusion, hvis Man\footnote{Som Hegel. Werke 4. B. Pg. 110—118.} lader begge disse
Hovedsider gaae under i en værende Umiddelbarhed.
% --- p. 143 ---
\setcounter{footnote}{0}%
\opage{143}Skal det absolut Ubetingede bestemmes som Det, \textbf{der} ophæver Urgrunden og dens Betingelser i en værende
Identitet af Væsen og Skin, saa bliver Grundbevægelsens
Fremgang til „Sagen selv“ atter en Tilbagegang
til Umiddelbarhedens absolute Primitivitet. Denne
uopløselige, adamantiske Umiddelbarhed nedsætter da den
forudgangne Mediation til en blot udvortes Reflexion,
en subjectiv Tankeb.vægelse. Alt ender saaledes i en % PRINTED AS IS: Tankeb.vægelse (a point for e, Tankebevægelse; two readers)
steril Ubetingethed, hvorfra Man hverken kan komme
frem eller tilbage. Det absolut Ubetingede kan ikke
fastholdes som \textbf{sat}, uden at det tillige forudsætter den
forudsættende Grund; det kan ikke tage sig sammen og
ophæve Grund og Betingelser, uden at det tillige maa
loslade Urgrunden, som den, der kan sætte andre Betingelser;
det kan ikke fremtræde som bestemt, uden at
reflectere den Grundbevægelse, hvorved det bestemmes.
„Sagen“ er altsaa den, at den uudtømmelige Urgrund
stedse afgiver bestemte Grunde, at Mediationens Strøm
bryder sine Bølger paa den Umiddelbarhed, i hvilken
den selv er chrystaliseret. Skal Grundens Indhold
nærmere bestemmes, da kan Bestemmelsen ikke længere
holdes i Mediationens rene Strømning, men maa sættes
og søges i det Medium, hvori Mediationen er gaaet
sammen med Umiddelbarheden, hvor Væsenets Identitet
er bleven identisk med Skinnets Forskjel, hvor den
bestemte Grund er gaaet op i det Begrundede, altsaa
i \textbf{en Væren}, \textbf{der saaledes følger af} \textbf{Væsenet},
\textbf{at Væsenet og Grunden} \textbf{selv følger med}.

Phænomenet er nu denne umiddelbare, og derved tillige
bestemte Identitet af Væsen og Skin. Idet Forskiellen
er ophævet, er Reflexionen forsvunden. Væsenet
er ikke længere Væsen, fordi dets Adskillelse fra
Skinnet er ophørt; af samme Grund er Skinnet heller
ikke længere Skin; begge Reflexer antage herved \textbf{et}
% --- p. 144 ---
\setcounter{footnote}{0}%
\opage{144}Præg af umiddelbar Væren. Men Identitetens Umiddelbarhed
er vunden gjennem Forskjellens og Modsætningens
Mediation; Skinnet trænger sig derfor frem
paany, giver sit Grundindhold et Anstrøg af Væsenhed,
og Reflexionen vender saaledes igjen tilbage, kun
med forandret Characteer. Væsenet bestemmes ved Skinnet,
og har sat sit Indhold i Skinnets Reflexion, og
Skinnet bestemmes umiddelbart som Det, der giver
Væsenet sin bestemte Skikkelse. Det ved Skinnet bestemte
Væsen holder sig nu ikke længere i en fjern
Baggrund, men er selv det Skin, hvori det skinner;
omvendt er Skinnet heller ikke længere den tomme Reflex,
der svinder og bedrager, men repræsenterer det grundige
Væsen: Phænomenets Reflexion reflecterer saaledes sin
egen Ophævelse. Saa umiddelbart, som det Bestemte
her falder sammen med Bestemmelsen, saa umiddelbart
adskille de sig igjen fra hinanden. Denne Vexel af
Identitet og Forskjel, ifølge hvilken Det, der skal bestemmes,
ligesaa umiddelbart har Bestemmelsen i sig
som udenfor sig, characteriserer Phænomenet som det i sig
Begrundede, som et Værende, der paa een Gang er i
Grunden og dog udenfor Grunden, som et Grundphænomen, hvis universelle Momenter ere \textbf{Materie} og
\textbf{Form}.

Materien er det værende Væsen, Væsenhedens
faste Forudsætning, der nærmer sig den palpable
Umiddelbarhed. I Modsætning hertil viser Formen sig
som Bestemmelighedens rene i sig reflecterede Activitet.
Materie og Form synes saaledes at forholde sig ligegyldige
mod hinanden.

Som væsentlig Umiddelbarhed er Materien ligegyldig
mod Reflexionen og Formen; derfor repræsenterer den
sig i Phænomenets Sphære som et Substrat, der antager
alle Værens Bestemmelser.
% --- p. 145 ---
\setcounter{footnote}{0}%
\opage{145}At Materien er det i Grundphænomenet Værende,
sees ogsaa deraf, at den paa een Gang gjengiver saavel Qvaliteten som Qvantiteten. Dens Grunddele,
Atomerne, udtrykke den oprindelige Qvalitetsforskjel, og
dens uendelige Delelighed gjengiver os den qvantitative
Discretion. Men Materien vorder ikke, uagtet den \textbf{er}
i en bestandig Vorden; den har Uforanderligheden tilfælleds
med Væsenet selv, uagtet den uafbrudt er Forandring
underkastet. Da den materielle Identitet ligesaa
lidt er istand til at sætte sin egen Forskjel, som
den materielle Forskjel formaaer at reflectere sin Identitet,
saa forbliver Materien, som saadan, \textbf{dorsk} og
dunkel. Materien er i sig selv ubestemmelig, og maa
dog bestemmes; dens Dunkelhed negerer Reflexionen,
og forudsætter dog Reflexionens Lys: Materiens egen
Ubestemmelighed udtrykker altsaa dens Krav paa at
bestemmes ved sit Andet ɔ: ved Formen.

Formen er, for sig tagen, den rene Bestemmen, der
endnu Intet har bestemt, efterdi det, den skal bestemme,
falder udenfor Bestemmelsen selv; den er den \textbf{klare},
\textbf{frit} svævende \textbf{Reflexion}, der i Materiens Væren
reflecterer sin egen Selvophævelse. Den formelle Identitet
gjenspeiler sin Forskjel; den viser derved ud over
sig selv, kommer i Bevægelse, og bestemmer sin egen
Ophævelse; af samme Grund reflecterer den formelle
Forskjel Identiteten som sin Negation og ophæver sig
selv. Formen bliver da herved sat som det active Princip,
som den uendelige Selvbevægelse.

I deres umiddelbare Modsætning ere Materie og
Form, hver for sig, lige endelige og uselvstændige.
Kunde Materien ikke bestemme sig selv, fordi dens
Identitet og Forskjel faldt ligegyldig ud fra hinanden,
saa kan den rene abstracte Form heller ikke fuldbyrde
sin Selvbestemmelse, efterdi dens Forskjel og Identitet
% --- p. 146 ---
\setcounter{footnote}{0}%
\opage{146}uafbrudt ophæve hinanden. Eenheden af den materielle
og formelle Forskjel er først den sande bestemte
Forskjel; thi forsaavidt Forskjellen har en materiel Side,
lader den sig ikke ophæve som et i Identiteten forsvindende
Skin, og forsaavidt den med det Samme
har en formel Side, viser den ud over sig selv, og bestemmer
sig til Identitet. Som Materien uden Form
er en Skygge, saa er Formen udenfor Materien et
Skin; kun ved deres gjensidige Affirmation erholde de
begge Selvstændighed. Materiens Negation af al Form
er dens egen væsentlige Form. Ved at holde sig som
det Ubestemte lige overfor den bestemmende Form giver
den sig en Bestemthed. Dens Ligegyldighed mod Formen
er dens egen Ligegyldighed mod sig selv; naar
den da optager Formen, og lader sig bestemme, optager
den ikke Mere, end hvad den allerede har; Formen udtrykker
kun Materiens egen Selvstændighed. Paa den
anden Side bliver Formen, naar den holder sig i sin
selvstændige Forskjel ligeoverfor Materien, reflecteret i
sin Andethed. Men nu er Formidentiteten just bestemt
som den, der i al Forskjel og Modsætning vinder sig
selv; Materien er da herved sat som det i Formens
Negativitet ophævede Andet ɔ: som Formens egen selvstændige
Bestaaen.

Denne udviklede Identitet af Materie og Form viser,
at den ene ikke begrunder den anden; thi da der
ingen Overgang er, saa kan der heller ikke være nogen
Grundrelation. Som der i Grunden selv ingen Materie
er, saa er der i Materien heller ingen Grund, og
som Formen ikke kan følge af Materien, efterdi den
selv er Materiens egen umiddelbare Forudsætning, saa
kan Materien heller ikke følge af Formen, da den abstracte
Form forudsætter Materiens Væren. Materiens
og Formens Bestaaen i og ved deres phænomenale Iden%
% --- p. 147 ---
\setcounter{footnote}{0}%
\opage{147}titet er ligesaa umiddelbart en Bestaaen i phænomenal Forskjel. Materien, der i sin Selvstændighed har ophævet
Formen, er derved allerede formet, ligesom Formen,
der har sat Materien som sit ophævede Andet, er
materialiseret. Begge ere da med eet Slag i Grundphænomenet
selv.

Bestemmes den phænomenale Identitet nærmere, da
viser den sig paa een Gang baade som væsentlig og
som uvæsentlig. Det er Materien væsentligt, at have
en bestemmende Form; men hvilken bestemt Form, den
skal have, er den endnu ganske uvæsentligt; ligeledes
er det Formen væsentligt, at materialiseres; men hvilken
Materie Formen skal bestemme, er aldeles ligegyldigt,
da Materiens Bestemmelse ene og alene afhænger
af Formen. Forsaavidt nu Materien reflecteres i
Formen som i sin Negation, bliver den derved begrændset;
men forsaavidt den selv bestaaer i Formen, ophæver
den ligesaa umiddelbart sin Ikke-Væren i Væren,
og sættes som den \textbf{bestemte} Materie, hvorved Formen
da ogsaa fixeres til umiddelbart \textbf{bestemt} Form. Denne
Bestemthed, hvormed al Mediation er ophørt, ligger i
\textbf{Overfladen}, der just viser sig som Materiens og Formens
rene Identitet og Forskjel. I den overfladiske
Bestemthed falder Materien sammen med sin Form; thi
den umiddelbare Form er Materiens Overflade; i Overfladen
ophører Materien, og forsvinder som Skin i Formens
Væren; som Materiens Bestemthed er Overfladen
Materiens Negativum, dens rene Andethed ɔ: Formen selv. Men paa den anden Side maa den overfladiske
Form nu ogsaa forsvinde som Reflex i Materiens
Væren. Paa Værens Standpunkt er al Bestemthed
kun en Grændse, der lader det Begrændsende falde
udenfor det Begrændsede; men Phænomenet, hvis Bestemthed
gjennem Væsenhedens Reflexion er hævet til
% --- p. 148 ---
\setcounter{footnote}{0}%
\opage{148}Værens Umiddelbarhed, identificerer derimod Grændsens
Sider, og sætter sig derved sin Overflade. Overfladen
er følgelig Materiens egen bestemte Væren; saa at det
ikke er Materiens Skin, men den værende Materie selv,
der i Overfladen viser sig paa det Allerbestemteste:\footnote{Synets og Følelsens Sands give os Overfladen fra tvende modsatte Sider. Synet seer egentlig kun Formen, ikke Materien; naar vi desuagtet synes at see Materien selv, da kommer det deraf, at Materiens Reflex er ophævet i Formens Væren. Følelsen refererer sig derimod umiddelbart til Materien, ikke til Formen; vi synes kun at føle Formen, fordi dens Reflex er ophævet i den bestemte Materie. Naar Synet sætter Formen som en Præmis, hvoraf man slutter til Materien, saa sætter Følelsen omvendt Materien som den Præmis, hvoraf man slutter til Formen. At Slutningen som saadan ikke kommer til Bevidsthed, hidrører derfra, at Reflexionen er ophævet i den phænomenale Identitet.}
den overfladiske Bestemthed af Materie og Form er
da, som begges satte Identitet, immanent og væsentlig.

Men da Mediationen er forsvunden, slaaer den overfladiske
Identitet ligesaa umiddelbart over i Forskjel,
hvorved den satte Bestemthed viser sig som uvæsentlig.
Materien, der er ligegyldig mod sin bestemte Form,
falder tillige udenfor Overfladen, og viser sig derved
som en uformet Masse, der kan modtage utallige andre
overfladiske Bestemmelser. Paa den anden Side
kan Formen, som det \textbf{bevægende}, \textbf{articulerende}
Princip, heller ikke acqviescere i Overfladens værende Bestemthed; thi da den formelle Identitet reflecterer sin
egen Forskjel, maa den, som Følge heraf, nedsætte sin
i Materien fixerede Overfladiskhed til Moment og Forudsætning
for andre bestemte Former; kun ved at ophæve de satte Bestemtheder som Reflexer i en Utallighed
af Formbestemmelser, kan Formen soutenere sin rene
% --- p. 149 ---
\setcounter{footnote}{0}%
\opage{149}Selvbestemmen. Da enhver bestemt Overflade saaledes
i det Uendelige kan ombyttes med en hvilkensomhelst
anden, er den overfladiske Forskjel imellem Materie og
Form poneret, og Overfladens Bestemthed reduceret til
en reen Uvæsentlighed.

Denne Forskjel imellem Materie og Form viser sig
endnu bestemtere, naar den opfattes i Relation til Overfladens
egen væsentlige og uvæsentlige Side. Naar
nemlig Materien uafbrudt tenderer til at hæve Overfladens
væsentlige Bestemthed, og som en „rudis indigestaque moles“ at synke i sit Chaos, saa ligger det derimod
i det formende og dannende Princip, at sætte den
overfladiske Bestemthed som et i den immanente Formation
opbevaret Moment, eller at omsætte Overfladens
Uvæsentlighed i dens Væsentlighed. Skjøndt Massen
i det Ubestemmelige unddrager sig den bestemte Overflade,
kan den desuagtet ikke undvige Formen, efterdi
denne selv griber ud over enhver fixeret Bestemthed.
Da Formen saaledes ved at ophæve sin overfladiske Forskjel
først bliver væsentlig og immanent, saa hjemfalder
Massen, idet den unddrager sig Overfladen, paa
ethvert Punkt til den væsentlige Form, og bliver derved
sat som dens \textbf{Stof}. Ved at give sig som Stof, er
Materien i enhver Henseende Formen underlagt. Forsaavidt
Overfladen endnu er ligegyldig og uvæsentlig,
ligner det ubestemte Stof den i sig uformelige Masse;
men da det er Formen selv, som ved at ophæve sin
egen Overfladiskhed har gjennemtrængt Massen, og behersker
Stoffet, saa er Overfladens Væsentlighed hermed
tillige sat; thi Stoffets Formation bestaaer netop
deri, at Formen og Overfladen bestemmes. Stoffets
væsentlige Form er overfladisk, efterdi det ligger i
Stoffets egen Elasticitet, at det ikke umiddelbart kan
have nogen fixeret Bestemthed. Formens Overfladisk%
% --- p. 150 ---
\setcounter{footnote}{0}%
\opage{150}hed er at sætte som Stoffets rene Formabilitet; men
allerede derved, at Materien er bestemt som Stof, har
den, som saadan, den bestemmende Form i sig; det specifiske
Udtryk for den bestemmende Form er umiddelbar
Formbestemthed, Overflade; følgelig er den overfladiske
Form ogsaa væsentlig.

Den alsidige Bestemmelse af Stof og Form udvikler
sig nærmere saaledes: α) Stof og Form ere umiddelbart
identiske; Identitetens Umiddelbarhed er den rene
Ubestemthed, og denne Ubestemthed er Overfladens eller
Bestemthedens Grund. Med Begrundelsen vender Mediationen
tilbage; idet Identiteten sætter sig som Grund,
ophæver den sin egen Umiddelbarhed, Formen kommer
i Bevægelse, og fixerer Overfladen som den immanente
Formations Resultat og Følge. Men naar Folgen er
sat, ophører Grundbevægelsen, og den immanente Formation
udpræger sig paa Overfladen\footnote{I denne Forudsætning er al Plastik begrundet.}. Stoffet selv bestemmes
herved fra en dobbelt Side: deels bliver det
nemlig i sig selv uendelig bestemt ved den immanente
Form, og derved hævet til \textbf{Formens Indhold}, deels
bliver det fra Endelighedens Side bestemt, det vil sige,
begrændset ved den overfladiske Form. Da nu Endelighedens
og Uendelighedens Side, ifølge den ophævede
Reflexion, ikke længere falde udenfor hinanden, da den
uvæsentlige og overfladiske Form umiddelbart udtrykker
den væsentlige og immanente, saa har Stoffet i Grunden
ingen anden Indholdsbestemmelse end den, der viser
sig paa Overfladen. Denne Eenhed af begge Formationens
Hovedsider er ikke et tautologisk „Einerlei“,
thi Mediationen er traadt imellem; men ligesaa lidt
er den at betragte som Reflex af en grundig Modsætning,
thi den forudsatte Mediation er ophævet; i begge
% --- p. 151 ---
\setcounter{footnote}{0}%
\opage{151}Siders umiddelbare Sammentræf ligger altsaa det bestemte
Stofs \textbf{specifiske Eiendommelighed}. β) Det
ene Stof træder med sin Eiendommelighed ligeoverfor
det andet; som Stoffets Eiendommelighed er given med
den bestemte Form, saa er omvendt Formens specifiske
Bestemthed fixeret i Stoffet. Da Forskjellen imellem
de heterogene Stoffer er sat overfladisk, saa viser Adskillelsen
sig endnu kun udvortes og ligegyldig; Formen
synes at paatrænge sig Stoffet, at komme til udenfra,
for at gjøre Skilsmisse imellem Stof og Stof. Desuagtet
er den overfladiske Form allerede forudsat som
den immanente Formations egen Bestemthed; Formadskillelsen begynder derfor ligesaa meget indenfra, saa at
Formen selv i ethvert eiendommeligt Stof forudsætter
et særeget Indhold. Hvorledes ethvert Stof i sig selv
er bestemt som Formens Indhold, kan det da kun vise
i Forhold til et andet eiendommeligt Stof, om hvilket
det Samme gjælder. Forsaavidt Adskillelsen er udvortes,
kan Forholdet kun bestemmes overfladisk; Stoffernes
Overflader berøre hinanden. Men i Berøringens Medium
falder det ene Stof umiddelbart sammen med det
andet; enhver Berøring frembringer en overfladisk Omformen,
idet den ophæver den ydre Forskjel; da nu Overfladen
selv er den immanente Forms egen Bestemthed,
saa frembringer enhver overfladisk Omformen tillige en
indre Transformation: de heterogene Stoffer gjennemtrænge hverandre indbyrdes\footnote{Disse tvende Synsmaader begyndte allerede at gjøre sig gjældende i den gamle ioniske Naturphilosophie. Den mechaniske Retning gik ud fra en Betragtning af den værende Mangfoldighed, og forklarede Tingenes Oprindelse af de i sig uforanderlige Urqvaliteters Adskillelse og Forbindelse ɔ: af Bestanddelenes alsidige Berørelse. Men ved denne Hypothese blev Qvaliteternes egen Oprindelse selv ligesaa uforklarlig som deres overfladiske Sondring og Forening. Den dynamiske Philosophie lod derimod Qvaliteternes egen Bestemthed forsvinde som Momenter i den uendelige Transformation uden at kunne give nogen tilstrækkelig Forklaring af Tingenes væsentlige og phænomenale Forskjel. Begge Theoriers Eensidighed ophæves i den dialectiske Identitet af alsidig Berøring og immanent Forvandling.}. Da al 
Gjennem%
% --- p. 152 ---
\setcounter{footnote}{0}%
\opage{152}trængen kun er en alsidig Berøring, er den materielle
Forskjel hermed vindiceret; den ene Partikel falder udenfor
den anden, saa at Stofferne ikke confunderes, og
flyde sammen i chaotisk Ubestemmelighed; men forsaavidt
den uendelige Berøring er en Gjennemtrængen, og
derved tillige en Transformation, er Formen selv sat
som det sondrende, articulerende Princip, og Materialismen
ophævet. Af denne Materiens og Formens
gjensidige Specification indlyser det nu, at, som Formen
formedelst Materien bliver forskjellig fra sig selv,
og derved adskilt i en Mangfoldighed af endelige Former,
saa bliver omvendt Materiens dunkle, indifferente Eenhed
adskilt og sondret formedelst Formen. Den i sin
rene Formation uendelige Formeenhed, der overalt vil
hæve enhver Modsætning til en flygtig Afglands af
Identitetens Lys, bindes og fixeres i den mørke, træge
Masse, forefinder i det Stof, den gjennemtrænger, uophørlige
Skranker, og falder saaledes ud fra sig selv i
en Række af materielle Former. I Materiens udvortes
ligegyldige Differents ligger altsaa Endelighedens
Grund. Men fra den anden Side betragtet, viser
Materien sig tillige som den Indifferents, der udsletter
al Forskjel. Som Materie og Masse er Grundstoffet
kun eet, og sig selv i alle Punkter liigt; Formen derimod,
som Væsenhedens rene Selvadskillelse, er den 
Re%
% --- p. 153 ---
\setcounter{footnote}{0}%
\opage{153}flexionsbevægelse, der driver det masseagtige Chaos over
i Forskjel, sondrer, articulerer og danner det. Det er
den ene, uendelige Form, der gjør Materien til de
mange særskilte Materier, til \textbf{et} discret Indhold af
heterogene Stoffer: Materiens endelige Bestemthed har sin
Grund i Formen. Ligesom Materien da fra sin
Side er Mediet imellem den bestemte Forms Identitet
med sig selv og Forskjel fra andre Former (et Medium,
der umiddelbart forsvinder i Formens Bestemthed), saaledes
er Formen igjen Mediet imellem Materiens specifiske
Eenhed med sig og Forskjel fra andre Materier;
herved er det i sig ubestemte Grundphænomen sondret
og spaltet i en \textbf{Mangfoldighed af bestemte} Phænomener.

Den bestemte Forms Reflexion i sig er dens Væren
i Materien, og den bestemte Materies I-sig-væren dens
Reflexion i Formen. Begges Reflexion i hinanden viser
sig da ligesaa umiddelbart som deres selvstændige Forsig-væren,
efterdi Reflexerne fra begge Sider hæves som
forsvindende Momenter; omvendt forsvinder derved tillige
deres For-sig-væren som Reflex af deres Væren i
hinanden: de ere, hver for sig, \textbf{reflexive}. Denne
Eenhed af Relation og Reflexion er igjen, naar den
nærmere bestemmes, at betragte som a) den \textbf{rene},
b) den \textbf{betingede} og c) den \textbf{ubetingede Reflexivitet}.
a) Ifølge den rene Reflexivitet er det den bestemte
Materie nok, at have en bestemt Form; som Formen altsaa bindes i Materien paa Grund af Bestemmelsens Umiddelbarhed, saa frigjøres den ligesaa
umiddelbart til For-sig-væren og Selvstændighed i
Kraft af Reflexionens Udvorteshed. Med Hensyn
til den overfladiske Formation, kan den ene og samme
% --- p. 154 ---
\setcounter{footnote}{0}%
\opage{154}Form meddeles til aldeles heterogene Materier\footnote{Den samme Jupitersform kan til Ex. udtrykkes i en Statue af Leer, Gips, Erts, Marmor o. s. v.},
og hvad den væsentlige Form angaaer, er den samme
Materie istand til at undergaae de allerforskjelligste
Metamorphoser, forsaavidt det nemlig er den overfladiske
Omformen, der, ifølge en uendelig alsidig
Berøring, tilveiebringer Transformationen, og forsaavidt
Stoffets Eiendommelighed ene og alene er
\textbf{at} udtrykke i Formens Bestemthed. Naar Stoffet
formedelst en Formforandring bliver et andet Stof,
da ligger Andetheden ikke i Materien, der fra sin
Side vedbliver at være sig selv liig, men i Formen.
Resultatet er, at i samme Grad, som den phænomenale
\textbf{Identitet} har vundet Selvstændighed, er ogsaa
\textbf{den phænomenale Adskillelse og} Forskjellighed
bleven \textbf{selvstændiggjort}\footnote{Jis, Vand, Damp kunne vise sig som det samme Stof i forskjellige Former.} Da Materiens og Formens phænomenale Adskillelse
imidlertid forudsætter Identitet, saa bestemme de selv
deres indbyrdes specifiske Bestemthed; denne Specification
kommer tilsyne i den betingede Reflexivitet.
Forsaavidt nemlig Formen selv baade er immanent
og overfladisk, saa adskilles de bestemte Phænomener
fra hinanden derved, at deres væsentlige Continuitet
reflecteres i en uvæsentlig Discretion; den phænomenale
Forskjel er kun en Skinforskjel, eftersom det
ene Phænomen ved Transformationen kan vorde det
andet. Men paa den anden Side kommer den væsentlige
Continuitet dog kun tilsyne i Discretionen; Phænomenernes Selvstændighed afhænger af deres
phænomenale Forskiel; i denne Selvstændighed er
% --- p. 155 ---
\setcounter{footnote}{0}%
\opage{155}Discretionens Skin væsentlig umiddelbart, og derved
tillige materielt; Continuiteten derimod, der ikke
kan sætte sig som bestemt Phænomen, forbliver uvæsentlig
og formel. Ifølge Discretionens Umiddelbarhed
er den phænomenale For-sig-væren sat som
Reflexion i andre Phænomener; enhver Formforskjel
er overfladisk, ligegyldig og foranderlig. De discrete
Middelpunkter blive staaende som Substrater,
hvis vexlende Former meddeles udenfra. Men naar
ethvert Phænomen henter sin Form udenfra, og der
udenfor det givne Phænomen ikke kan gives Andet
end Phænomener, saa er Phænomen uafbrudt bestemt
ved Phænomen, deres Exclusion er en Assimilation,
saa at den Udvorteshed, ifølge hvilken Phænomenet
maatte hente sin Bestemmelse fra andre Phænomener,
altsaa er en Andethed, der ligger i det bestemte
Phænomens egen Bestemmelse. Herved er den discrete
Form sat immanent, saa at ethvert Phænomen
nu kan siges at bestemme sin egen Form. Men
saaledes er ogsaa den phænomenale For-sig-væren
igjen forvandlet til en Væren for Andet, thi Phænomenet
bestemmer kun sig selv ved at bestemme
andre Phænomener. Resultatet bliver altsaa dette:
ethvert Phænomen bestemmes af andre Phænomener,
saa at det derved sættes som et mod sin overfladiske
Form ligegyldigt Stof, men ligesaa umiddelbart er
Phænomenet bestemt i og ved sig selv, det er et i
sig specificeret Indhold, der sætter Formationen sit
bestemte Maal: Formens og Stoffets gjensidige Uafhængighed
og Reflexion i hinanden, eller deres Reflexivitet,
er alsidig betinget. Phænomenet er kun bestemt, forsaavidt det kommer
tilsyne (thi hvad der ikke kommer tilsyne, kan heller
ikke være Phænomen); men Phænomenet kommer
% --- p. 156 ---
\setcounter{footnote}{0}%
\opage{156}kun tilsyne i den overfladiske Form, i andre Phænomener,
dets Grund er opløst i lutter Betingelser; Phænomenets Selvstændighed er saaledes negativ;
den viser sig umiddelbart som Uselvstændighed. Men
denne Uselvstændighed er igjen negativ, ethvert Phænomen
ophæver Betingelserne i sin egen Grund;
Grunden selv forsvinder i den phænomenale Selvstændighed,
og sætter sit Phænomen ubetinget.

Forskjellen er nu dreven til sit umiddelbare Extrem.
Formadskillelsen er materiel, thi den er udvortes
og umiddelbar; men Materien er tillige alsidig
formet, efterdi al Forskjel er bleven overfladisk.
Formernes Overfladiskhed er ikke Andet end
\textbf{deres} materielle Væren. I enhver bestemt Form er
Materiens Bestemthed nu tillige sat saa umiddelbart,
at enhver Forandring ophæver Phænomenets
Selvstændighed. Reflexiviteten er følgelig ubetinget.
Phænomenet kan altereres, deri viser sig dets
betingede Selvstændighed; men det altereres kun da,
naar de givne Betingelser ere ophævede, og nye indtraadte;
de forandrede Betingelser give sig hver
Gang som Momenter i det Ubetingedes absolute Bestemthed.
Som Betingelserne ere, saa er ogsaa den
Ubetingethed, hvori de ophæves, og saaledes er ogsaa
Phænomenet bestemt. Da der nu ikke kan gives
nogen phænomenal Forandring uden forandrede Betingelser,
og de forandrede Betingelser stedse concentreres
i det Ubetingede, saa er Overgangen fra Phænomen
til Phænomen en Overgang fra det Ubetingede
til det Ubetingede. Denne Overgang er formel,
tautologisk og usand, Forskjellen forsvinder som
et Skin, da det Ubetingede stedse er identisk med
sig selv. Men ligesaa umiddelbart repræsenterer den
sig som udvortes og materiel; Overgangen skeer ved
% --- p. 157 ---
\setcounter{footnote}{0}%
\opage{157}et Spring, efterdi det ene Ubetingede \textbf{er} ligegyldigt
mod det andet; denne materielle Adskillelse forsvinder
igjen som Skin i det Ubetingedes væsentlige Identitet, og saa fremdeles i det Uendelige. Da
Materien og Formen, som Følge heraf, ikke længere
kunne adskilles, maae de nødvendigviis tabe deres
særskilte Betydning, og opgaae i en Verden af selvstændige, frit udtrædende \textbf{Phænomener}.

\underafsnit{b) Existents.}
\paragraf{25}{Ting, Egenskaber, Tilsyneladelse.}

Idet det bestemte Phænomen udtræder af Grundens Inderlighed, og kommer til \emph{Existents}, bliver
dets Reflexion-i-sig en ligesaa umiddelbar For-sig-væren,
som dets Reflexion-i-Andet en Væren-for-Andet; det kommer derved til at bestaae som \emph{Tingen}, der har sine \emph{Egenskaber}. Egenskaberne
vise sig at være Tingens egen Existents: Tingen
lader sig tilsyne i Egenskaberne; omvendt forsvinde
Egenskaberne i Tingens Bestaaen: Tingens Væren i dens Egenskaber er kun \emph{tilsyneladende}. Denne existentielle Antinomie imellem Immanents og
Transcendents ophæves i \emph{Tilsyneladelsen}. Modsætningen imellem det primitive og consecutive
Væsen har nu udviklet sig til en Modsætning imellem

Grundens i sig værende Eenhed og de selvstændige
Phænomeners Mangfoldighed. Denne Modsætning er
absolut; thi idet den ophæves, bliver den sat, og hver
Gang den forsvinder, kommer den igjen tilsyne. I Me%
% --- p. 158 ---
\setcounter{footnote}{0}%
\opage{158}diationen har Grunden sin Selvstændighed, men den
medierende Grundbevægelse maa stedse afgive \textbf{et} værende Resultat, hvori Umiddelbarheden vender tilbage.
Overgangen fra den tilstrækkelige Grund til det selvstændige
Phænomen er en Omsætten af Mediation og
phænomenal Umiddelbarhed, der fortsætter sig i det
Uendelige. Den phænomenale Selvstændighed er grundet
i den ophævede Mediation; men Mediationens Ophævelse
ligger i Mediationen selv; denne bliver saaledes
det Overgribende, ophæver sin egen Ophævelse, og
nedsætter Phænomenet til et i Grundens Selvstændighed
forsvindende Moment. Identiteten af Grundens
og Phænomenets Selvstændighed er at sætte som
\textbf{Existents}.

Den existentielle Eenhed af Grund og Phænomen
er for det Første umiddelbar, og som saadan phænomenal.
Naar Grunden har adskilt sig som den absolute
i -sig-værende Urgrund fra den tilværende og relative
Grund, naar den under denne Selvadskillelse har
forudsat og ophævet sit bestemte Antal af ydre Betingelser,
da slutter den alle sine Extremer sammen, og
fixerer sig som den tilstrækkelig bestemte Grund. Den
tilstrækkelige Grund er i Indbegrebet af sine ophævede
Betingelser ubetinget, dens Grundelementer sondres til
særlig Bestemthed, for i gjensidig Bestaaen og negativ
For-sig-væren at reflectere hverandre; den tilstrækkelige
Grund kommer derved selv tilsyne i slaaende Umiddelbarhed:
\textbf{den existerer som det} \textbf{selvstændige} Phænomen.

Paa den anden Side er den Umiddelbarhed, hvormed
Grundphænomenets Materie og Form falde udenfor
hinanden, kun en Reflex af Grundens overgribende
Mediation. Skjøndt Materien synes at være ligegyldig
mod alle Formens Bestemmelser, saa er denne Li%
% --- p. 159 ---
\setcounter{footnote}{0}%
\opage{159}gegyldighed dog kun en væsentlig Uvæsentlighed. Som
Materien ikke er af sig selv, saa kan den heller ikke
være i sig selv; den maa blive Grundlag for \textbf{den} bestemmende
Form, og viser sig derved som et Mediationsmoment,
der tilhører Grunden. Den absolute
Urgrund behøver selv Phænomenernes Mangfoldighed,
for deri at grunde sin egen Eenhed, den forudsætter
det hele overfladiske Spil af heterogene Stoffer og materielle
Former, for deri at klare sin egen immanente,
affirmative Selvstændighed: Phænomenet existerer
\textbf{altsaa som et Gjenskin} \textbf{af den existerende
Grund}.

Grunden og Phænomenet have da, som Følge heraf,
samme Indhold, Forskjellen er kun formel: \textbf{alle}
Momenter i det Absolute ere at bestemme som Grundmomenter,
thi der gives ingen phænomenal Bestemthed,
uden at den jo er sat af Grunden, ligesaa lidt
som der gives nogen Grundbestemmelse, der ikke tillige
maa vise sig som phænomenal. Betragtes den absolute
Existents fra Mediationens Hovedside, saa existerer
det Absolute selv som Urgrundens esoteriske Indhold,
i hvilket de phænomenale Modsætninger forsvinde
som immanente Reflexer; bestemmes den absolute Existents derimod fra den anden Hovedside som ophævet
Mediation, som reflecteret Væren, da existerer det Absolute
selv i Grundens exoteriske Indhold, i de selvstændige
Phænomeners ydre Mangfoldighed, saa at
den primitive Væsensfylde er absorberet i det consecutive
Væsen.

Men i denne existentielle Identitet ligger en existentiel
Forskjel, ifølge hvilken enhver af Hovedsiderne selv
bliver at sætte som den absolute Existents.
% --- p. 160 ---
\setcounter{footnote}{0}%
\opage{160}1) \textbf{Det} Absolute existerer som Phænomen.

I det selvstændige Phænomen udtømmer den tilstrækkelige
Grund sit bestemte Indhold; dens dunkle prægnante
Inderlighed klarer sig, Alt kommer tilsyne i
overfladisk Umiddelbarhed. Da alle Phænomenets Betingelser ere ophævede i dets ubetingede Selvstændighed,
saa refererer det existerende Phænomen sig til sit eget
relative Centrum, som til et i sig absolut Middelpunkt,
og beroer paa sig selv alene. Med samme ubetingede
Selvstændighed træder det ene Phænomen op efter det
andet, og da ethvert af disse har sin tilstrækkelige
Grund i sig selv, saa forsvinder Urgrundens absolute
Centralitet i de utallige Phænomeners relative Centra.
Phænomenernes For-sig-væren udgjør deres umiddelbare
ubetingede Existents, ethvert Phænomen kommer tilsyne
i sit eget Lys; dets Relation til andre Phænomener,
der skulde vise dets uendelige Betingethed, er et uvæsentligt
Skin, eftersom Betingelserne overalt ere ophævede
i ubetinget Selvstændighed. Phænomenernes gjensidige
Forhold, deres Væren for hinanden, er da, fra
denne Side betragtet, kun en udvortes Reflexion.
Men den phænomenale For-sig-væren er ligesaa umiddelbart
en Væren for Andet, Reflexionen i sig er selv
overfladisk; den er væsentlig en Reflexion i Andet,
efterdi den er phænomenal. Ifølge Phænomenets umiddelbare
Reflexion-i-sig blive Betingelserne ophævede i
dets ubetingede For-sig-væren, ifølge den umiddelbare
Reflexion-i-Andet opløses den ubetingede For-sig-væren
i lutter Betingelser. Kun i denne bratte umedierede
Omvexling kommer den overfladiske Existents tilsyne.
Formen for det Absolutes Existents i Phænomenernes
Række er altsaa den uendelige Progres. Det Absolute
kan som saadant ikke existere i noget enkelt Phænomen,
thi Phænomenets ubetingede Selvstændighed slaaer umid%
% --- p. 161 ---
\setcounter{footnote}{0}%
\opage{161}delbart over i Uselvstændighed, ligesaa lidt kan en absolut
Existents udtrykkes i Phænomenernes Sum, thi
ethvert existerende Phænomen er behæftet med samme
væsentlige Modsigelse. Den phænomenale Existents
er saaledes uendelig fragmentarisk; den maa uophørlig
begynde forfra, og den Progres, der skulde indeholde
den nye Begyndelse, er kun en Skinprogres, i hvilken
den phænomenale Uselvstændighed overalt kommer tilsyne.
Som det Tilværende maatte gaae tilgrunde,
fordi det i sin værende Umiddelbarhed manglede al
Grund, saa maa det existerende Phænomen nu ogsaa
gaae tilgrunde, fordi det vil ophæve og i overfladisk
Bestemthed udtømme den absolute Grund. Jo klarere
den particulære Selvstændighed fremtræder, desto skarpere
viser sig den Overgang, hvori det Selvstændige
finder sin Opløsning. Idet nu det Absolute saaledes
synes at forløbe sig i idel Relativitet, og sprede sit
Indhold i peripheriske Centra, forraader den phænomenale
Existents med det Samme sin uendelige Negativitet;
Urgrunden optræder paany som det nivellerende
Princip, og tager sin absolute Centralitet tilbage.

2) Det Absolute existerer som Grund.

Kun i Forhold til Phænomenernes umiddelbare Selvstændighed
kommer Urgrunden tilsyne som den altopløsende
Grund; men Alt, hvad den herved opløser, har
den selv poneret; den tager altsaa kun sit Eget tilbage,
for at affirmere sig selv. Det har allerede viist
sig, at Grundens esoteriske og exoteriske Indhold reflectere
hinanden gjensidig; det esoteriske Indhold er i sig
Grundidentitet, og har saaledes med eet Slag alle Forskjelsmomenter
i sin indre Uendelighed, alle Tilværelsesprinciper
i sit ene Grundprincip; men i denne sin
umiddelbare Primitivitet har Urgrunden endnu ingen
% --- p. 162 ---
\setcounter{footnote}{0}%
\opage{162}Existents; den er endnu kun forudsat som Alværens
reale Begyndelse, som den væsentlige Vorden, som Oprindeligheden
selv. I Modsætning hertil frembyder
det exoteriske Indhold en Evolution af den grundige
Forskjel. Her er Eenheden sat i Mangfoldighedens
Form, her er en sporadisk Fleerhed af Principer og relative
Begyndelser, en Utallighed af tilværende Grunde
og endelige Betingelser. Det har fremdeles viist sig,
at ethvert af Leddene i den centrale Modsætning selv
maa forudsættes som absolut, for at det i sig fordoblede
Grundindhold alsidig kan udvikles. Den absolute
Urgrund er i sin oprindelige Væsenhed uendelig Mediation
af abstract Identitet, negativ, transcendent Eenhed,
og abstract Forskjellighed, negativ, transcendental
Mangfoldighed. Denne Mediation grunder sig i dens
egen primitive Bevægelse. --- Kun idet Urgrunden tager
sin concrete Eenhed tilbage gjennem den exponerede
Mangfoldighed, bliver dens esoteriske Indhold udviklet;
men for at den alsidige Sondring kan realiseres, maa
den exoteriske Proces af endelige Grunde og Betingelser
forudsættes som det Medium, igjennem hvilket Tilbagevendelsen
uophørlig skeer. Derfor maa den consecutive
Grundbevægelse ogsaa fuldende sit Kredsløb,
og dens exoteriske Indhold uafbrudt restaureres gjennem
Urgrundens Medium. Var den exoteriske Progres af relative Grunde og særskilte Betingelser kun
et flygtigt Skin, indeholdt den ikke et i sig relativt
Absolutum, da vilde Urgrunden mangle det reale Medium
for sin immanente Explication, da vilde den savne
den discrete Værens Andethed, hvori Væsenhedens og
Oprindelighedens rene Lys kunde bryde sine Straaler;
men omvendt, hvis den primitive Grundbevægelse gik
under i den exoteriske og consecutive, hvis de brudte
Straaler ikke igjen kastedes tilbage i Urgrundens Dyb,
% --- p. 163 ---
\setcounter{footnote}{0}%
\opage{163}da vilde Grundmysteriet forblive i evig Dunkelhed, da
maatte den oprindelige Eenhed afbrydes i Grundens
Alteration, og det Absolute fornegte sig selv.

I den phænomenale Existents har nu Grundens exoteriske
Mediation culmineret. Mangfoldighedens rige
Afvexlinger ere komne tilsyne; den positive Umiddelbarhed
har fixeret sig i de existerende Phænomeners relative
Centra; det exoteriske Grundindhold er sat med
absolut Overfladiskhed, og derved forudsat som det i sig
Absolute. Men denne umiddelbare Positivitet er tillige
uendelig Negativitet. α) For at existere maa det ene
selvstændige Phænomen forudsætte det andets Existents;
gives der eet eneste existerende Phænomen, saa gives
der med det Samme en uendelig Mangfoldighed. Det
selvstændige Phænomen reflecterer altsaa i sin umiddelbare
Existents en Utallighed af selvstændige Phænomener.
At ethvert Phænomen bestaaer i sig, vil da sige,
at det bestaaer i alle de andres Bestaaen. β) Men
Phænomenernes Reflexion i hverandre er overfladisk, de
relative Centra frastøde hinanden gjensidig; det ene
Phænomen er umiddelbart det andets Negation, forsaavidt
det ene existerer, kan det andet ikke existere; bestaaer
det ene Phænomen altsaa i de andres Bestaaen, saa bestaaer
det eo ipso i sin egen Ikke-Bestaaen. Den Negation,
hvoraf den existentielle Umiddelbarhed fremgaaer,
er selv en existerende Umiddelbarhed, hvoraf Negationen
resulterer. γ) Denne Modsigelse finder ikke sin Løsning
i den consecutive Proces, der tvertimod maa betragtes
som dens uendelige Fortsættelse. Det Væsentlige
er, at den existentielle Forskjellighed kommer tilsyne;
men idet den bliver phænomenal, bliver den tillige
umiddelbar og derved uvæsentlig. Det er saaledes kun
dets uvæsentlige Side, der kan forsvinde i den opløsende
Grund, og det er denne Opløsning, hvorved Ur%
% --- p. 164 ---
\setcounter{footnote}{0}%
\opage{164}grunden ophæver sin Alteration, og vender tilbage i
sig selv. Denne Tilbagevenden er først den sande,
væsentlige Fremskriden. Det pleromatiske Centrum er
under sin indre Selvadskillelse slaaet ind i de relative
Centra, fordi Grunden først kommer til sand Bestaaen
i Existentsen. Men i den phænomenale Existents er
Urgrunden fortabt, og dens Centralitet ophævet. Den
pleromatiske Centralitet bestaaer altsaa, som det phænomenale
Centrum bestaaer ɔ: den bestaaer i dens Ikke-Bestaaen,
den existerer i Negationen af dens Existents.
Først idet den fuldbyrder denne Negation, affirmerer
den alle existentielle Momenter, og sætter sig derved
som det i sig selv Existerende, som den overgribende
Centralitet, som de bestaaende Phænomeners primitive
Forudsætning, som det existentielle \textbf{Dyb}.
3) Det Absolute existerer i Mediationen af Centralexistentsen og den phænomenale Progres.

Den i Urgrundens Eenhed tilhyllede Mangfoldighed
er en negativ Identitet, der ophæver sig selv, for at
sætte sin Forskjel. At Alidentiteten har i sig uendelig
Forskjel, og at omvendt Mangfoldigheden og den uen,
delige Forskjel reflecteres i sig ved at reflectere Identiteten,
viser den rene Mediation. Desuagtet er denne
Mediation kun Formen for det Absolutes uendelige
Selvbestemmen, og som saadan uden al Bestemthed;
den er kun formel. Men den formelle Bestemmen sætter
og forudsætter Bestemthed, Andetheden dukker op af
Reflexionens Eenhed, og Formen afsætter sit Indhold,
Mangfoldigheden skal altsaa bestemmes og sættes som
Mangfoldighed, og Forskjellene udpræges i deres Forskellighed;
omvendt skal ogsaa Grundeenheden sættes
og articuleres i sin egen concrete Selvreflexion, saa at
alle Forskjelsmomenter vende tilbage og forsvinde i
% --- p. 165 ---
\setcounter{footnote}{0}%
\opage{165}den som reale Skin, som immanente Reflexer: Overfladen
skal reflecteres i Dybet, og Dybet i Overfladen.
Idet nu den ydre, i sig reflecterede Mangfoldighed
bliver \textbf{sat}, saa bliver med det Samme den umiddelbare
Aleenhed \textbf{ophævet}, Grunden gaaer tilgrunde i de existerende Phænomener; Grundidentitetens Reflexion i sig
forsvinder som Reflexion i Mangfoldigheden. Her er
altsaa en væsentlig Overgang, en Universalbegyndelse,
hvorved den phænomenale Mangfoldighed ophæver
den ubestemte Grundeenhed, og begynder den exoteriske
Progres. Men paa den anden Side er Mangfoldighedens
Reflexion i sig gjennem den ophævede
Grundidentitet kun Forudsætning for Aleenhedens Selvreflexion
i den ophævede Mangfoldighed. Naar Ur-Eenheden
\textbf{sættes} efter sin concrete Identitet med sig,
saa er de existerende Phænomeners ydre Uendelighed
med eet Slag \textbf{forsvunden} af Rækken, og de bestemte
Forskjelsreflexer maae søges i Grundexistentsens klare
Dyb. Den immanente Centralexistents resulterer altsaa
af Phænomenernes Tilbagegang i Grunden, den fremgaaer
af phænomenale Betingelser, og har saaledes selv
en \textbf{relativ Begyndelse}. I \textbf{Mediationen af den
universelle} \textbf{og relative Begyndelse er} \textbf{først den
grundige og} \textbf{sande Realbegyndelse at sætte}.
a) Ethvert existerende Phænomen er behæftet med den
Modsigelse, at dets Existents paa een Gang er uendelig betinget ved det Absolutes Universalmediation,
og dog i sin umiddelbare Bestemthed indskrænket til
en særegen Kreds af endelige Betingelser. At ethvert
Phænomen existerer i ubetinget For-sig-væren, følger
altsaa deraf, at den ene phænomenale Betingelseskreds
er begrændset ved den anden, de specifiske
Grundrelationer afbryde Universalgrundens Mediation
i det Uendelige, og \textbf{de} relative, partielle Begyndelser
% --- p. 166 ---
\setcounter{footnote}{0}%
\opage{166}nedsætte Universalbegyndelsen til en udvortes Reflexion.
Men netop derved bliver den partielle Begyndelse
en Selvmodsigelse; i den phænomenale Række:
\dots{} B, C\dots{}. kan C ikke begynde uden under Forudsætning
af B; den Kreds af Betingelser, hvori C
existerer, er selv betinget ved den Betingelseskreds,
i hvilken B har sin Existents; B forudsætter igjen
A, og saaledes i det Uendelige. Kunde blot eet
eneste Phænomen begynde sin Existents, da vilde
med det Samme alle øvrige begynde: enhver umiddelbart
bestemt Begyndelse er altsaa umiddelbart betinget
ved Universalbegyndelsen. Som nu den phænomenale
Række er adsplittet i ydre Mangfoldighed,
i partielle Begyndelser, hvori der ingen Universalbegyndelse
er, saa er Universalbegyndelsen selv kun
forudsat som Grundeenhedens tilhyllede Mangfoldighed,
i hvilken der ingen Partialbegyndelse er, fordi
der endnu ikke er sat nogen existentiel Forskjel.
Overgangen fra den universelle til den partielle Begyndelse
kan desuagtet ikke skee ved et Spring; idet
den phænomenale Række synker som Skin i Urgrundens
Væsen uden at kunne begynde, saa forudsætter
den sit forudsættende Princip ɔ: den sætter alle sine
partielle Begyndelser tilbage i Universalbegyndelsens
concrete Inderlighed, og forudsætter denne prædisponeret
til at sætte den specifiske Forskjel. Forsaavidt
det bestemte Phænomen begynder sin Existents under
Forudsætning af særskilte Betingelser, gaaer Udvortesheden
og Andetheden tilbage i en uendelig Regres,
og Betingelsernes Betingelser bortflygte i det Ubestemmelige;
men forsaavidt denne udvortes Reflexion
i Rækken af de partielle Begyndelser tillige har sin
Forudsætning i den prædisponerede Universalbegyndelse,
reflecteres alle ydre phænomenale Betingelser
% --- p. 167 ---
\setcounter{footnote}{0}%
\opage{167}i Centralexistentsens eget Dyb. Betingelsernes immanente
Reflexion i Universalbegyndelsen og deres
udvortes Reflexion i Rækkens partielle Begyndelser
ere altsaa een og samme Reflexion. For at en Existents
kan begynde, kræves da, at dens særegne Betingelser
forefindes i Phænomenernes Række, dernæst
at den Afbrydelse og Fraværelse, der ligger i den
phænomenale Andethed, integreres ved den væsentlige
Nærværelse af Universalbegyndelsens immanente Betingelser,
der forudsættes i den existentielle Centralitet.
Som Betingelserne og de relative Grunde ere,
saaledes er ogsaa Centralexistentsen i sin Urgrund
disponeret til at sætte den nye Begyndelse; men
omvendt ere de ydre specifiske Betingelser og relative
Grunde Centralexistentsens egne Forudsætninger,
ved hvis Ophævelse den sætter det Ubetingede. De
endelige Grunde og phænomenale Betingelser udspinde
sig som en tautologisk Fortsættelse af den ene relative
Existents i den anden; men først naar Centralitetens
Lyn slaaer ned i Betingelsernes Kreds, bliver
Realbegyndelsen sat, og en ny Existents dukker
op af Mediationens Strøm.
Centralexistentsen har nu sin Modsigelse deri, at den
paa een Gang skal indeholde den i sig specificerede
Universalbegyndelse, Mediationen af alle Betingelser
i een absolut Ubetingethed, og saaledes subordinere
sig den phænomenale Progres, medens den dog tillige
maa forudsætte Phænomenernes umiddelbare
Adskillelse, for ved deres Ophævelse at vinde i existentiel
Gehalt, og disponeres til at fortsætte Phænomenernes Række. Denne Modsigelse finder sin
Løsning i den immanente Selvbevægelse, hvorved
Grundexistentsen uafbrudt løslader sin Andethed.
Det Absolutes Selvudvikling til Existents er en sig
fortsættende Explication af Urgrundens dobbelte Ind%
% --- p. 168 ---
\setcounter{footnote}{0}%
\opage{168}hold. Idet Grundeenheden oplader sit Mysterium,
og sætter sit Indhold til i Phænomenernes Cyclus,
da viser dette Indhold sig for det Første som dens
\textbf{eget} Indhold; thi det er nu blevet klart, hvorledes
den opløser Phænomenernes Existents, for at tage
de specificerede Indholdsmomenter tilbage i sin egen
existentielle Selvreflexion; dernæst er det et \textbf{andet}
Indhold, thi det har ligeledes viist sig, at de bestemte
Phænomeners Opløsning afgive Stof og Materiale
for andre Phænomener i den exoteriske Række.
Umiddelbarhedens „\textbf{Vorden}“ er i Centralexistentsen
sat som et Moment, der ophæves i Reflexionens
„\textbf{Har været}“, medens derimod, hvad den phænomenale
Progres angaaer, det væsentlige „Har været“
viser sig som en Forsvundethed, der idelig maa
erstattes paany i Partialbegyndelsens umiddelbare
Vorden; naar det existerende Phænomen, ifølge sin
rene Overfladiskhed, uafbrudt tager sin Begyndelse
fra Andet, deels fra andre Phænomener, og deels
fra Ur-Existentsens Dyb, saa tager Centralexistentsen
gjennem Andethedens ɔ: gjennem de partielle Begyndelsers
forudsatte Ophævelse sin immanente Realbegyndelse
fra sig selv; hvad der altsaa med Hensyn
til Rækken er uendelig Regres og Progres, det
er med Hensyn til Centralexistentsen uendelig Selvfordybelse
og Selvudvikling. Ved denne immanente
Selvudvikling vindicerer Centralexistentsen sin absolute
Dignitet, og subordinerer sig den phænomenale Række, saa at Centralitetens dialektiske Puls slaaer
i ethvert existerende Phænomen. Urgrundens dobbelte Grundlag krævede for sin Udvikling
en esoterisk og exoterisk Mediation, en primitiv
og consecutiv Grundbevægelse. Formedelst
Indholdets Dobbelthed er altsaa Mediationen og
% --- p. 169 ---
\setcounter{footnote}{0}%
\opage{169}Formationen bleven fordoblet, og denne udviklede
Formforskjel sat som Moment i den ene absolute
Formation. Indholdet udformer sin egen Modsætning
af Grundens Eenhed; Modsætningen er altsaa
ikke umiddelbart dualistisk. Men det samme Grundindhold,
der i Formationen sætter sin alsidig udviklede
Identitet, udvikler tillige denne Indholdsidentitet
i en dobbelt Formation. Den primitive og
consecutive Grundbevægelse ere i deres gjensidige
Reflexion begge absolute. Ved denne immanente
Formforskjel bliver Indholdet forskjelligt fra sig selv.
Den udvortes Formforskjel lader Indholdet i det
Væsentlige uforandret, den endelige Form forudsætter
et givet Indhold som Stof og Materiale; men
den immanente Form, der gjennemtrænger, former
og bestemmer alle Indholdsmomenter, giver derved
Indholdet sin absolute Bestemthed, d. e. den bestemmer
det som sit bestemte Indhold, saa at det ene
Indhold kun formedelst Formen kan adskille sig fra
det andet. Har den absolute Existents altsaa en i
sig redupliceret Formation, har dens Udvikling tvende
immanente Urformer, da har den ogsaa i sin concrete Indholdsidentitet tvende modsatte Indhold.

Kun i Mediationsbevægelsen kan Centralexistentsen
adskille sig fra de phænomenale Existentser; men
idet nu den alsidige Existentialmediation er gjennemført,
vender Umiddelbarheden tilbage. Centralexistentsen
existerer i Mediationen af alle det Absolutes Forudsætninger
i Identiteten af alle phænomenale og immanente
Betingelser, den existerer umiddelbart som det
absolut Ubetingede. Men det Samme gjælder om den
phænomenale Existents; ogsaa den tager sin Realbegyndelse
i alle Forudsætninger af det Absolute; dens
særskilte Kreds af Betingelser er ligeledes medieret
% --- p. 170 ---
\setcounter{footnote}{0}%
\opage{170}gjennem alle Centralexistentsens immanente Betingelser,
ogsaa den existerer i sin Umiddelbarhed som det absolut
Ubetingede. Den absolute Form er da, som Følge
heraf, sat i ethvert Indhold, og omvendt er ethvert
Indhold eiendommelig bestemt ved den alsidig bestemmende
Form. Saaledes er det dobbelte Grundindhold,
formedelst den ene bestemmende Form, umiddelbart taget,
eet og samme Indhold, og dog er det ligesaa
umiddelbart et aldeles forskjelligt Indhold paa Grund
af de utallige Formbestemmelser, der sættes i den absolute
Formation. Det Absolute er i alle Formationer
stedse identisk med sig selv, og derfor eet eneste i
sig Ubetinget; men er det exoteriske Indhold ubetinget
ophævet i det esoteriske, formedelst den umiddelbare,
existentielle Eenhed, saa er det esoteriske ligesaa ubetinget
forsvundet i det exoteriske paa Grund af den umiddelbare,
existentielle Mangfoldighed og Forskjel. Idet
den absolute Ubetingethed lader sig tilsyne i Centralexistentsen,
kommer den med samme Slag til Gjennembrud
i alle relative Existentser, og tindrer i utallige
Centra. I denne umiddelbare Identitet af den centrale
og phænomenale Existents lader \textbf{det} Existerende sig
tilsyne, deels som \textbf{Tingen} i \textbf{sig}, deels som \textbf{Tingene}
i \textbf{deres} Egenskaber.

Efter sin umiddelbare I-sig-væren er Centralexistentsen
at fastholde som den abstracte Væsensexistents,
som den blinde Transcendents, det rene Negativum af
alle Phænomener; som Følge heraf kan den da heller
ikke nærmere bestemmes; thi alle Bestemmelser forsvinde
som Reflexer af en ydre Mangfoldighed, der slet
ikke vedkommer dens egen væsentlige Reflexion i sig;
ifølge denne abstracte For-sig-væren er det Existerende
altsaa sat som det Transcendente, som Tingen i sig
(„Das Ding an sich“). Naar det hedder, at Tingen
% --- p. 171 ---
\setcounter{footnote}{0}%
\opage{171}i sig ikke kan erkjendes, da vil dette med andre Ord
sige, at der slet ikke er Noget i den at erkjende, efterdi
den udelukker al Forskjel og Relation; Tingen i sig
er saaledes at betragte som den existerende Tautologie,
hvorved Sætningen: A \textbf{er} = A forvandles til:
A \textbf{existerer} = A. At nu Tingen i sig, trods denne
Tautologie, desuagtet ikke aldeles forsvinder og bliver
til ingen Ting, skylder den sin Relation til de fornegtede
Phænomener. Som Eenheden og Identiteten tilkommer
Tingen i sig, saa falder derimod Mangfoldigheden
og Forskjellen over paa Phænomenerne; disse ere
da ikke blot forskjellige fra Tingen, men ogsaa fra hverandre
indbyrdes. Men paa den anden Side er den
existerende Identitet stedse betinget ved den existerende
Forskjel; Tingen i sig bestemmes selv derved, at den
udelukker alle phænomenale Bestemmelser. Den distingverende
Reflexion, hvorved Phænomenerne adskilles fra
Tingen i sig, er altsaa den selvsamme, hvorved de adskille
sig fra hverandre. Det, der giver ethvert Phænomen
den existerende Identitet, den ubetingede For-sig-væren,
der holder sig ligegyldig mod enhver Relation
til de øvrige Phænomener, det i Phænomenet Existerende,
der bliver tilbage, naar vi abstrahere fra alle
relative Bestemmelser, er følgelig ogsaa en Ting i sig.
Skal Identiteten ubetinget existere i og for sig selv, da
maa Forskjellen udelukkes, men ved denne Exclusion
fixeres Forskjellen selv i umiddelbar For-sig-væren,
og kommer derved til at existere som umiddelbar Identitet.
Den ene Ting i sig bliver saaledes med eet
Slag til en Mangfoldighed af Ting i sig.

Tingens Reflexion i dens For-sig-væren reflecterer
tillige dens Væren i Andet. Dette Andet er nu selv
en Ting; at Tingen reflecteres i sig, som i sit Andet,
vil da sige, at den reflecteres i en anden Ting. For%
% --- p. 172 ---
\setcounter{footnote}{0}%
\opage{172}saavidt de Existerende i deres Mangfoldighed ere Ting,
holde de sig, hver især, i Forskjellens og Modsætningens
Extrem; men idet de bestemmes ved at bestemme
hinanden gjensidig, identificere de tillige deres existentielle Forskjel: i denne gjensidige Reflexion \textbf{have} Tingene
\textbf{Egenskaber}.

Egenskaben er for det Første Tingens \textbf{egen} Bestemthed:
Tingens Bestemmelse i Relation til Andet
kommer i Egenskaben tilsyne som dens egen umiddelbare
Selvbestemmelse: som Egenskaben er, saaledes er
ogsaa Tingen beskaffen; dernæst er Egenskaben et Udtryk
af Tingens Bestemthed i Relation til en anden
Tings Egenskab. Hvad der gjælder om Tingene selv,
gjælder da ogsaa om Egenskaberne: deres Reflexion-i-Andet
er deres egen Reflexion-i-sig. Ifølge denne
Dobbeltreflexion have Tingene og Egenskaberne, hver
for sig, en vis Selvstændighed, en relativ Bestaaen.

Ethvert Existerende er altid et Noget; men ethvert
Noget kan ikke derfor siges at existere. Hvad „Noget“
er i den umiddelbare Væren, det er \textbf{Tingen} som
det phænomenale Væsen; \textbf{Noget} er, som saadant, en
Umiddelbarhed, der endnu ikke har gjennemløbet Reflexionen,
\textbf{Tingen} er derimod en Væsensexistents, hvis
Umiddelbarhed reflecterer den ophævede Reflexion; Noget er, som det Tilværende, den rene Negation af
\textbf{Andet}; \textbf{Tingen} er derimod, i Identiteten med sin
Beskaffenhed, desuagtet forskjellig fra \textbf{Egenskaben};
Noget \textbf{er} Qvalitet, men Tingen har Egenskab, deri
ligger Forskjellen.

En væsentligere Lighed finder Sted imellem Tingen
og \textbf{Phænomenet}. Det existerende Phænomen
er den umiddelbare Eenhed af Ting og Egenskaber;
forsaavidt Tingen er Phænomen, kan den opløses i
Materier, og gjennemløbe en Række af endelige For%
% --- p. 173 ---
\setcounter{footnote}{0}%
\opage{173}mer: \textbf{Tingen bestaaer af Materier}\footnote{Hegels Logik, 2tes Buch. Pag. 155 (Werke 4de Band) og J. L. Heibergs speculative Logik Pag. 51.}; men Phænomenet
staaer desuagtet i existentiel Udvikling under
Tingen, thi den phænomenale Ubetingethed har endnu
ikke reflecteret \textbf{Tingens og Egenskabernes} væsentlige
Forskjel. Paa Grund af sin rene Overfladiskhed
kan den phænomenale Existents ikke tabe en
eneste af sine Bestemmelser, uden at den derved tillige
altereres; Phænomenet er i denne Henseende et Noli
me tangere; imellem Ting og Egenskaber finder derimod
den Adskillelse Sted, at Tingen selv er sat som
det Væsentlige, medens Egenskaberne i deres indbyrdes
Forskjellighed vise sig som det Uvæsentlige; derfor kan
Tingen meget vel miste en af dens Egenskaber, uden
at den derfor strax opløses\footnote{Naar et firkantet Bord forandres til et rundt, bliver det et andet Phænomen, men er desuagtet den samme Ting. En Blomst uden Duft er ogsaa en Blomst.}. At Tingen har Egenskaber,
udtrykkes ogsaa derved, at den \textbf{bestaaer} i sine
Egenskaber, eller, hvad der er det Samme, at den bestaaer
i Egenskabernes Bestaaen. Forsaavidt her nu
er en dobbelt Bestaaen, gaaer Tingens Existents ikke
op i nogensomhelst enkelt Egenskab; men da den dobbelte
Bestaaen desuagtet grunder sig i en Identitet,
kan Tingen, som saadan, dog heller ikke taale en Forandring
af alle sine Egenskaber, uden at den derved
selv tillige maa altereres.

Under den skarpe Sondring af Ting og Egenskaber
driver den distingverende og fixerende Reflection sit allerfrieste
Spil. I Forskjellen mellem Tingene og Egenskaberne
fastholdes Tingenes egen indbyrdes Forskjel.
I Egenskabernes Forandring foregaaer der vel en
For%
% --- p. 174 ---
\setcounter{footnote}{0}%
\opage{174}andring med Tingene; men da disse ved at bestaae i
Egenskaberne tillige have en særskilt Bestaaen i sig selv,
saa ere de ikke \textbf{for hverandre} indbyrdes, og følgelig
uafhængige af de i Egenskaberne indtrædende Forandringer.
I et saadant Selvstændighedsextrem kunne
Tingene fixeres, just fordi de have Egenskaberne, i hvis
Bestemmelse det netop ligger, at lade Tingene være
for hverandre, og derved reflectere den Strøm af Forandringer,
i hvilke Tingene selv holde sig uforandrede.
Som den hele Vexel af relative Bestemmelser er Tingene
selv en ligegyldig og udvortes Reflexion, saaledes
ere ogsaa Tingene ligegyldige mod hinanden indbyrdes;
de kunne mødes i de allerforskjelligste Relationer, uden
at deres faste Bestaaen lider derved. De i Egenskaberne
reflecterede Forandringer vedrøre altsaa ikke Tingenes
Væsen, men kun deres phænomenale Side.

Fastholdes derimod Identiteten af Ting og Egenskaber,
da bliver Tingenes Væsen reflecteret i dens Phænomen,
og begges Eenhed sat som Tingenes sande Existents.
Den Forskjel, der kommer tilsyne imellem Væsen
og Phænomen, falder da sammen med Tingenes
egen Bestaaen, og som Tingen selv existerer i umiddelbar
Eenhed af dens væsentlige og uvæsentlige Side,
saaledes have ogsaa Egenskaberne i sig en væsentlig og
uvæsentlig Existents. Ved at tabe en uvæsentlig Egenskab
lider Tingen selv en uvæsentlig Forandring, hvorimod
den væsentlig altereres, naar en væsentlig Egenskab
forsvinder. Tingenes og Egenskabernes Reflexion
er altsaa i Grunden een og den selvsamme, og det
Eiendommelige i Reflexionen ligger deri, at den baade
\textbf{er} ophævet og dog ikke ophævet.

Er Tingens Reflexion-i-sig \textbf{Eet} med dens Reflexion
i Egenskaberne, saa falder Tingens concrete Bestemthed
umiddelbart sammen med dens Egenskaber, og disse
% --- p. 175 ---
\setcounter{footnote}{0}%
\opage{175}forsvinde som Reflexer i Tingens „Hæcceitet“, Phænomenet er da negeret som et blot Phænomen, og derved
ophævet i Bestemmelsen af sit „Noumenon“: Tingen
\textbf{lader sig} tilsyne i sine Egenskaber. Men da Egenskabernes
Reflexion i Tingen nu ogsaa er \textbf{Eet} med
deres egen Reflexion i sig, falder Selvstændigheden over
paa Egenskaberne; Tingens bestemte Phænomen bliver
saaledes sat i dets Forskjel fra Tingen selv, og denne
derimod forudsat som et i sig ubestemmeligt „Noumenon“:
Tingens Existents i Egenskaberne er \textbf{kun} tilsyneladende.
Den existentielle Antinomie, at Tingen paa een
Gang baade er transcendent og immanent, ligger i
Tilsyneladelsen \textbf{selv}.

Tingene i deres phænomenale Bestemthed ere hinanden
umiddelbart modsatte, og danne saaledes en Række
af existerende Phænomener. Kunde den væsentlige Existents
udtømmes i phænomenal Umiddelbarhed, da vilde
Tingen selv være identisk med sin Tilsyneladelse, og
enhver Illusion forsvinde. Men nu er Tingen i sit
Phænomen tillige forskjellig fra Phænomenet, den er
da, som Følge heraf, sin egen Tilsyneladelse modsat, og
dens umiddelbare Bestaaen i Phænomenet illusorisk.
Det Væsentlige i Tilsyneladelsen, nemlig Tingen i sig,
er altsaa just \textbf{Det}, der ikke lader sig tilsyne; derved
ophæver Tilsyneladelsen selv sin egen Væsentlighed.
Paa den anden Side bliver denne Væsentlighed tillige
sat i dens Ophævelse; thi hvor forskjellig Tilsyneladelsen
end er fra Tingen selv, saa lader den dog Tingens
abstracte For-sig-væren komme tilsyne; hvad den derimod
synes at skjule, er Tingens egen concrete Bestemthed.
Nu har det imidlertid viist sig, at Hæcceiteten,
hvorved Tingen er bestemt i sig, og forskjellig fra enhver
anden Ting, nødvendigviis reflecteres i en Cyclus
af væsentlige og uvæsentlige Egenskaber. Enten maa
% --- p. 176 ---
\setcounter{footnote}{0}%
\opage{176}da Tingen i sig forblive en tom Abstraction, der mangler
al Existents, efterdi den mangler Tilsyneladelse,
eller ogsaa dens hiinsidige Existents selv maa blive en
Tilsyneladelse i hiinsidige Phænomener. Illusionen, der
ligger i den abstracte Modsætning mellem Ting og
Phænomen, forsvinder, idet den lader sig tilsyne, eller
med andre Ord: Tilsyneladelsens Antinomie er kun tilsyneladende.

I Tilsyneladelsen culminerer den existentielle Reflexion.
Det Existerende er kun et Tilsyneladende; thi det bestaaer
altid i en anden Bestaaen; denne anden Bestaaen
er selv negativ, og bestaaer tillige i sin Ophævelse
som et Tilsyneladende. Men da enhver Existents, just
fordi den er tilsyneladende, bestaaer i Negationen af
dens Bestaaen, og det ene Existerende er, i Forhold til
det andet, denne Negation, saa bestaaer den ene Existents
i og ved den anden, og erholder saaledes, paa
Grund af Tilsyneladelsen, en i sig reflecteret Bestaaen.
α) Tingen i sig lider af den Modsigelse, at der skal
være en abstract Selvstændighed uden Skin af Reflexion,
og just derved bliver et abstract Skin uden al Selvstændighed.
Tingens Eenhed i sig er kun en Skineenhed,
thi den existerer i en Mangfoldighed af Ting i sig;
men Tingens væsentlige Mangfoldighed er ogsaa en
Skinfleerhed, thi den existerer kun i den ene ureflecterede
Ting i sig. Tingens absolute Væsenhed er da et i sig
selv tilsyneladende Væsen, og maa derfor lade sig tilsyne
i \textbf{sit Andet} ɔ: i Egenskaber og Phænomener.
8) Men ved at komme tilsyne i Egenskaberne, geraader
Tingen selv i en ny Modsigelse. Dens Existents oscillerer
nu imellem de to Extremer, enten at sætte sin
Selvstændighed til i Egenskabernes Bestaaen, eller at
absorbere Egenskaberne i sin Bestaaen. Skeer det Første,
saa udtømmes Grundexistentsen i de existerende 
% --- p. 177 ---
\setcounter{footnote}{0}%
\opage{177}mener, og da Phænomenerne negere enhver Modsætning imellem en væsentlig og uvæsentlig Existents, forsvinde Tingene selv i Materiens vexlende Former; skeer det Sidste, saa falde Tingenes ubestemmelige \textit{Noumena} tilbage i en blot tilsyneladende Væsenhed, og henflyde som uselvstændige Momenter i den absolute Centralexistents. \textit{γ}) Denne Modsigelse opløser sig i Phænomenernes egen tilsyneladende Væren; thi ogsaa denne bestaaer i Reflexionen af abstract Selvstændighed og Uselvstændighed. Som Tingene først komme tilsyne i sig ved at lade sig tilsyne i Phænomenerne, saa blive Phænomenerne omvendt, ved at lade sig tilsyne i hverandre, Tingenes egne Tilsyneladelser. Tilsyneladelsen er den \textbf{satte} Eenhed af Existentsens Identitet og Forskjel. Forskjellen existerer, thi enhver Existents har sin særlige Bestaaen, dens bestemte Form og Indhold adskiller sig fra andre Indhold i andre Former; dens Selvstændighed udelukker enhver anden Selvstændighed, den er, som saadan, umiddelbar og positiv. I den absolute Negativitet existerer altsaa en Utallighed af positive Indhold, i den absolute Selvstændighed bestaae utallige Selvstændige: den existentielle Fleerhed \textbf{er} den Form, i hvilken \textbf{Eenheden \textbf{lader sig tilsyne}.} Men i Fleerhedens Positivitet ligger dens Negativitet; den existerende Forskjel er i sig selv tilsyneladende, derfor kommer Identiteten til Existents. Da Modsætningen imellem de forskjellige Indhold er i sig tilsyneladende, maa den \textbf{lade sig \textbf{tilsyne}, \textbf{som en} Mangfoldighed} af Forskjelsmomenter, i den absolute Indholdsidentitet; der existerer altsaa kun eet absolut Indhold i alle relative Indhold, kun een absolut Form i alle relative Former, kun een absolut Positivitet i al Negativitet, kun een absolut Selvstændighed i alle Selvstændige. Idet nu begge Former for Alexistentsen ere i Tilsyneladelsen
% --- p. 178 ---
\setcounter{footnote}{0}%
\opage{178}satte med samme ubetingede Gyldighed, saa har den existentielle Modsætning imellem Centralitetens Dyb og de overfladiske Phænomener ophørt at være en Modsigelse. Som den immanente og væsentlige Tilsyneladelse bestaaer i Mediationen, hvis Aleenhed er Dybet, saa existerer den uvæsentlige Tilsyneladelse i Umiddelbarheden, hvis phænomenale Forskjel er overfladisk. Tilsyneladelsen er selv Identitet af Mediation og Umiddelbarhed. Tingene i deres Phænomener ere umiddelbare Existentser; de komme tilsyne som Tilværende. Tingens Tilværen ligger ogsaa i dens Tilsyneladelse; thi hvad der skal lade sig tilsyne, maa alt have gjennemgaaet sin Vorden. Men denne Tilværen er tillige væsentlig, og idet Vorden reflecteres i Tilværen, kommer Oprindelsen tilsyne. Tingens Modsætning til dens Phænomen og umiddelbare Tilsyneladelse er herved reflecteret som en væsentlig Forskjel imellem dens \textbf{Tilværelse} og dens \textbf{Oprindelse.} Ifølge Tilsyneladelsens Umiddelbarhed synes den ene Ting at oprinde af den anden, og dog kommer Oprindelsen ikke selv tilsyne; thi Mediationen, der her viser sig som det urolige, overfladiske Omslag af forskjellige Indhold og Former, brydes uafbrudt med Umiddelbarheden, og er som saadan blot tilsyneladende. Idet Tingenes phænomenale Existents skal gaae sammen med sig selv til absolut Bestaaen, støder den sig selv fra sig. Men ved at lade sig tilsyne i sin egen Andethed ɔ: i sin blot tilsyneladende Oprindelighed, lader den tillige sin \textbf{væsentlige Oprindelse,} det absolute Andet, hvoraf den selv er sat, komme tilsyne, og bestaaer i sin existentielle Mangfoldighed som et Resultat, der forudsætter sit forudsættende Princip ɔ: \textbf{Centralitetens \textbf{Dyb}.}
% --- p. 179 ---
\setcounter{footnote}{0}%
\opage{179}I Centralexistentsen er Umiddelbarheden omvendt  et
Gjenskin af Mediationen. Det existentielle Dyb kan
ikke umiddelbart komme tilsyne, thi derved vilde det
selv forvandles til et overfladisk Phænomen, ligesaalidt
kan det være en Ting i sig, da Tingens Forskjel fra
Phænomenet selv er en Reflex af den i sig reflecterede
Væsenhed; Centralexistentsen existerer altsaa i Form af
væsentlig Mediation. Ifølge denne Mediation er dens
umiddelbare Selvstændighed dens egen uendelige Negativitet;
herved er dens Indhold forudsat som Grundens
absolute Oprindelighed; men da denne Oprindelighed
er i sig tilsyneladende, saa giver den Tingene deres
Oprindelse, og sætter den ydre existentielle Mangfoldighed
som det Medium, hvori den lader sig tilsyne; som
negativ Existents bestaaer den altsaa i sin egen Negation
ɔ: i Forskjellens og Mangfoldighedens positive
Umiddelbarhed; men da denne ydre umiddelbare Existents nu ogsaa bestaaer i sin Negativitet, da den selv
kun er en Tilsyneladelse af den tilsyneladende Eenhed,
saa er dens Selvstændighed sat i dens Ophævelse. Ved
at bestaae i den ydre Existents bestaaer Centralexistentsen
i Ophævelsen af den existentielle Mangfoldighed;
idet den herved negerer sin egen Negation, existerer den
altsaa i Tilsyneladelsen af alle umiddelbare Existentser,
som den i sig affirmative, totaliserede Identitet. Naar
Centralexistentsen da bestaaer i Mediationens positive
Eenhed, der i sit Indhold ophæver alle Forskjelsmomenter,
som negative og tilsyneladende, saa bestaaer
Phænomenalexistentsens Indhold i en umiddelbar Forskjel,
hvis Eenhed i sig er tilsyneladende og negativ.
Ved at sætte denne dobbelte Bestaaen i een alsidig
Tilsyneladelse, gaaer den absolute Existents over i
Totalitet.
% --- p. 180 ---
\setcounter{footnote}{0}%
\opage{180}\underafsnit{c) \textbf{Totalitet}.}
\paragraf{26}{Indre og Ydre, Kraft og Product, Præformation.}

I Tilsyneladelsen komme alle Existentialmomenter
tilsyne, saavel i deres Umiddelbarhed, som i
deres gjensidige Overgang og immanente Reflexion,
hvorved Alexistentsen sættes som \emph{Totalitet}. Efter
sin væsentlige Side er Totaliteten bestemt som
et Rige af Love, efter sin Umiddelbarhed derimod
som en Verden af Phænomener: en abstract Modsætning, der ophæves i det væsentlige Forhold ɔ: i
Reflexionen af den reflecterede og umiddelbare Totalitet. Ifølge den totaliserede Formforskjel, der
tillige er en Forskjel i Indhold, forudsætter det væsentlige
Forhold den i Mysteriets Inderlighed og
Selvhedens Dyb reflecterede Totalitet paa den ene
og den i en udvortes Række udfoldede Umiddelbarhedstotalitet
(Hele og Dele) paa den anden Side,
og bestemmer sig derved til et Forhold imellem det Indre og det \emph{Ydre}. Den med sine ydre Former
incommensurable Inderlighed ytrer sig i uudtømmelig Productivitet; Identiteten af Indre og Ydre reflecteres i Kraften, og Totaliteten udvikler
sig til et Forhold imellem Kraft og \emph{Product}. Som Formationen af en ny Gestalt forudsætter den
producerende Kraft, saa forudsætter Kraftens specifiske
Productivitet producerede Gestalter i det Uen%
% --- p. 181 ---
\setcounter{footnote}{0}%
\opage{181}delige, saa at enhver Realformation reflecterer sin
\textbf{Præformation}. Ifølge Tilsyneladelsens Reflexion er ethvert Moment i det Absolutes Existents saaledes gjennemkla-

ret, at Grundexistentsen kan sættes som Medium for
Phænominalexistentsens uendelige Progres, ligesom % PRINTED AS IS: Phænominalexistentsens (sweep reading)
den phænomenale Existents omvendt bliver sat som
det Medium, hvorigjennem Grundexistentsen vinder
sin Selvfordybelse: det esoteriske og exoteriske Indhold
gjennemskinne hinanden. Som der altsaa Intet
gaaer tabt, naar kun Grundexistentsen fastholdes, efterdi
den selv indeholder alle Phænomenalexistentsens
Momenter, og sætter dem ud af sig, for i deres Negativitet
at affirmere sin egen Selvstændighed, saa
tabes der heller intet væsentligt Led af den absolute
Existents, hvis Phænomenalexistentsen sættes i sin
Selvreflexion. Thi idet den er sat af Andet, reflecterer
den fra alle Sider sin absolute Andethed, og er saaledes
gjennemtrængt af Grundexistentsens Indhold, at Dybets
Eenhed lader sig tilsyne i dens overfladiske
Mangfoldighed. Ved at sættes i Form af Tilsyneladelse
er Alexistentsen udviklet til \textbf{Totalitet}.

Totaliteten er da nu den meest adæqvate Form for
det Absolutes Identitet i alle selvstændige Forskjelle.
Det ligger i Totaliteten, at hver Existents har sin
Bestaaen, at Sondringen og Adskillelsen gjennemføres;
men det følger paa den anden Side ogsaa af
Totalitetens Væsen, at ingen Existents kan isoleres
til abstract For-sig-væren og falde udenfor Heelheden.
Begge Totalitetens Fordringer ere fyldestgjorte i Tilsyneladelsen,
der paa een Gang sætter enhver Selvstændighed
som uendelig \textbf{negativ} og uendelig positiv.
% --- p. 182 ---
\setcounter{footnote}{0}%
\opage{182}Men for saavidt Totaliteten kan udvikles af en
hvilkensomhelst Existents, er den endnu, ifølge Tilsyneladelsen,
aldeles formel, og mangler den særegne
Bestemthed. I den rene Tilsyneladelse ere alle Existentser
ophævede: de ere derved alle lige endelige og
uendelige, lige middelbare og umiddelbare, lige
centrale og peripheriske. Udgangspunktet for Totalitetsudviklingen er som Følge heraf aldeles ligegyldigt.
Tilsyneladelsens Categori giver sig som en
gjennemsigtig Totalitetsform; men Totaliteten selv
kræver en nærmere Bestemmelse.

Totalitetens rene Extremer ere: paa den ene Side
Tilsyneladelsens immanente Væsen, en Identitet, der
ophæver enhver phænomenal Forskjel, og affirmerer
Totaliteten som Tingenes altomfattende og altbestemmende Eenhed, paa den anden Side Tilsyneladelsens
Selvadskillelse i Phænomenernes concrete Umiddelbarhed,
hvorved Sondringen fuldbyrdes, og Totaliteten
viser sig i den hele Mangfoldighed af særlige Existentser.
Da fremdeles enhver af disse Sider med
al sin Selvstændighed er et Tilsyneladelsens Moment,
og følgelig selv kun har en tilsyneladende Selvstændighed,
saa ere de ikke udvortes sammensatte, ei heller behøves der et Tredie, for at danne deres Forening:
Tilsyneladelsens egen Reflexion udgjør deres
sande Bestaaen. Begge Siders reflecterede Bestaaen
bliver nu nærmere at bestemme som \textbf{Loven} og Hæcceiteten.

a) I Totaliteten af sine væsentlige og uvæsentlige
Egenskaber, d. e. i sin umiddelbart bestemte Tilsyneladelse, er Tingen sat med dens Hæcceitet; men
som \textbf{denne} Ting (hæc res) er den derved tillige
begrændset og forskjellig fra alle andre bestemte
Ting, af hvilke enhver især udtrykker sin Hceccei%
% --- p. 183 ---
\setcounter{footnote}{0}%
\opage{183}tet. Den hele Totalitet udpræger sig 1 Hæcceiteternes udvortes, uendelige Mangfoldighed. I
Tingenes umiddelbare Bestemthed og væsentlige
Begrændsning ligger deres Endelighed; den ene
Ting opstaaer og forgaaer efter den anden: Hæcceiteten
er forgængelig og omskiftelig. Men som
Tingenes umiddelbare Selvstændighed er tilsyneladende,
saa er deres Foranderlighed og Omskiftelighed
ogsaa tilsyneladende; hvormange Ting der end
opstaae og forgaae, en Utallighed af bestemte Ting
vil der dog altid blive. Antallet af de bestaaende
Ting er uendeligt, deres Antal, der opstaae og
forgaae, er følgelig ogsaa uendeligt: det er den
hele Totalitet, der uophørlig opstaaer, og den hele
Totalitet, der uophørlig forgaaer; thi i Væsenet
er Qvantiteten ophævet, den absolute Reflexion
regner med uendelige Størrelser. Da Hæcceitetens ydre Mangfoldighed og Forskjel saaledes i det
Uendelige forvandles til et ligegyldigt og uvæsentligt
Skin, saa er Tingenes Foranderlighed negeret
og tillige optagen i deres positive Uforanderlighed,
og Loven sat som deres \textbf{væsentlige} Grundlag,
som deres \textbf{immanente Bestaaen}.

Ved at sættes som Tingenes Grundlag faaer
Loven en vis Lighed med den væsentlige Ting,
med „Tingen i sig“, da begge fremstille sig som
de bestemte Tings sande, almindelige Væsen. Desuagtet
kan Loven ikke hypostaseres til Ting, efterdi
den negerer enhver selvstændig Umiddelbarhed.
Naar „\textbf{Tingen} i \textbf{sig}“ er bestemt som det værende
Grundlag, der forholder sig ligegyldigt til
Existentsens Phænomener, saa har „\textbf{Loven}“ derimod
sin Bestaaen i \textbf{de} bestaaende Ting; „Tingen
i \textbf{sig}“ afviser Tilsyneladelsen, og viser sig
% --- p. 184 ---
\setcounter{footnote}{0}%
\opage{184}derved som det \textbf{transcendente} Væsen; „\textbf{Loven}“
lader sig derimod tilsyne i Tingenes Tilsyneladelse,
og gjælder kun som det \textbf{immanente} Væsen.
Hvad „\textbf{Tingen} i \textbf{sig}“ er i Form af \textbf{ophævet}
Reflexion, det er \textbf{Loven} i Form af \textbf{sat} Reflexion;
naar Tingen i sig ved en hiinsidig abstract Umiddelbarhed støder de phænomenale Bestemmelser
bort fra sig, saa kommer Loven selv aabenlyst
frem, som Væsenheden i Phænomenernes Vexel.\footnote{Uagtet Kant demonstrerede, at vi ikke kunne erkjende „das Ding an sich“ (det transcendente Væsen), saa var han dog overbeviist om, at vi kunne efterforske Tingenes Love; thi Loven (det transcendentale ɔ: tilsyneladende og immanente Væsen) gjælder ikke \textit{Noumena}, men \textit{Phænomena}.} Naar Tingenes vexlende Mangfoldighed sættes
som Reflex i deres immanente Eenhed, saa er herved
tillige forudsat, at denne Væsensidentitet opgaaer
som Reflex i den bestemte Umiddelbarhed.
Da nu Loven selv er Tingenes Væsen, og da Tingenes
selvstændige Umiddelbarhed just fremgaaer
af deres egen Væsenhed, saa udspringer Hæcceiteten
af Lovens egen sættende Reflexion. Denne
Modsigelse, at Hæcceiteten paa een Gang fremgaaer
af Loven og dog falder udenfor Loven, paa
een Gang bestemmes ved Loven og dog er Loven
ligegyldig, med eet Slag sættes og ophæves af
Lvven, udtrykkes i den Sætning, at \textbf{Tingen er % PRINTED AS IS: Lvven (for Loven; S211)
Loven undergiven}. Tingens og Lovens Bestaaen
er herved fordoblet: Tilsyneladelsen er Tingens
Sandhed, den lader derfor Lovens Væsen
komme tilsyne som den Andethed, i hvilken den
selv bestaaer; desuagtet bestaaer den ikke som Lov,
% --- p. 185 ---
\setcounter{footnote}{0}%
\opage{185}men som umiddelbar Selvstændighed. Hvad
Loven angaaer, saa har ogsaa den Hæcceiteten
som det Andet, i hvilket den kommer til Existents,
men bestaaer dog ikke som \textbf{denne} Ting,\footnote{Stenen er tung; den har altsaa en væsentlig Bestaaen i sin Tyngde, thi hvordan den forresten er beskaffen, saa er den dog Tyngdens Lov undergiven. Imidlertid existerer den dog ikke som Tyngde, men som Steen ɔ: som det Andet, hvori Tyngdens Lov kan komme tilsyne. Hvis Nogen seer et Menneske døe blot som dette Menneske, da sandser han kun et Factum, hvis hans Tanke derimod seer Menneskeheden døe i dette Menneskes Død, erkjender han Dødelighedens Lov.} men
som dens Negation. Begge Hovedsider ere saaledes,
ifølge Tilsyneladelsens Reflexion, lige positive og
lige negative, men ifølge den umiddelbare Forskjel
imellem Lov og Hæcceitet hinanden abstract modsatte.
Sættes nu Tingens Bestaaen i Loven som
det Positive, saa er Hæcceiteten et forsvindende
Skin, et Negativum, gjennem hvilket Totalitetens
\textbf{Eenhed} affirmeres; sættes Tingens Bestaaen i
umiddelbar Selvstændighed derimod som det Posir
tive, saa er Loven den Negation, hvorved Hæcceiteten
bestemmes, og Totalitetens \textbf{Alhed} affirmeres.
Medens Loven altsaa i første Henseende er
forudsat som Tingenes Grundlag, saa er den i
sidste Henseende sat som Tingenes \textbf{immanente
og væsentlige Grændse}. Tingenes Bestemthed
er deres i Loven reflecterede Endelighed, der
saaledes bliver sat i en dobbelt Form. I Form
af væsentlig Identitet kommer Tingenes Endelighed
tilsyne som en Fælledsforsvinden i Lovens
Væsen, som i deres rene Uendelighed og
Alminde%
% --- p. 186 ---
\setcounter{footnote}{0}%
\opage{186}lighed: \textbf{de ere alle lige} \textbf{for Loven}; i Form
af væsentlig Forskjel kommer Endeligheden tilsyne
som den Grændse, Tingen ikke kan overskride, uden
at dens Bestaaen altereres; Loven sondrer og adskiller
Tingene: \textbf{hertil og ikke længere}! er det
Lovbud, der gjælder hver enkelt Existents. Tingenes
Adskillelse og Overgang i hinanden gjør dem
til Sider i den immanente Relation, der er bestemt
som deres Lov. Men Tingene ere tillige meer end
blotte Sider i Forholdets Lov; de ere Hæcceiteter, hvis umiddelbare Selvstændighed viser sig ligegyldig
imod den gjensidige Relation. --- Mediet
imellem Lov og Hæcceitet er den væsentlige Egenskab.
Som \textbf{Egenskab} er den nemlig et Moment
i Hæcceiteten, en Bestemmelse i den for-sig-værende
Umiddelbarhed, hvorved Tingen selv falder
udenfor Loven og Forholdet; som \textbf{væsentlig} Egenskab
er den selv ligegyldig mod Hæcceiteten, reflecterer
sin Uforanderlighed i dennes Foranderlighed,
og viser Tingens Bestemthed ved og Afhængighed
af Loven. Da nu den samme Ting har en Fleerhed
af Egenskaber, da enhver Egenskab viser, at
Tingens Reflexion i sig er sat som dens Reflexion
i andre Ting, og Tingens Selvstændighed som
Side i den immanente Reflexion, saa kan enhver
Egenskab for sig sættes som væsentlig, medens de
øvrige ophæves som uvæsentlige, hvorved Tingens
Hæcceitet viser sig at være en Fleerhed af forskellige
Love underkastet.\footnote{Mennesket er fornuftigt, og som saadan underkastet Pligtens Love; Mennesket er legemligt, og ifølge denne Egenskab Tyngdens Love underkastet.} Ved at reflectere sin
væsentlige Eenhed i Tingenes umiddelbare Forskjel,
% --- p. 187 ---
\setcounter{footnote}{0}%
\opage{187}kommer Loven tilsyne i en \textbf{Mangfoldighed af
forskjellige Love}..
I Tingenes umiddelbare Forskjellighed kommer Lovenes
væsentlige Forskjel tilsyne; som Tingenes
væsentlige Relationer udtrykke Lovenes Mangfoldighed,
saa fremstille de ydre, uvæsentlige Relationer
derimod Hæcceitetens Ligegyldighed mod Loven
selv. Begge Relationer ere imidlertid Tilsyneladelser,
og følgelig Momenter i een og samme Reflexion.
Skal den bestemte Lov komme tilsyne,
da forudsætter den et Sammentræf af givne Betingelser;
den reflecterer sin Væsenhed i ligegyldige
Relationer. Heraf følger Lovens \textbf{thetiske} og
\textbf{hypothetiske} Beskaffenhed.\footnote{I den experimenterende Physik gjælder det just om at træffe de rette Betingelser, for at Loven i sin rene Almindelighed kan komme tilsyne.} α) Loven
er abstract thetisk; idet den, som Tingenes eget
Væsen, mangler al indre Reflexion i sig selv,
mangler den tillige al Mediation, og giver ingen
Grund for sin Gyldighed. Desuagtet er Loven
dog en Væsenhed; den er Uforanderlighedens immanente
Reflexion i Tingenes Foranderlighed, og
kan som Følge heraf ikke afbrydes. Er Loven een
Gang given som Forudsætning, saa ere med det
Samme alle dens Conseqventser givne. Lovenes
Udvikling af Loven har saaledes en \textbf{analytisk}
Characteer. β) Men som det tilsyneladende Væsen
har Loven ingen Gyldighed udenfor dens Tilsyneladelse,
og udenfor Tingenes Existents aldeles ingen
Opfyldelse. Som Tingenes umiddelbare Bestemthed
er væsentlig bestemt og begrændset ved
Loven, saa er omvendt Lovens Bestemthed igjen
% --- p. 188 ---
\setcounter{footnote}{0}%
\opage{188}bestemt ved Tingenes Umiddelbarhed, og dens Opfyldelse
uendelig betinget. De modsatte Sider i
det lovbestemte Forhold maae factisk gives; de
sættes ikke af Loven selv; men de forudsættes, for
at Loven kan sættes. Derfor komme de mangfoldige
Love sporadisk tilsyne i Tingenes existentielle Mangfoldighed; den ene Lov kan ikke umiddelbart
deduceres af den anden; men Lovenes Tilbageførelse
til Loven maa skee ad den \textbf{synthetiske} Vei.\footnote{Den mathematiske Erkjendelse er paa een Gang analytisk og synthetisk.} γ)

Desuagtet er Lovens thetiske og hypothetiske Bestemthed
dog i det Væsentlige een og samme Reflexion.
At Loven kommer tilsyne i thetisk Form, hidrører fra
dens Mangel paa Selvfordybelse; dens Reflexion
i sig, hvorved den nivelerer enhver umiddelbar
Forskjel, er umiddelbart dens Reflexion i de existerende
Ting, hvorved den selv gaaer over i Fleerhed
og Forskjel. Med den ene thetiske Lov er da
given en Mangfoldighed af thetiske Lovbestemmelser.
Alle Theses falde her umiddelbart sammen,
thi de ere med eet Slag. Den ene Lov kan altsaa
ikke hente sin Betydning fra den anden; den ene
har ingen Gyldighed at laane den anden, men enhver
har Gyldigheden i sig selv. Adskillelsen imellem
Grundloven og de afledede Love ligger altsaa
slet ikke i Lovenes Væsen; thi enhver bestemt Lov
er Grundloven selv i særlig Bestemthed. Forsaavidt
den ene Lov har lige Gyldighed med den anden,
vise de sig indbyrdes ligegyldige; den ene synes
ikke at forudsætte den anden. Men da Lovenes
Mangfoldighed er en Tilsyneladelse af den
samme Thesis, saa er deres gjensidige 
Ligegyldig%
% --- p. 189 ---
\setcounter{footnote}{0}%
\opage{189}hed selv kun tilsyneladende, de forudsætte hverandre
indbyrdes: hvis een Lov kunde bortfalde, saa
kunde de alle forsvinde, hvis den ene Lov vedbliver,
saa bestaae de alle; ved denne gjensidige Forudsætten
slaaer den thetiske Bestemthed over i den
hypothetiske, og denne i hiin. Just fordi Thesis
saaledes paa een Gang baade er opgiven i Hypothesis,
og tillige sat som dennes Extrem, er Lovenes indbyrdes Sammenhæng af sig selv saa indlysende, og dog en Gaade.\footnote{I den pythagoræiske Læresætning er Thesis umiddelbart given med Hypothesis; desuagtet behøvedes der dog en Pythagoras for at deducere hiin af denne.}

Lovene ere Tingenes Baand; det Gaadefulde i
Lovenes indbyrdes Sammenhæng er tillige Hemmeligheden
i Tingenes væsentlige Forbindelse. At Tingen,
der i sin umiddelbare Selvstændighed kun synes
at beroe paa sig selv, desuagtet maa staae og falde
med andre fra dens ydre Berøring og umiddelbare
Omgivelse langt bortfjernede Existentser, udtrykker Forbindelsens Gaade i den meest overraskende Form.\footnote{Derfor hørte der en Newtons Genie til for at see Klodernes Lov i et Æbles Fald.}
Gaaden ligger i den fixerede Adskillelse imellem Lovenes
Væsen og Tingenes Existents. Det Medium,
gjennem hvilket Lovenes Forskjel reflecterer sin Identitet,
er endnu forudsat som en abstract Andethed, der
vel er Loven undergiven, men dog falder udenfor Loven.
Det er Tingenes umiddelbare Hæcceitet, der
træder imellem Lov og Lov, og derved skjuler Forbindelsen.
Af samme Grund danner nu ogsaa Loven
fra sin Side en abstract Modsætning til Tingene;
thi dens Eenhed og Forskjel holde sig udelukkende i
Reflexionen, og udvikle sig til en Mangfoldighed af
% --- p. 190 ---
\setcounter{footnote}{0}%
\opage{190}\textbf{væsentlige} Relationer. Tingenes Existents er derimod
nedsat i Umiddelbarheden; den spredes og samles
i \textbf{uvæsentlige} og ligegyldige Relationer. Men
saa væsentlig denne Adskillelse end fra een Side synes
at være, saa væsentlig er tillige dens Ophævelse
i Totalitetens alsidige Udvikling, hvorved Gaaden
løses.\footnote{Det forstaaer sig af sig selv, at vi her ikke tænke paa den empiriske, men logiske Løsning.}.
 Den i sin rene Tilsyneladelse formelle Totalitet
har ved Lovenes og Tingenes Modsætning efterhaanden
vundet et concret Indhold. Enhver af dens
Hovedsider lader sig tilsyne i eiendommelig Form, og
da Formen paa begge Sider \textbf{er} væsentlig og immanent,
saa har dens Indhold tillige en væsentlig Bestemthed.
Som Lovenes Rige er Phænomenalverdenens
rolige Forbillede, beskuet i Reflexionens Speil,
saa er den phænomenale Verden omvendt Lovenes
urolige Afbillede, betragtet i Umiddelbarhedens Lys.
Det staaer endnu tilbage at vise, hvorledes disse
tvende Sider forholde sig til hinanden indbyrdes.
Ere de, hver for sig, kun Sider, hvor er da det Medium,
i hvilket de have deres Selvstændighed? Ere
de hver for sig Totaliteter, har den ene Totalitet
under sin Udvikling adskilt sig i tvende lige selvstændige
Totaliteter, hvorledes vil det da være muligt
at føre en saa absolut Dualisme tilbage til den oprindelige
Grundeenhed? a) Totaliteten er for det Første bestemt som Phænomenalexistents, og Loven er dens egen væsentlige
Side. Som Tingenes Grundlag falder Lovens

Væsen sammen med Existentsens egen immanente
Bestaaen. De Indholdsbestemmelser, der forsvinde
% --- p. 191 ---
\setcounter{footnote}{0}%
\opage{191}som Sider i det lovbestemte Forhold, \textbf{ere} Reflexer
af Tingenes egen umiddelbare Bestemthed, Væsensmomenter
i Hæcceiteten selv. Da nu Tingene
saaledes ikke blot ophæves som Sider i Loven,
men tillige sættes som selvstændige Phænomener i
Hæcceiteten, saa har den phænomenale Verden ikke
alene Lovens Indhold i sig selv, men desforuden
en Mangfoldighed af uvæsentlige, mod Loven ligegyldige
Bestemmelser, altsaa det hele i sig fordoblede
Indhold. Saaledes bliver Phænomenalexistentsen
det Overgribende, den subsisterer i Lovenes
Totalitet, som i sin Andethed; denne Bestaaen er
desuagtet dens egen Reflexion-i-sig, den har selv
Andetheden i sig, og bestaaer i væsentlig Reflexion.
Men idet Phænomenalexistentsen gjennem sin egen
Andethed vender tilbage i sig selv som reflecteret
Verden, kommer den fra den modsatte Side eo
ipso fremmed til sig selv; thi den Umiddelbarhed,
hvormed den fra først af var identisk, er nu bleven
dens Andethed; for at bestaae i denne Andethed,
ophæver den igjen sin reflecterede Bestaaen,
og kommer tilsyne i Tingenes ydre Mangfoldighed.
b) Totaliteten maa dernæst bestemmes som reflecteret
Totalitet, saa at Phænomenalexistentsen bliver en
Side af Lovens eget overgribende Indhold. Da
alle Ting bestaae formedelst Loven, der sætter deres
immanente Grændse, saa er den Bestaaen,
hvorved de falde udenfor hinanden, reflecteret som
Moment i den Bestaaen, ifølge hvilken de falde
sammen med hverandre, og identificeres med Lovens
Væsen. Fra den phænomenale Verdens Synspunkt er hiin den positive, denne den negative
Bestaaen; fra Lovens Synspunkt forholder det sig
derimod omvendt. Idet nu denne Forskjel imellem
% --- p. 192 ---
\setcounter{footnote}{0}%
\opage{192}det Positive og Negative igjen forsvinder i Lovens
egen Dialektik, bliver Lovens Indhold selv sat som
den hele Totalitet. Tingenes umiddelbare Existents
er deres positive Tilværen i Endelighedens Form.
Men da Loven herved er sat som den immanente
Grændse, saa er Tingens positive Bestemthed tillige
et Udtryk for Lovens egen negative Bestemmelse;
ved at komme tilsyne i Tingen som i et
umiddelbart Andet, bliver den sig selv \textbf{et} Andet ɔ:
den adskiller sig i en Fleerhed af bestemte Love..
Loven optager herved Andetheden som integrerende
Moment i sit eget Væsen. Den ene Ting er i
sin Bestaaen afhængig af den anden, thi Loven er
deres Grundlag, den er tillige forskjellig fra den
anden, thi Loven er deres Grændse. Idet nu den
ene Ting subsisterer i sig ved at subsistere i den
andens Bestaaen, saa bestaaer den i Opgivelsen af
sin Selvstændighed og Overgang i en anden Existents;
men da enhver anden Existents, til hvilken
Overgangen skeer, igien bestaaer i sin Andethed,
saa vender Tingenes umiddelbare Selvstændighed
uafbrudt tilbage gjennem dens Negation; den har
altsaa sin Umiddelbarhed i Reflexionen og sin Reflexion
i Umiddelbarheden, saa at denne immanente
Overgang og Tilbagevenden er Lovens Selvbevægelse.
Foruden sit eget væsentlige Indhold omfatter
Lovenes Rige da tillige Phænomenalexistentsens
vexlende Mangfoldighed, og vinder først Væsenhedens
Ro i Bevægelsens Uro.
c) Forsaavidt den \textbf{umiddelbare Totalitet} i Kraft
af sin egen Reflexion sætter sig selv fast i Umiddelbarhedens
Form, er et for-sig-værende Rige af
Love en uvæsentlig Abstraction: \textbf{kun den} phænomenale \textbf{Verden existerer}. Forsaavidt den

\clearpage
% --- PDF 202: the title of 3die Hæfte, bound at the end of KB's copy (r-struct.tsv, F202).
\begin{center}
\vspace*{4em}
{\Huge\textbf{Den speculative Logik}}\\[1em]
{\Large i dens Grundtræk}\\[1em]
af\\[1em]
{\Large\textbf{Lic. Rasmus Nielsen,}}\\[0.5em]
Prof. extr. i Philosophien ved Kjøbenhavns Universitet.\\[2em]
{\large\textbf{3die Hæfte.}}\\[1em]
\textit{1842.}
\end{center}
\end{document}
